\section{Exact circuits for matrix product unitaries}

The constructions of~\cite{RankTwoMPUCircuits2026} give exact circuits
for unitaries with bounded consecutive operator Schmidt ranks. The
rank-two arguments are followed by a general construction using positive
metrics in the real affine spaces of prefix Gram matrices. The latter
removes the spectral conditioning hypothesis at every fixed bond bound.

\begin{theorem}[Contraction does not increase a cut rank]
    \label{thm:mpu_schmidt_span_contraction}
    \lean{Matrix.rank_contract_right_le, Matrix.rank_contract_left_le}
    \leanok
    Let $K$ be a field and let $J,A$ be finite sets. For a matrix
    $M:I\times(J\times A)\to K$ and $f:A\to K$, put
    $M_f(i,j)=\sum_a M(i,(j,a))f(a)$. Then
    $\operatorname{rank}M_f\leq\operatorname{rank}M$.
    Similarly, for finite $I$ and a matrix
    $M:(A\times I)\times J\to K$, contraction
    $M_f(i,j)=\sum_a f(a)M((a,i),j)$ does not increase rank.
\end{theorem}
\begin{proof}\leanok
    In the first case write $M_f=MQ$, where
    $Q_{(j',a),j}=\delta_{j',j}f(a)$.
    In the second write $M_f=PM$, where
    $P_{i,(a,i')}=\delta_{i,i'}f(a)$.
    Rank cannot increase under multiplication on either side.
\end{proof}

\begin{theorem}[A constant number of recursive calls]
    \label{thm:mpu_constant_branching_cost}
    \lean{MPUCircuit.cost_le_pow_of_constant_branching}
    \leanok
    Let $T_j$ be a real sequence, let $a\in\mathbb R$, $C\geq0$ and
    $r\geq2$. Suppose $T_0\leq a$ and $T_{j+1}\leq rT_j+C$ for
    every nonnegative integer $j$. Then
    $T_L\leq(a+C)r^L-C$ for every nonnegative integer $L$.
\end{theorem}
\begin{proof}\leanok
    Induct on $L$. At the step use
    $r((a+C)r^L-C)+C\leq(a+C)r^{L+1}-C$, since $(r-2)C\geq0$.
\end{proof}

\begin{theorem}[Commuting unitary polar factors]
    \label{thm:mpu_two_kraus_polar}
    \lean{Matrix.exists_unitary_polar_factor,
        Matrix.commute_of_unitary_conj_eq,
        Matrix.mul_conjTranspose_eq_of_unitary_polar_factor,
        Matrix.unitary_polar_conj_eq_of_two_kraus,
        Matrix.exists_commuting_unitary_polar_factors}
    \leanok
    Let square complex matrices $E,F$ satisfy
    $E^\dagger E+F^\dagger F=EE^\dagger+FF^\dagger=I$.
    There are unitaries $V,W$ such that
    $E=V\sqrt{E^\dagger E}$, $F=W\sqrt{F^\dagger F}$, and
    $E^\dagger E$ commutes with $V^\dagger W$.
\end{theorem}
\begin{proof}\leanok
    The equality-of-Gram-matrices extension theorem supplies unitary
    polar factors, including at zero singular values. Their output
    Gram matrices are $V(E^\dagger E)V^\dagger$ and
    $W(F^\dagger F)W^\dagger$. The two completeness identities imply
    $V(E^\dagger E)V^\dagger=W(E^\dagger E)W^\dagger$.
    Multiplication by $V^\dagger$ on the left and $W$ on the right
    gives the commutation assertion.
\end{proof}

\begin{theorem}[Coherent recording and erasure of two sectors]
    \label{thm:mpu_two_sector_merging}
    \lean{Matrix.twoSectorEncoder_mul_self,
        Matrix.twoSectorEncoder, Matrix.zeroFlagProjection,
        Matrix.twoSectorEncoder_conjTranspose,
        Matrix.twoSectorEncoder_mem_unitaryGroup,
        Matrix.twoSectorEncoder_merge,
        Matrix.one_kronecker_commute_kronecker_one,
        Matrix.twoSectorEncoder_merge_kronecker}
    \leanok
    Let $P$ be an orthogonal projection and put
    $\mathcal E(P)=\left(\begin{smallmatrix}P&I-P\\I-P&P\end{smallmatrix}\right)$.
    Then $\mathcal E(P)$ is a Hermitian unitary.
    Let $Q_0=\left(\begin{smallmatrix}I&0\\0&0\end{smallmatrix}\right)$
    be the zero-record projection. For any square matrices $X_0,X_1$
    on another factor,
    \[
      \mathcal E(I\otimes P)
      \begin{pmatrix}X_0\otimes I&0\\0&X_1\otimes I\end{pmatrix}
      \mathcal E(I\otimes P)Q_0
      =\begin{pmatrix}X_0\otimes P+X_1\otimes(I-P)&0\\0&0\end{pmatrix}.
    \]
    Thus every input with its record qubit zero has its record zero
    at the end.
\end{theorem}
\begin{proof}\leanok
    Multiply the blocks using $P^2=P$,
    $P(I-P)=(I-P)P=0$, and $(I-P)^2=I-P$.
    Matrices acting on distinct tensor factors commute.
\end{proof}

\begin{theorem}[Combining sector erasure with its final factor]
    \label{thm:mpu_two_sector_cancellation}
    \lean{Matrix.isStarProjection_half_one_add, Matrix.flagHadamard,
        Matrix.flagHadamard_mul_diagonal_mul,
        Matrix.flagHadamard_mul_controlled_mul,
        Matrix.diagonal_mul_twoSectorEncoder,
        Matrix.diagonal_mul_twoSectorEncoder_of_mem_unitaryGroup}
    \leanok
    Put $\mathcal H=2^{-1/2}
      \left(\begin{smallmatrix}I&I\\I&-I\end{smallmatrix}\right)$.
    Then $\mathcal H\operatorname{diag}(I,R)\mathcal H
      =\mathcal E((I+R)/2)$.
    If $R$ is a Hermitian involution, $(I+R)/2$ is an orthogonal
    projection. For a unitary $B$ and any square matrix $C$,
    \[
      \operatorname{diag}(B,B)\mathcal E((I+B^\dagger C)/2)
        =\mathcal H\operatorname{diag}(B,C)\mathcal H.
    \]
\end{theorem}
\begin{proof}\leanok
    The first identity follows by multiplying the Hadamard blocks.
    For the second, multiply those same blocks with $B$ and
    $BB^\dagger C=C$. The projection assertion follows from
    $R^2=I$ and $R^\dagger=R$.
\end{proof}

% The structural subcase arguments immediately below are mathematical proofs.
% The general exact circuit theorem at the end of this section is formalized.
\begin{lemma}[Two-operator simultaneous singular-value decomposition]
    \label{lem:mpu_two_operator_svd}
    Let square complex matrices $E,F$ satisfy
    $E^\dagger E+F^\dagger F=EE^\dagger+FF^\dagger=I$.
    There are unitaries $L,R$ and diagonal matrices $D_E,D_F$ with
    $E=LD_ER$ and $F=LD_FR$.
\end{lemma}
\begin{proof}
    \uses{thm:mpu_two_kraus_polar}
    Put $H=E^\dagger E$, and choose unitary polar extensions
    $E=V\sqrt H$ and $F=W\sqrt{I-H}$.
    The second normalization gives $VHV^\dagger=WHW^\dagger$.
    Thus $V^\dagger W$ commutes with $H$; diagonalize both on the
    eigenspaces of $H$. The resulting left and right bases give the
    simultaneous singular-value decomposition.
\end{proof}

\begin{lemma}[Two controlled sectors in the Schmidt spaces]
    \label{lem:mpu_two_schmidt_sectors}
    Let $U$ be a bipartite unitary of operator Schmidt rank two.
    After possibly interchanging the factors, there are unitaries
    $A_0,A_1$ and matrices $B_0,B_1$ such that
    $U=A_0\otimes B_0+A_1\otimes B_1$.
    The matrices $B_\pm=B_0\pm B_1$ are unitary, and
    $R=B_+^\dagger B_-$ is a Hermitian involution.
    The $A_s$ belong to the first Schmidt space of $U$ and the
    $B_\pm$ to the second.
\end{lemma}
\begin{proof}
    \uses{lem:mpu_two_operator_svd, thm:mpu_circle_factorization}
    Take a Hilbert--Schmidt orthogonal Schmidt decomposition.
    Normalizing one pair and tracing the two unitary identities of
    $U$ gives the hypotheses of Lemma~\ref{lem:mpu_two_operator_svd}
    for the opposite pair. Repeat on the other side. Thus $U$ is
    equivalent under local left and right unitary multiplication to
    a diagonal operator with a unit-modulus coefficient matrix of rank two.
    The circle--line argument gives two projective row types or two
    projective column types. This expresses the diagonal operator as
    two controlled phase terms. Restoring the local factors gives the
    claimed decomposition; $R$ is unitarily conjugate to the difference
    of the two complementary sector projections. The two factor pairs
    are independent, so they span the original Schmidt spaces.
\end{proof}

\begin{theorem}[Exact circuits at cut rank two]
    \label{thm:mpu_rank_two_exact_circuits}
    Let $N\geq1$. Every $N$-qubit unitary whose operator Schmidt rank
    is at most two at each consecutive cut has an exact circuit of
    $O(N^{\log_2 6})$ one- and two-qubit gates, with
    $O(\log(N+1))$ auxiliary qubits initialized and returned to zero.
    Arbitrary pairs may interact. Routing on a line gives gate count
    and depth $O(N^{1+\log_2 6})$ with the same auxiliary-space bound.
\end{theorem}
\begin{proof}
    \uses{thm:mpu_schmidt_span_contraction,
        thm:mpu_constant_branching_cost, lem:mpu_two_schmidt_sectors,
        thm:mpu_two_sector_merging, thm:mpu_two_sector_cancellation}
    At each balanced split use the two-sector decomposition.
    A clean qubit records the eigenvalue of $R$ by a Hadamard,
    controlled $R$, and a second Hadamard. Apply $A_0,A_1$ conditioned
    on the record, erase it, and apply $B_+$.
    Since $R=B_+^\dagger B_-$ and $R=R^\dagger$, the erasure and the
    final $B_+$ combine into two controlled child calls. Recording
    takes two calls, and the $A_s$ take two more, for six in total.
    An external control is retained by computing its conjunction with
    the record into two reusable clean qubits.

    Every child is in an original Schmidt space, so contraction on the
    opposite interval preserves all internal rank bounds.
    Adjoint also preserves these ranks. The recursion has
    $O(6^{\lceil\log_2N\rceil})$ nodes and uses three auxiliary qubits
    per level, with storage reused between calls. At rank one there are
    only two unitary factors. At one site a controlled unitary is an
    allowed two-qubit gate. The local diagonalizing unitaries used in
    the structural proof are never separately implemented.
\end{proof}

\section{Limits of recursion inside the original Schmidt spaces}

The following finite-dimensional examples are discussed in
\cite{MPUSchmidtObstructions2026}. They constrain a possible recursive
proof, rather than providing a circuit complexity lower bound.

\begin{theorem}[A spin-one reflection]
    \label{thm:mpu_spin_one_reflection}
    \lean{MPUSchmidtObstruction.coupling_conjTranspose,
        MPUSchmidtObstruction.spinOne, MPUSchmidtObstruction.coupling,
        MPUSchmidtObstruction.reflection, MPUSchmidtObstruction.paddedSpinOne,
        MPUSchmidtObstruction.qubitReflection,
        MPUSchmidtObstruction.coupling_sq,
        MPUSchmidtObstruction.reflection_mem_unitaryGroup,
        MPUSchmidtObstruction.qubitReflection_mem_unitaryGroup}
    \leanok
    Let $(J_a)_{a=1}^3$ be the Cartesian spin-one matrices,
    $(\sigma_a)_{a=1}^3$ the Pauli matrices, and
    $H=\sum_aJ_a\otimes\sigma_a$. Then
    $H=H^\dagger$ and $H^2=2I-H$.
    Consequently $U=(I+2H)/3$ is unitary.
    Adjoining the identity on a two-dimensional sector gives
    the three-qubit unitary $\widehat U=U\oplus I_2$.
\end{theorem}
\begin{proof}\leanok
    Direct multiplication gives the two identities for $H$.
    Substitution yields $U^\dagger U=U^2=I$.
    The direct-sum extension is unitary on both orthogonal sectors.
\end{proof}

\begin{theorem}[Scalar-plus-skew-symmetric unitary matrices]
    \label{thm:mpu_spin_one_unitary_space}
    \lean{MPUSchmidtObstruction.eq_zero_of_scalarSkew_mem_unitaryGroup,
        MPUSchmidtObstruction.scalarSkew,
        MPUSchmidtObstruction.scalarSkewLinearMap,
        MPUSchmidtObstruction.leftOperatorSpace,
        MPUSchmidtObstruction.scalarSkew_mulVec,
        MPUSchmidtObstruction.eq_smul_one_of_mem_leftOperatorSpace_of_unitary,
        MPUSchmidtObstruction.spinOne_zero_mem_leftOperatorSpace,
        MPUSchmidtObstruction.leftOperatorSpace_not_spanned_by_unitaries}
    \leanok
    Put $W=\{aI_3+K:a\in\mathbb C,\ K^T=-K\}$.
    Every unitary matrix in $W$ is scalar, and $W$ is not spanned
    by its unitary matrices.
\end{theorem}
\begin{proof}\leanok
    A skew-symmetric matrix of odd size has a nonzero kernel vector.
    Norm preservation on that vector forces $|a|=1$.
    Taking the trace of $(aI+K)^\dagger(aI+K)=I$ gives
    $3=3|a|^2+\operatorname{tr}(K^\dagger K)$, hence $K=0$.
    The nonzero spin matrix $J_1$ belongs to $W$ and is not scalar.
\end{proof}

\begin{theorem}[An operator space with no unitary matrix]
    \label{thm:mpu_qubit_unitary_space_obstruction}
    \lean{MPUSchmidtObstruction.paddedScalarSkew_not_mem_unitaryGroup,
        MPUSchmidtObstruction.paddedScalarSkew,
        MPUSchmidtObstruction.paddedScalarSkewLinearMap,
        MPUSchmidtObstruction.qubitLeftOperatorSpace,
        MPUSchmidtObstruction.not_mem_unitaryGroup_of_mem_qubitLeftOperatorSpace}
    \leanok
    Put $D=\operatorname{diag}(1/3,1/3,1/3,1)$ and
    $\widehat W=\{aD+(K\oplus0):a\in\mathbb C,\ K^T=-K\}$.
    No matrix in $\widehat W$ is unitary.
\end{theorem}
\begin{proof}\leanok
    \uses{thm:mpu_spin_one_unitary_space}
    The last coordinate forces $|a|=1$.
    The first three coordinates would give the unitary matrix
    $aI_3/3+K$. The preceding theorem forces $K=0$, so its diagonal
    forces $|a|=3$, a contradiction.
\end{proof}

% The Schmidt-space identifications in the next proposition are not yet
% formalized. The matrix identities and obstructions above are formalized.
\begin{proposition}[A Schmidt-space obstruction at rank four]
    \label{prop:mpu_spin_one_schmidt_obstruction}
    The spin-one reflection has operator Schmidt rank four and left
    Schmidt space $W$. Its three-qubit extension has the same rank
    and left Schmidt space $\widehat W$.
    Thus an original Schmidt space of a unitary can fail to have
    even one unitary element.
\end{proposition}
\begin{proof}
    \uses{thm:mpu_spin_one_reflection,
        thm:mpu_qubit_unitary_space_obstruction}
    The four left factors $I,J_1,J_2,J_3$ are linearly independent,
    as are the four right factors $I,\sigma_1,\sigma_2,\sigma_3$.
    In the extension, replace the left identity coefficient by $D$
    and the spin matrices by $J_a\oplus0$.
    Both factor families remain independent. Their spans are precisely
    the stated spaces.
\end{proof}

\begin{theorem}[The rational controlled Pauli operator]
    \label{thm:mpu_pauli_control_unitary}
    \lean{MPUPauliControl.oneSite_mem_unitaryGroup,
        MPUPauliControl.coordinateInt, MPUPauliControl.pauliCombination,
        MPUPauliControl.oneSite, MPUPauliControl.witness,
        MPUPauliControl.pauliCombination_mul_self,
        MPUPauliControl.pauliCombination_conjTranspose,
        MPUPauliControl.coordinateInt_sum_sq,
        MPUPauliControl.witness_mem_unitaryGroup}
    \leanok
    For a real vector $v$, the matrix $v\cdot\sigma$ is Hermitian
    and its square is $\|v\|_2^2I$.
    Let the eight vectors $n_j$ be the rows of
    \[
      \frac1{15}
      \begin{pmatrix}
       15&0&0\\0&15&0\\0&0&15\\9&12&0\\
       5&-10&-10\\-10&5&-10\\-10&-10&5\\-5&-2&-14
      \end{pmatrix}.
    \]
    Each vector has Euclidean norm one. Hence each $n_j\cdot\sigma$
    and the controlled operator
    $V=\sum_j(n_j\cdot\sigma)\otimes|j\rangle\langle j|$ are unitary.
\end{theorem}
\begin{proof}\leanok
    Expand the two-dimensional Pauli matrices. Each of the eight
    integer rows has squared length $225$.
    The controlled operator is block diagonal in its control label,
    so unitarity follows from unitarity of each block.
\end{proof}

\begin{theorem}[The three-dimensional control coefficient space]
    \label{thm:mpu_pauli_control_space_obstruction}
    \lean{MPUPauliControl.not_forall_normSq_controlLinearCombination,
        MPUPauliControl.controlLinearCombination,
        MPUPauliControl.controlDiagonal, MPUPauliControl.controlCoefficientMap,
        MPUPauliControl.controlCoefficientSpace, MPUPauliControl.coefficient,
        MPUPauliControl.controlDiagonal_not_mem_unitaryGroup,
        MPUPauliControl.controlCoefficientSpace_not_mem_unitaryGroup,
        MPUPauliControl.controlCoefficientMap_injective,
        MPUPauliControl.finrank_controlCoefficientSpace,
        MPUPauliControl.controlCoefficientMap_apply,
        MPUPauliControl.witness_eq_sum_pauli_kronecker,
        MPUPauliControl.coefficient_linearIndependent,
        MPUPauliControl.coefficient_mem_controlCoefficientSpace,
        MPUPauliControl.controlCoefficientSpace_eq_span}
    \leanok
    Let $T(v)=\operatorname{diag}(v\cdot n_j)$ for $v\in\mathbb C^3$,
    and put $B_a=T(e_a)$. The map $T$ is injective, its range has
    complex dimension three, and its range is
    $\operatorname{span}_{\mathbb C}\{B_1,B_2,B_3\}$.
    It contains no unitary matrix. The controlled Pauli operator
    satisfies $V=\sum_{a=1}^3\sigma_a\otimes B_a$.
\end{theorem}
\begin{proof}\leanok
    \uses{thm:mpu_pauli_control_unitary}
    The first three diagonal entries recover the coordinates of $v$.
    Unit modulus at the first six entries would force those
    coordinates to have modulus one and pairwise zero real inner
    products. Three nonzero pairwise orthogonal vectors in the real
    plane $\mathbb C$ would be linearly independent, contradicting
    its real dimension two. The range and tensor expansion follow
    by linearity of the displayed coefficients.
\end{proof}

\begin{proposition}[A Schmidt-space obstruction at bond dimension three]
    \label{prop:mpu_pauli_control_schmidt_obstruction}
    There is a four-qubit unitary with an open-boundary matrix product
    representation of bond dimensions $3,3,2$ whose right Schmidt
    space at the first cut contains no unitary matrix.
\end{proposition}
\begin{proof}
    \uses{thm:mpu_pauli_control_unitary, thm:mpu_pauli_control_space_obstruction}
    Let $q_0=e_1$, $q_1=e_2$, $q_2=e_3$,
    $q_3=(3e_1+4e_2)/5$, and
    $O=I-2\mathbf 1\mathbf 1^T/3$, with $\mathbf 1=(1,1,1)^T$.
    Put $n_j=q_j$ and $n_{j+4}=Oq_j$ for $0\leq j<4$.
    The controlled Pauli operator
    $V=\sum_{j=0}^7(n_j\cdot\sigma)\otimes|j\rangle\langle j|$
    is unitary. At the second cut its coefficient vectors
    $(q_j,Oq_j)$ span a space of dimension three; at the last cut
    the right factors are diagonal one-qubit matrices.
    Its right Schmidt space at the first cut consists of
    $\operatorname{diag}(v\cdot n_j)$, for $v\in\mathbb C^3$.
    Unitarity would give $|v_1|=|v_2|=|v_3|=1$ and
    $\operatorname{Re}(v_a\overline{v_b})=0$ for all distinct $a,b$.
    The mixed vector $q_3$ gives the first orthogonality equation,
    and $Oe_1,Oe_2$ give the other two. Three nonzero pairwise
    orthogonal vectors in the real plane $\mathbb C$ are impossible.
\end{proof}

\begin{theorem}[Exact controlled Pauli circuits]
    \label{thm:mpu_orthogonal_walk_circuits}
    Let $n_0\in\mathbb R^3$ be a unit vector and $O_i\in O(3)$.
    Put $n(x)=O_1^{x_1}\cdots O_n^{x_n}n_0$ and
    $U=\sum_x(n(x)\cdot\sigma)\otimes|x\rangle\langle x|$.
    This unitary has an open-boundary matrix product representation
    of maximal bond dimension at most three, and an exact circuit
    with $2n$ two-qubit gates, at most $n+1$ one-qubit gates, and
    no auxiliary qubits. With its target at one end of a line,
    it has an adjacent-gate circuit of size and depth $O(n)$.
\end{theorem}
\begin{proof}
    \uses{thm:mpu_pauli_control_unitary}
    Write $s_i=\det O_i$ and $R_i=s_iO_i\in SO(3)$.
    A rotation through angle $\theta$ about unit axis $u$ lifts to
    $V_i=\cos(\theta/2)I-i\sin(\theta/2)(u\cdot\sigma)$.
    The Pauli product rule gives
    $V_i(v\cdot\sigma)V_i^\dagger=(R_iv)\cdot\sigma$.
    Thus, with $W(x)=V_1^{x_1}\cdots V_n^{x_n}$,
    $n(x)\cdot\sigma=(\prod_i s_i^{x_i})
      W(x)(n_0\cdot\sigma)W(x)^\dagger$.
    Apply the two controlled products and the central target gate;
    each $s_i=-1$ contributes a $Z$ on its control.
    On a line, move the target past the controls during the first
    product and return it during the second, using two sweeps of
    adjacent swaps. Every logical position is restored.
\end{proof}

\begin{theorem}[Approximation with a small quantum target]
    \label{thm:mpu_controlled_target_approximation}
    Let $n,t\geq1$ and
    $U=\sum_{y\in\{0,1\}^n}|y\rangle\langle y|\otimes V(y)$
    be a unitary with an open-boundary matrix product representation
    of maximal bond dimension $D$, where $V(y)$ acts on $t$ qubits.
    Put $N=n+t$. For $0<\varepsilon<1$, there is a nonuniform
    circuit of $\operatorname{poly}(N,D,2^t,\log(1/\varepsilon))$
    one- and two-qubit gates whose implemented physical unitary
    $\widetilde U$ satisfies $\|\widetilde U-U\|\leq\varepsilon$.
    The number of auxiliary qubits has the same polynomial bound,
    and all auxiliaries return exactly to zero.
    In particular, $t\leq c\log_2N$ for a fixed $c$ gives polynomial
    resources in $N,D,\log(1/\varepsilon)$.
\end{theorem}
\begin{proof}
    Put $q=2^t$. A right-canonical representation of the normalized
    Choi vector has letter matrices of norm at most one.
    Compute the $q^2$ block entries reversibly using matrix--vector
    contractions. Largest-entry-pivoted QR decomposes each computed
    block into $O(q^2)$ two-level rotations and row permutations.
    Compare the numerical procedure to exact QR with the same pivot
    choices. Pivots remain at least $1/(4\sqrt q)$; conservative
    error amplification is at most $(Cq)^{q^2}$.
    Its logarithm requires only polynomially many precision bits:
    \[
      O\bigl(N+\log(ND+1)+q^2\log(q+1)
        +\log(1/\varepsilon)\bigr).
    \]
    Convert the recorded rotations to binary angles, retaining all
    diagonal phases. A Gray path of length at most $t$ implements
    each two-level operation using basis-pattern controls, with
    reusable conjunction qubits. Uniform block approximation gives
    $\|\widetilde U-U\|=\max_y\|\widetilde V(y)-V(y)\|$.
    Since every computed bit depends only on unchanged controls,
    reversing the arithmetic erases all auxiliaries exactly.
    The detailed finite-precision proof is in
    \cite{MPUSchmidtObstructions2026}.
\end{proof}

\section{Completions and general approximation}

\begin{theorem}[Exchanging orthogonal isometric images]
    \label{thm:mpu_orthogonal_isometry_swap}
    \lean{Matrix.gram_sub_of_orthogonal_isometries,
        Matrix.orthogonalIsometrySwap,
        Matrix.orthogonalIsometrySwap_mem_unitaryGroup,
        Matrix.orthogonalIsometrySwap_mul_left,
        Matrix.orthogonalIsometrySwap_mul_right}
    \leanok
    Let $J,V:\mathcal H\to\mathcal K$ be isometries between
    finite-dimensional complex Hilbert spaces with $J^\dagger V=0$.
    Then $(V-J)^\dagger(V-J)=2I$, and
    $S=I-(V-J)(V-J)^\dagger$ is a unitary satisfying $SJ=V$ and $SV=J$.
\end{theorem}
\begin{proof}\leanok
    Orthogonality gives both cross Gram terms zero. Put $W=V-J$;
    then $W^\dagger W=2I$ and $(WW^\dagger)^2=2WW^\dagger$.
    Thus $S=S^\dagger$ and $S^2=I$. Multiplication by $J$ or $V$
    gives the two asserted actions.
\end{proof}

For a Kraus isometry
$V\psi=\sum_{a=1}^r A_a\psi\otimes|a\rangle$, take
$J\psi=\psi\otimes|0\rangle$. The completion uses only
$I,A_a,A_a^\dagger,A_aA_b^\dagger$. If the collective isometry has
internal operator cut ranks at most $D$ and the flag is at an endpoint,
the completion has ranks at most $(D+1)^2$. This bounds one completion;
it does not bound repeated completion by a constant independent of
recursion depth.

\begin{theorem}[A positive defect can increase bond dimension]
    \label{thm:mpu_positive_defect_rank_growth}
    For every $m\geq1$ there is a contraction $E$ which is an operator
    Schmidt factor of an open-boundary bond-two MPU with an exact
    $2m$-gate circuit, such that the positive defect
    $\sqrt{I-E^\dagger E}$ has operator cut rank $2^m$.
    The same contraction admits a signed unitary completion of bond two.
\end{theorem}
\begin{proof}
    Put $q=2^m$, $h=\pi/(2q)$ and
    $\Theta|i,j\rangle=h(i-j)|i,j\rangle$ on two $m$-qubit blocks.
    Take $E=\cos\Theta$, $F=\sin\Theta$, and
    $U=E\otimes I+iF\otimes X$ with its target at an endpoint.
    Controlled target rotations give the exact circuit.
    The positive defect has coefficient matrix
    $K_{ij}=\sin(h|i-j|)$. If $Kv=0$, its interior second differences
    give $2\sin(h)v_a=0$; the first and last rows then kill the
    endpoint coordinates. Thus $K$ is invertible.
    Both $E$ and $F$ have rank two at this cut, and
    $\left(\begin{smallmatrix}E&-F\\F&E\end{smallmatrix}\right)$
    supplies the signed completion.
\end{proof}

\begin{theorem}[General quasipolynomial approximation]
    \label{thm:mpu_general_quasipolynomial_approximation}
    Let $N\geq1$ and let $U$ be an $N$-qubit unitary of consecutive
    operator cut ranks at most $D$. For $0<\varepsilon<1$, there are
    a nonuniform circuit unitary $W$ and a clean auxiliary embedding
    $J\psi=\psi\otimes|0\cdots0\rangle$ with
    $\|WJ-JU\|\leq\varepsilon$.
    For fixed $D$, both gate and auxiliary-qubit counts are at most
    \[
      \exp\!\left\{C_D\log(N+1)
        \log\bigl(\log(N+1)+\log(1/\varepsilon)+2\bigr)\right\}.
    \]
    This bounds the auxiliary output error, without asserting exact
    auxiliary erasure or a bound on arbitrary auxiliary inputs.
\end{theorem}
\begin{proof}
    \uses{thm:mpu_schmidt_span_contraction}
    At each balanced cut, an Auerbach basis of the left Schmidt space
    and its norm-one coordinate functionals write a contraction as
    $A=\sum_{i=1}^r A_i\otimes B_i$, with $r\leq D$ and all factors
    contractions of inherited internal cut ranks.
    Recursively implement compressed blocks $A_i/s,B_i/s$ for a fixed
    numerical $s$. Their linear combination has ideal block
    $A/(rs^2)$. An odd sine Fourier series of coefficient absolute
    sum at most one implements linear amplification on the required
    singular-value interval through Grover powers. A compact convolution
    kernel gives truncation degree $O_D((1+\log(1/\eta))^2)$.
    The Chebyshev perturbation estimate
    $\|T_n(X)-T_n(Y)\|\leq n^2\|X-Y\|$, applied to Hermitian dilations,
    controls child errors. The combined gate and auxiliary cost obeys
    $F_{\rm parent}\leq C_D M^c\sum F_{\rm child}$.
    Along $\lceil\log_2N\rceil$ levels the required logarithmic
    precision grows at most polynomially in the initial precision
    and the number of levels, giving the displayed bound.
    Final constant-cost amplification of the unitary block yields
    a block $B$ with $\|B-U\|=O(\varepsilon^2)$; unitarity then gives
    $\|WJ-JU\|\leq\sqrt{2\|B-U\|}$.
    The complete Fourier and circuit estimates are in
    \cite{MPUSchmidtObstructions2026}, Section~7.
\end{proof}

\section{Uniform families without exact-length full rank}

\begin{theorem}[Uniform all-length polynomial circuits]
    \label{thm:mpu_uniform_averaged_boundary_circuits}
    Fix a physical dimension $d\geq2$, a finite-dimensional bond space,
    matrices $A^{ij}$ and boundary vectors $l,r$. Suppose
    \[
      (U_N)_{\mathbf i,\mathbf j}
        =lA^{i_1j_1}\cdots A^{i_Nj_N}r
    \]
    is unitary for every $N\geq1$. Then $U_N$ has an exact circuit
    of polynomial size and depth, with $O(N)$ auxiliary qudits which
    begin and end in zero. The gates may be restricted to adjacent
    sites on a line. The constants and polynomial exponent depend
    on the fixed presentation; no exact-length full-rank assumption
    is required.
\end{theorem}
\begin{proof}
    Let $E_+$ be the span of all positive-word columns $A_wr$ and
    quotient by the vectors annihilated by all positive-word rows
    $lA_u$. This preserves the operators at lengths at least two,
    using endpoint tensors $lA^{ij}$ and $A^{ij}r$ on the quotient.
    Write $w$ for its dimension. Cumulative cap lengths $1,\ldots,w$
    span the quotient and its dual: a plateau in either cumulative
    span is invariant and therefore terminal.

    Choose Hilbert coordinates on this quotient. For each cap length
    let $P_m$ and $Q_n$ be its left and right Gram matrices, with
    cap inputs averaged and cap outputs summed. Unitarity at length
    $m+k+n$ gives, for a central interval tensor $T$ of length $k$,
    \[
      \sum_{\mathbf a}\operatorname{Tr}
        (T_{\mathbf a,\mathbf b}^\dagger P_m
         T_{\mathbf a,\mathbf b'}Q_n)
        =\delta_{\mathbf b,\mathbf b'}.
    \]
    This identity is separately linear in the two Gram matrices.
    Thus $P=w^{-1}\sum_{m=1}^wP_m$ and
    $Q=w^{-1}\sum_{n=1}^wQ_n$ satisfy every interval identity and
    are positive definite. One-sided cap traces give the endpoint
    identities. The Grams are averaged before taking square roots.

    Put $L=P^{1/2}$ and $R=Q^{1/2}$. The amplitudes $LTR$ define
    fixed interval isometries with two outer flags. Their flags merge
    by the contraction $R^{-1}L^{-1}$, whose finite norm is a constant
    of the family. A fixed unitary dilation and exact subspace
    amplification give a bounded number of child invocations per
    merge. On a balanced tree the cost satisfies
    $H(n)\leq c(H(\lceil n/2\rceil)+H(\lfloor n/2\rfloor))+cn$,
    and is polynomial. The tree has $O(N)$ logical flags, and the
    image reflections reuse one $O(N)$ reversible scratch pool.
    All inner flags and scratch are erased exactly; the root has
    no outer flags. Adjacent swaps add at most a linear cost factor.
    The complete quotient, normalization, amplification and auxiliary
    count are proved in \cite{RankTwoMPUCircuits2026}, Section~4.
\end{proof}

This removes Assumption~1 from the uniform-bulk circuit theorem of
\cite{Styliaris2025MPUCircuits}. The averaging argument uses the same
fixed presentation at every positive length. The following affine
construction removes that restriction.
Fixed trace boundaries $\operatorname{Tr}(bA_{\rm word})$ are included:
use the enlarged bond space of matrices, with each letter acting by
left multiplication, right vector $b$ and left functional the trace.

\section{Polynomial exact circuits at arbitrary fixed cut rank}

\begin{theorem}[Real affine identities of cap Grams]
    \label{thm:mpu_real_affine_gram_identities}
    \lean{MPUCircuit.isHermitian_of_mem_real_affineSpan,
        MPUCircuit.trace_mul_eq_of_mem_real_affineSpan}
    \leanok
    The real affine hull of a set of Hermitian matrices consists of
    Hermitian matrices. If $\operatorname{Tr}(XQ)=z$ on the original
    set, with $Q$ fixed, the same identity holds on its real affine hull.
\end{theorem}
\begin{proof}\leanok
    Both properties are preserved by
    $X,Y,Z\mapsto c(X-Y)+Z$ for every real $c$.
\end{proof}

\begin{theorem}[Compactness and a positive determinant maximizer]
    \label{thm:mpu_compact_affine_gram_slice}
    \lean{Matrix.isCompact_setOf_posSemidef_trace_mul_eq_one,
        Matrix.isCompact_affineSubspace_inter_posSemidef,
        Matrix.PosSemidef.posDef_of_isMaxOn_det_re,
        Matrix.exists_posDef_isMaxOn_det_re_of_isCompact,
        Matrix.exists_posDef_det_re_max_of_trace_mul_eq_one}
    \leanok
    Let $\mathcal C$ be a real affine space of complex square matrices,
    and let $Q>0$ satisfy $\operatorname{Tr}(XQ)=1$ for all
    $X\in\mathcal C$. Then $\mathcal C\cap\{X\succeq0\}$ is compact.
    If this slice contains a positive definite matrix, its determinant
    has a positive definite maximizer. More generally, this last
    assertion holds for any compact set of positive semidefinite
    matrices containing a positive definite matrix.
\end{theorem}
\begin{proof}\leanok
    The entire set $\{X\succeq0:\operatorname{Tr}(XQ)=1\}$ is the
    continuous image of the compact set of density matrices under
    $Y\mapsto Q^{-1/2}YQ^{-1/2}$. Intersect it with the closed affine
    space $\mathcal C$. The determinant is continuous and real on
    positive semidefinite matrices. Its maximum is strictly positive
    by comparison with the assumed positive definite matrix, and a
    positive semidefinite matrix with nonzero determinant is positive
    definite.
\end{proof}

\begin{definition}[The dimension-normalized dual metric]
    \label{def:mpu_dual_gram_metric}
    \lean{MPUCircuit.dualGramMetric}
    \leanok
    For a positive definite $r\times r$ matrix $P$, with $r\geq1$,
    put $Q=P^{-1}/r$.
\end{definition}

\begin{theorem}[Determinant stationarity on a Hermitian affine space]
    \label{thm:mpu_determinant_stationarity}
    \lean{Matrix.hasDerivAt_det_one_add_smul,
        Matrix.hasDerivAt_det_add_smul,
        Matrix.IsHermitian.star_trace_mul_eq,
        Matrix.trace_nonsing_inv_mul_eq_zero_of_isLocalMax_det,
        Matrix.PosDef.eventually_posDef_add_real_smul,
        Matrix.PosDef.trace_nonsing_inv_mul_sub_eq_zero_of_isMaxOn_det}
    \leanok
    For a nonsingular complex square matrix $P$ and any matrix $H$,
    \[
      \left.\frac{d}{dz}\det(P+zH)\right|_{z=0}
        =\det(P)\operatorname{Tr}(P^{-1}H).
    \]
    If $P>0$ maximizes the determinant on the positive semidefinite
    part of a real affine space $\mathcal C$ of Hermitian matrices,
    then $\operatorname{Tr}(P^{-1}(X-P))=0$ for every
    $X\in\mathcal C$.
\end{theorem}
\begin{proof}\leanok
    Differentiate the determinant polynomial of $I+zP^{-1}H$ and
    multiply by $\det P$. The positive definite matrices form the
    interior of the positive semidefinite cone within the Hermitian
    matrices, so both signs of each Hermitian affine direction are
    feasible near $P$. Fermat's theorem gives a zero real derivative.
    The trace of a product of Hermitian matrices is real by cyclicity
    of the trace and conjugate transposition. Cancel the nonzero
    real determinant to obtain the complex trace identity.
\end{proof}

\begin{theorem}[Algebra of the balanced dual metric]
    \label{thm:mpu_dual_gram_metric_algebra}
    \lean{MPUCircuit.dualGramMetric_posDef,
        MPUCircuit.trace_mul_dualGramMetric_eq_one,
        MPUCircuit.dualGramMetric_mul, MPUCircuit.mul_dualGramMetric,
        MPUCircuit.dualGramMetric_inv,
        MPUCircuit.trace_dualGramMetric_inv_mul_inv}
    \leanok
    \uses{def:mpu_dual_gram_metric}
    The dual metric is positive definite and satisfies
    \[
      PQ=QP=I/r,\qquad Q^{-1}=rP,\qquad
      \operatorname{Tr}(Q^{-1}P^{-1})=r^2.
    \]
    If $\operatorname{Tr}(P^{-1}(X-P))=0$, then
    $\operatorname{Tr}(XQ)=1$.
\end{theorem}
\begin{proof}\leanok
    Use $PP^{-1}=P^{-1}P=I$ and $\operatorname{Tr}I=r$.
\end{proof}

\begin{theorem}[A balanced positive metric on an affine space]
    \label{thm:mpu_affine_gram_balancing}
    \lean{MPUCircuit.exists_posDef_dualGramMetric_on_affineSubspace}
    \leanok
    Let $\mathcal C$ be a real affine subspace of Hermitian $r\times r$
    matrices, with $r\geq1$. Suppose it contains a positive definite matrix, and
    there is a positive definite $\overline Q$ with
    $\operatorname{Tr}(P\overline Q)=1$ for all $P\in\mathcal C$.
    Then there are $P_*,Q_*>0$ such that
    \[
      P_*\in\mathcal C,\qquad
      Q_*=P_*^{-1}/r,\qquad
      \operatorname{Tr}(PQ_*)=1\quad(P\in\mathcal C),\qquad
      \operatorname{Tr}(Q_*^{-1}P_*^{-1})=r^2.
    \]
\end{theorem}
\begin{proof}\leanok
    \uses{thm:mpu_compact_affine_gram_slice,
        thm:mpu_determinant_stationarity,
        thm:mpu_dual_gram_metric_algebra}
    The compactness lemma supplies a positive definite determinant
    maximizer $P_*$. Both signs of every affine
    tangent direction $H$ are feasible near $P_*$, so
    $\operatorname{Tr}(P_*^{-1}H)=0$ by the derivative of the determinant.
    Apply this to $H=P-P_*$ and use
    $\operatorname{Tr}(P_*^{-1}P_*)=r$. The final trace identity
    follows from the algebra of the dual metric.
\end{proof}

\begin{definition}[Contraction of balanced square-root weights]
    \label{def:mpu_balanced_gram_contraction}
    \lean{MPUCircuit.balancedGramContraction}
    \leanok
    \uses{def:mpu_dual_gram_metric}
    For positive $P$ and its dual $Q=P^{-1}/r$, the joining
    contraction is $C=(\sqrt Q)^{-1}(\sqrt P)^{-1}$.
\end{definition}

\begin{theorem}[The contraction of balanced metrics]
    \label{thm:mpu_balanced_gram_contraction}
    \lean{MPUCircuit.balancedGramContraction_eq,
        MPUCircuit.balancedGramContraction_gram,
        MPUCircuit.balancedGramContraction_vec_inner}
    \leanok
    \uses{def:mpu_balanced_gram_contraction}
    For $P>0$ and $Q=P^{-1}/r$, put
    $C=(\sqrt Q)^{-1}(\sqrt P)^{-1}$. Then
    $C=\sqrt r\,I$, $C^\dagger C=rI$, and
    $\langle\operatorname{vec}C,\operatorname{vec}C\rangle=r^2$.
    Thus the corresponding contraction bra has norm $r$.
\end{theorem}
\begin{proof}\leanok
    The positive square root commutes with inversion and with
    multiplication by nonnegative real scalars. Therefore
    $(\sqrt Q)^{-1}=\sqrt{rP}=\sqrt r\sqrt P$.
    Cancel $\sqrt P$ and take the trace of $C^\dagger C$.
\end{proof}

\begin{theorem}[Optimality within positive-metric joining]
    \label{thm:mpu_balanced_gram_contraction_optimality}
    \lean{MPUCircuit.card_sq_le_trace_inv_mul_inv_of_trace_mul_eq_one,
        MPUCircuit.trace_dualGramMetric_inv_mul_inv_re,
        MPUCircuit.dualGramMetric_minimizes_inverse_trace}
    \leanok
    \uses{def:mpu_dual_gram_metric}
    Let $P,Q$ be positive definite metrics on a bond of dimension
    $r\geq1$, with $\operatorname{Tr}(PQ)=1$. Then
    \[
        r^2\leq\operatorname{Re}\operatorname{Tr}(Q^{-1}P^{-1}).
    \]
    For fixed $P$, the dual metric $Q_*=P^{-1}/r$ is positive,
    satisfies the normalization and attains equality. Hence the
    joining contraction norm $r$ is optimal among these normalized
    positive metrics. This assertion concerns this joining method;
    it gives no lower bound on arbitrary circuit complexity.
\end{theorem}
\begin{proof}\leanok
    \uses{thm:mpu_dual_gram_metric_algebra}
    Set $X=\sqrt Q\sqrt P$ and $Y=\sqrt{Q^{-1}}\sqrt{P^{-1}}$.
    Then $X^\dagger Y=I$,
    $\operatorname{Tr}(X^\dagger X)=1$ and
    $\operatorname{Tr}(Y^\dagger Y)=\operatorname{Tr}(Q^{-1}P^{-1})$.
    Hilbert--Schmidt Cauchy--Schwarz gives the inequality.
    The dual-metric identity gives equality for $Q_*$.
\end{proof}

\begin{definition}[Density-weighted interval Grams]
    \label{def:mpu_density_interval_grams}
    \lean{MPUCircuit.intervalGramTransfer,
        MPUCircuit.prefixInputGram, MPUCircuit.intervalGramTransferLinearMap,
        MPUCircuit.intervalInputGram}
    \leanok
    For interval matrices $A_{ab}:\mathbb C^s\to\mathbb C^r$, put
    \[
      \mathcal T_\rho(P)=\sum_{a,b,c}\rho_{cb}A_{ab}^\dagger PA_{ac},
      \qquad
      K(P,Q)_{bc}=\sum_a\operatorname{Tr}(A_{ab}^\dagger PA_{ac}Q).
    \]
    The first map is linear in $P$. For prefix rows $F_{ab}$, their
    density-weighted Gram is $P_F(\rho)=\mathcal T_\rho(I)$, where
    the left virtual space is one dimensional.
\end{definition}

\begin{theorem}[Concatenation of density-weighted transfers]
    \label{thm:mpu_density_transfer_concatenation}
    \lean{MPUCircuit.intervalGramTransfer_concat,
        MPUCircuit.prefixInputGram_concat}
    \leanok
    \uses{def:mpu_density_interval_grams}
    For consecutive intervals $A,B$, the transfer of their concatenation
    with input $\rho\otimes\sigma$ is
    $\mathcal T^B_\sigma(\mathcal T^A_\rho(P))$.
    In particular,
    $P_{FA}(\rho\otimes\sigma)=\mathcal T^A_\sigma(P_F(\rho))$.
\end{theorem}
\begin{proof}\leanok
    Expand the finite sums, use $(AB)^\dagger=B^\dagger A^\dagger$,
    and permute the independent physical indices.
\end{proof}

\begin{theorem}[Density matrices determine the entire interval Gram]
    \label{thm:mpu_density_gram_separation}
    \lean{MPUCircuit.trace_intervalGramTransfer,
        MPUCircuit.eq_one_of_trace_density_mul_eq_one,
        MPUCircuit.intervalInputGram_eq_one_of_density_tests}
    \leanok
    \uses{def:mpu_density_interval_grams}
    One has $\operatorname{Tr}(Q\mathcal T_\rho(P))
    =\operatorname{Tr}(\rho K(P,Q))$ for every $\rho$.
    If this value is one for every density matrix $\rho$, then
    $K(P,Q)=I$. More generally, any complex matrix $K$ satisfying
    $\operatorname{Tr}(\rho K)=1$ for every density matrix is the
    identity; no Hermitian hypothesis on $K$ is required.
\end{theorem}
\begin{proof}\leanok
    Trace cyclicity and permutation of the finite sums give the
    pairing identity. For every nonzero vector $v$, test against
    $|v\rangle\langle v|/\langle v,v\rangle$ to obtain
    $\langle v,(K-I)v\rangle=0$. Complex polarization gives $K=I$,
    including all off-diagonal entries.
\end{proof}

\begin{definition}[The real affine hull of physical prefix Grams]
    \label{def:mpu_prefix_gram_affine_hull}
    \lean{MPUCircuit.densityPrefixGrams,
        MPUCircuit.prefixGramAffineHull,
        MPUCircuit.intervalGramTransferAffineMap}
    \leanok
    \uses{def:mpu_density_interval_grams}
    Let $\mathcal S_F=\{P_F(\rho):\rho\succeq0,
    \operatorname{Tr}\rho=1\}$, and let $\mathcal C_F$ be its
    real affine hull. The interval transfer is also an affine map
    over the real scalars.
\end{definition}

\begin{theorem}[Positivity of physical prefix Grams]
    \label{thm:mpu_prefix_gram_positivity}
    \lean{MPUCircuit.prefixInputGram_eq_sum_sandwich,
        MPUCircuit.prefixInputGram_posSemidef,
        MPUCircuit.posSemidef_of_mem_densityPrefixGrams,
        MPUCircuit.isHermitian_of_mem_prefixGramAffineHull}
    \leanok
    \uses{def:mpu_prefix_gram_affine_hull}
    Each physical prefix Gram is positive semidefinite, and every
    matrix in $\mathcal C_F$ is Hermitian. In the convention
    $(F_a)_{bv}=F_{ab}(v)$, one has
    $P_F(\rho)=\sum_aF_a^\dagger\rho^T F_a$.
\end{theorem}
\begin{proof}\leanok
    Expand the matrix entries to obtain the sandwich formula.
    The transpose of a positive semidefinite matrix is positive
    semidefinite. Each sandwich and their sum are therefore positive
    semidefinite. Apply the real affine Hermitian closure lemma.
    \uses{thm:mpu_real_affine_gram_identities}
\end{proof}

\begin{theorem}[Compatibility of the entire affine hull across intervals]
    \label{thm:mpu_prefix_affine_hull_transfer}
    \lean{MPUCircuit.intervalGramTransfer_mem_densityPrefixGrams_concat,
        MPUCircuit.intervalGramTransfer_mem_prefixGramAffineHull_concat}
    \leanok
    \uses{def:mpu_prefix_gram_affine_hull}
    For every interval density matrix $\sigma$,
    $\mathcal T^A_\sigma(\mathcal S_F)\subseteq\mathcal S_{FA}$ and
    $\mathcal T^A_\sigma(\mathcal C_F)\subseteq\mathcal C_{FA}$.
    The density factors may each be entangled throughout their
    respective intervals.
\end{theorem}
\begin{proof}\leanok
    \uses{thm:mpu_density_transfer_concatenation}
    The concatenation identity sends $P_F(\rho)$ to
    $P_{FA}(\rho\otimes\sigma)$. Positivity and unit trace are
    preserved by the tensor product. An affine map sends the affine
    hull of a set into the affine hull of its image.
\end{proof}

\begin{definition}[A vectorized weighted interval]
    \label{def:mpu_vectorized_weighted_interval}
    \lean{MPUCircuit.vectorizedWeightedInterval}
    \leanok
    For square bond weights $L,R$, define an interval map whose
    amplitudes, including the two virtual output indices $u,v$, are
    $(LA_{ab}R)_{uv}$.
\end{definition}

\begin{theorem}[The full isometry identity from density tests]
    \label{thm:mpu_vectorized_density_isometry}
    \lean{MPUCircuit.vectorizedWeightedInterval_gram,
        MPUCircuit.vectorizedWeightedInterval_isometry_of_density_tests}
    \leanok
    \uses{def:mpu_density_interval_grams,
        def:mpu_vectorized_weighted_interval}
    The vectorized map $V$ has Gram $V^\dagger V=K(L^\dagger L,RR^\dagger)$.
    Hence $V$ is an isometry if
    $\operatorname{Tr}(RR^\dagger\mathcal T_\rho(L^\dagger L))=1$
    for every input density matrix $\rho$.
\end{theorem}
\begin{proof}\leanok
    \uses{thm:mpu_density_gram_separation}
    The inner product of two vectorized matrices is their
    Hilbert--Schmidt trace pairing. Sum this identity over the
    physical output index, and apply density separation.
\end{proof}

\begin{theorem}[Interval isometries from affine Gram normalization]
    \label{thm:mpu_affine_interval_isometry}
    \lean{MPUCircuit.intervalInputGram_eq_one_of_prefixGramAffineHull,
        MPUCircuit.vectorizedWeightedInterval_isometry_of_prefixGramAffineHull}
    \leanok
    \uses{def:mpu_prefix_gram_affine_hull,
        def:mpu_vectorized_weighted_interval}
    Suppose $P\in\mathcal C_F$ and $P,Q\succeq0$, with
    $\operatorname{Tr}(XQ)=1$ for every $X\in\mathcal C_{FA}$.
    The interval map with amplitudes $\sqrt P A_{ab}\sqrt Q$ is
    an isometry.
\end{theorem}
\begin{proof}\leanok
    \uses{thm:mpu_prefix_affine_hull_transfer,
        thm:mpu_vectorized_density_isometry}
    For every density matrix $\rho$, transfer membership gives
    $\mathcal T_\rho(P)\in\mathcal C_{FA}$, so its trace pairing
    with $Q$ is one. Apply the full density-test isometry identity
    and the positive square-root identities.
\end{proof}

\begin{theorem}[Balanced metrics at two consecutive cuts]
    \label{thm:mpu_two_cut_balanced_isometry}
    \lean{MPUCircuit.exists_balanced_prefix_metrics_interval_isometry}
    \leanok
    \uses{def:mpu_prefix_gram_affine_hull,
        def:mpu_dual_gram_metric, def:mpu_vectorized_weighted_interval}
    Suppose the two affine spaces $\mathcal C_F$ and $\mathcal C_{FA}$
    each contain a positive definite matrix and each admit a positive
    definite trace normalizer. Their bond dimensions are $r,s\geq1$.
    There exist positive definite $P\in\mathcal C_F$ and
    $S\in\mathcal C_{FA}$ whose dual metrics normalize their respective
    affine spaces and satisfy the inverse-metric trace identities
    $r^2,s^2$. The interval map with amplitudes
    $\sqrt P A_{ab}\sqrt{S^{-1}/s}$ is an isometry.
\end{theorem}
\begin{proof}\leanok
    \uses{thm:mpu_affine_gram_balancing,
        thm:mpu_affine_interval_isometry,
        thm:mpu_prefix_gram_positivity}
    Apply determinant balancing to both Hermitian affine spaces,
    then use the earlier chosen metric and the later dual metric in
    the affine interval-isometry theorem.
\end{proof}

\begin{definition}[Prefix and suffix Grams at a finite cut]
    \label{def:mpu_finite_cut_cap_grams}
    \lean{MPUPrefixGram.factorizedOperator, MPUPrefixGram.prefixGram,
        MPUPrefixGram.suffixGram}
    \leanok
    Let $F_a:\mathbb C^r\to\mathbb C^B$ and
    $G_c:\mathbb C^E\to\mathbb C^r$. Define
    $U_{(a,c),(b,e)}=(F_aG_c)_{be}$ and
    \[
      P_F(\rho)=\sum_aF_a^\dagger\rho^T F_a,\qquad
      Q_G(\tau)=\sum_cG_c\tau G_c^\dagger.
    \]
\end{definition}

\begin{theorem}[Global unitarity normalizes the actual cap Grams]
    \label{thm:mpu_global_cap_gram_normalization}
    \lean{MPUPrefixGram.trace_prefixGram_mul_suffixGram,
        MPUPrefixGram.prefixGram_posSemidef,
        MPUPrefixGram.suffixGram_posSemidef,
        MPUPrefixGram.trace_prefixGram_mul_suffixGram_of_isIsometry,
        MPUPrefixGram.trace_prefixGram_mul_suffixGram_normalizedIdentity,
        MPUPrefixGram.prefixGram_eq_prefixInputGram}
    \leanok
    \uses{def:mpu_finite_cut_cap_grams, def:mpu_density_interval_grams}
    For arbitrary matrices $\rho,\tau$, one has
    \[
      \operatorname{Tr}(P_F(\rho)Q_G(\tau))
        =\operatorname{Tr}(U^\dagger U(\rho\otimes\tau)).
    \]
    Both cap Grams are positive semidefinite for positive semidefinite
    inputs. If $U^\dagger U=I$, their pairing equals
    $\operatorname{Tr}\rho\operatorname{Tr}\tau$.
    In particular, the maximally mixed suffix input normalizes every
    trace-one prefix input. The rectangular prefix formula agrees
    with the one-row prefix Gram above.
\end{theorem}
\begin{proof}\leanok
    Write $U^\dagger U$ as the sum of the rank-one Grams of its rows.
    The vectorization trace identity turns each product-density
    pairing into $\operatorname{Tr}((F_aG_c)^\dagger\rho^T
    F_aG_c\tau)$. Trace cyclicity and the finite sums give the
    displayed identity. Positivity follows from the matrix
    sandwiches. Substitute $U^\dagger U=I$ and take the trace of a
    tensor product to obtain the normalized formula.
\end{proof}

\begin{theorem}[Minimal spans supply positive definite cap witnesses]
    \label{thm:mpu_minimal_cap_gram_witnesses}
    \lean{MPUPrefixGram.prefixGram_normalizedIdentity_posDef_of_span_eq_top,
        MPUPrefixGram.suffixGram_normalizedIdentity_posDef_of_span_eq_top,
        MPUPrefixGram.exists_posDef_prefixInputGram_normalizer_of_isIsometry,
        MPUPrefixGram.exists_posDef_mem_densityPrefixGrams_of_span_eq_top,
        MPUPrefixGram.exists_posDef_mem_prefixGramAffineHull_of_span_eq_top}
    \leanok
    \uses{def:mpu_finite_cut_cap_grams, def:mpu_prefix_gram_affine_hull}
    Suppose the physical input spaces are nonzero. If the prefix rows
    span $\mathbb C^r$, their maximally mixed Gram is positive definite,
    providing a positive definite point in both $\mathcal S_F$ and
    $\mathcal C_F$. If the suffix columns span $\mathbb C^r$, their
    maximally mixed Gram is positive definite. Under global isometry,
    it provides a positive definite trace normalizer for every
    trace-one prefix input.
\end{theorem}
\begin{proof}\leanok
    \uses{thm:mpu_global_cap_gram_normalization}
    The unnormalized suffix Gram is the frame operator of its
    columns; the prefix Gram is the transpose of the frame operator
    of its rows. A finite frame operator is positive definite
    precisely when its vectors span. Scale by the reciprocal input
    dimension and apply the global normalization identity.
\end{proof}

\begin{definition}[The two cuts of a three-part factorization]
    \label{def:mpu_three_part_factorization}
    \lean{MPUCircuit.prefixIntervalProduct, MPUCircuit.intervalSuffixProduct}
    \leanok
    For a prefix $F$, middle interval $A$, and suffix $G$, define
    the longer prefix $FA$ and longer suffix $AG$ by ordinary
    contraction of their shared virtual indices. The resulting
    physical operator has coefficients $F_{ab}A_{cd}G_{ef}$.
\end{definition}

\begin{theorem}[Reassociation of the physical indices]
    \label{thm:mpu_factorization_reassociation}
    \lean{MPUCircuit.factorizedOperator_reassociate,
        MPUCircuit.factorizedOperator_isometry_reassociate}
    \leanok
    \uses{def:mpu_three_part_factorization, def:mpu_finite_cut_cap_grams}
    The operators obtained by contracting $(FA)G$ and $F(AG)$
    differ only by reassociation of their physical index triples.
    Global isometry is preserved by this reassociation.
\end{theorem}
\begin{proof}\leanok
    Expand the two finite virtual sums and apply associativity of
    scalar multiplication. Reindex the physical coordinates by
    the natural product associativity bijections.
\end{proof}

\begin{theorem}[An affine normalizer from global isometry]
    \label{thm:mpu_affine_normalizer_global_isometry}
    \lean{MPUCircuit.exists_posDef_prefixGramAffineHull_normalizer_of_isIsometry}
    \leanok
    \uses{def:mpu_finite_cut_cap_grams, def:mpu_prefix_gram_affine_hull}
    Suppose $U=FG$ is globally isometric, its suffix input space is
    nonzero, and the suffix columns span the bond space. Then a
    positive definite matrix $Q$ satisfies
    $\operatorname{Tr}(XQ)=1$ for every $X\in\mathcal C_F$.
\end{theorem}
\begin{proof}\leanok
    \uses{thm:mpu_minimal_cap_gram_witnesses,
        thm:mpu_real_affine_gram_identities}
    Take the maximally mixed suffix Gram and extend its physical
    trace normalization to the real affine hull.
\end{proof}

\begin{theorem}[Balanced interval isometries from global isometry]
    \label{thm:mpu_balanced_interval_global_isometry}
    \lean{MPUCircuit.exists_balanced_interval_isometry_of_factorizedOperator_isometry}
    \leanok
    \uses{def:mpu_three_part_factorization, def:mpu_prefix_gram_affine_hull,
        def:mpu_dual_gram_metric, def:mpu_vectorized_weighted_interval}
    Let $U=FAG$ be globally isometric, with nonzero physical input
    spaces and bond dimensions $r,s\geq1$. Suppose the rows of $F$
    and $FA$, and the columns of $AG$ and $G$, span their respective
    bond spaces. There are balanced positive metrics
    $P\in\mathcal C_F$ and $S\in\mathcal C_{FA}$, with their dual
    metrics normalizing the entire affine spaces and their
    inverse-metric traces equal to $r^2,s^2$. The middle interval
    with amplitudes $\sqrt P A_{cd}\sqrt{S^{-1}/s}$ is an isometry.
\end{theorem}
\begin{proof}\leanok
    \uses{thm:mpu_factorization_reassociation,
        thm:mpu_affine_normalizer_global_isometry,
        thm:mpu_minimal_cap_gram_witnesses,
        thm:mpu_two_cut_balanced_isometry}
    Apply global normalization and the spanning criterion at both
    cuts. These supply positive definite normalizers and positive
    definite points in both affine spaces. Apply the two-cut
    determinant-balancing theorem.
\end{proof}

\begin{definition}[Physical coefficient flattenings]
    \label{def:mpu_physical_coefficient_flattenings}
    \lean{MPUCircuit.coefficientFlattening,
        MPUCircuit.firstCutCoefficientFlattening,
        MPUCircuit.secondCutCoefficientFlattening}
    \leanok
    \uses{def:mpu_three_part_factorization}
    For $U=FG$, its cut coefficient matrix groups output and input
    indices on each side:
    $W_{(a,b),(e,f)}=(F_aG_e)_{bf}$.
    For a three-part factorization $FAG$, the first and second cut
    matrices are those of $F(AG)$ and $(FA)G$, respectively.
\end{definition}

\begin{theorem}[Minimal cut ranks give spanning caps]
    \label{thm:mpu_minimal_cut_ranks_spanning_caps}
    \lean{MPUCircuit.spans_eq_top_of_rank_mul_eq_card,
        MPUCircuit.coefficientFlattening_eq_mul,
        MPUCircuit.spans_eq_top_of_coefficientFlattening_rank_eq_card,
        MPUCircuit.secondCutCoefficientFlattening_apply,
        MPUCircuit.spans_eq_top_of_twoCutCoefficientFlattening_ranks}
    \leanok
    \uses{def:mpu_physical_coefficient_flattenings}
    If $X$ has $r$ columns, $Y$ has $r$ rows and
    $\operatorname{rank}(XY)=r$, then the rows of $X$ and columns
    of $Y$ span $\mathbb C^r$. The physical cut coefficient matrix
    factors through the virtual bond in precisely this way.
    Consequently, cut ranks equal to the two virtual dimensions
    imply all four spanning conditions for $F,AG,FA,G$.
    The two cut matrices contain the same physical coefficients,
    grouped at their respective cuts.
\end{theorem}
\begin{proof}\leanok
    Rank of a product is bounded by the rank of each factor, and
    each factor has rank at most $r$. Both ranks therefore equal
    $r$, giving the two spans. Apply this at both cuts.
    Matrix multiplication gives the coefficient factorization,
    and associativity gives the equality of regrouped coefficients.
\end{proof}

\begin{theorem}[Balanced interval isometries from minimal cut ranks]
    \label{thm:mpu_balanced_interval_minimal_cut_ranks}
    \lean{MPUCircuit.exists_balanced_interval_isometry_of_coefficientFlattening_ranks}
    \leanok
    \uses{def:mpu_physical_coefficient_flattenings,
        thm:mpu_balanced_interval_global_isometry}
    Let $FAG$ be globally isometric, with nonzero physical input
    spaces and virtual dimensions $r,s\geq1$. If its two bounding
    physical cut matrices have ranks $r,s$, then the balanced
    positive metrics, affine normalization identities, contraction
    traces $r^2,s^2$ and middle-interval isometry of
    Theorem~\ref{thm:mpu_balanced_interval_global_isometry} exist.
\end{theorem}
\begin{proof}\leanok
    \uses{thm:mpu_minimal_cut_ranks_spanning_caps,
        thm:mpu_balanced_interval_global_isometry}
    Derive the four spanning conditions from the two rank equalities,
    and apply the global-isometry result.
\end{proof}

\begin{theorem}[Simultaneous balanced metrics at all minimal cuts]
    \label{thm:mpu_simultaneous_balanced_cut_metrics}
    \lean{MPUCircuit.exists_balanced_prefix_metric_of_factorizedOperator_isometry,
        MPUCircuit.exists_simultaneous_balanced_prefix_metrics_of_factorizedOperator_isometries,
        MPUCircuit.exists_simultaneous_balanced_prefix_metrics_of_coefficientFlattening_ranks}
    \leanok
    \uses{thm:mpu_affine_normalizer_global_isometry,
        thm:mpu_minimal_cap_gram_witnesses,
        thm:mpu_affine_gram_balancing}
    Consider an indexed family of globally isometric two-part
    factorizations $F_jG_j$ with spanning prefix rows and suffix
    columns. The physical spaces and positive bond dimensions $r_j$
    may depend on $j$. There is one chosen metric $P_j$ at each index
    such that
    \[
        P_j\in\mathcal C_j,\quad P_j>0,\quad Q_j=P_j^{-1}/r_j>0,
        \qquad \operatorname{Tr}(XQ_j)=1\quad(X\in\mathcal C_j),
        \qquad \operatorname{Tr}(Q_j^{-1}P_j^{-1})=r_j^2.
    \]
    The same conclusion holds when each physical coefficient flattening
    has rank exactly $r_j$, since this implies both spanning conditions.
    All these identities hold for the same chosen family. In
    particular, the metrics may be chosen at every cut of a finite
    minimal chain before its interval isometries are assembled.
\end{theorem}
\begin{proof}\leanok
    Global isometry and minimality give a positive normalizer and
    a positive point at each cut. Apply affine determinant balancing
    separately, and choose one of its metric witnesses at each index.
\end{proof}

\begin{definition}[Scalar endpoint caps]
    \label{def:mpu_scalar_endpoint_caps}
    \lean{MPUCircuit.scalarPrefixOperator, MPUCircuit.emptyPrefixCap}
    \leanok
    A prefix with scalar virtual bonds has physical matrix
    $U_{ab}=F_{ab}(1,1)$. The empty prefix has one input, one output,
    and scalar coefficient one.
\end{definition}

\begin{theorem}[The endpoint affine spaces and metrics]
    \label{thm:mpu_endpoint_gram_normalization}
    \lean{MPUCircuit.intervalInputGram_scalar_bonds,
        MPUCircuit.prefixInputGram_eq_trace_smul_one_of_isIsometry,
        MPUCircuit.densityPrefixGrams_eq_singleton_one_of_isIsometry,
        MPUCircuit.prefixGramAffineHull_eq_singleton_one_of_isIsometry,
        MPUCircuit.emptyPrefixCap_isIsometry,
        MPUCircuit.emptyPrefixCap_prefixGramAffineHull,
        MPUCircuit.dualGramMetric_one_of_card_eq_one,
        MPUCircuit.dualGramMetric_one_unit}
    \leanok
    \uses{def:mpu_scalar_endpoint_caps, def:mpu_prefix_gram_affine_hull,
        def:mpu_dual_gram_metric}
    If the physical matrix of a scalar-bond prefix is an isometry,
    then $P_F(\rho)=\operatorname{Tr}(\rho)I$ for every complex
    input matrix $\rho$. For a nonzero input space, both its
    physical Gram set and real affine hull equal $\{I\}$. This
    also holds for the empty prefix. The one-dimensional identity
    is its own dual metric.
\end{theorem}
\begin{proof}\leanok
    \uses{thm:mpu_density_gram_separation}
    Closing the scalar bonds gives $K(I,I)=U^\dagger U$.
    Trace pairing therefore gives the scalar Gram identity for
    every $\rho$. A faithful density matrix witnesses the identity
    in the physical Gram set. Take its affine hull, and use the
    dimension-one dual-metric formula.
\end{proof}

\begin{definition}[Normalized joining contraction]
    \label{def:mpu_normalized_interval_merger}
    \lean{MPUCircuit.normalizedBalancedGramContraction,
        MPUCircuit.contractedWeightedIntervals}
    \leanok
    \uses{def:mpu_dual_gram_metric}
    For a positive metric $P$ on an $r$-dimensional bond, put
    $S=\sqrt P$, $T=\sqrt{P^{-1}/r}$ and
    $C=r^{-1}T^{-1}S^{-1}$.
    Given weighted coefficients $X_{ab}$ and $Y_{cd}$ on adjacent
    intervals, their contracted coefficient is $X_{ab}CY_{cd}$,
    vectorized over the two remaining virtual bonds.
\end{definition}

\begin{theorem}[The balanced merger has uniform success amplitude]
    \label{thm:mpu_balanced_merger_success_amplitude}
    \lean{MPUCircuit.normalizedBalancedGramContraction_vec_inner,
        MPUCircuit.weighted_interval_normalized_contraction,
        MPUCircuit.contractedWeightedIntervals_balanced,
        MPUCircuit.contractedWeightedIntervals_balanced_gram}
    \leanok
    \uses{def:mpu_normalized_interval_merger,
        def:mpu_vectorized_weighted_interval}
    The joining contraction $C$ has squared Hilbert--Schmidt norm one.
    For arbitrary outside weights $L,R$ and adjacent coefficients
    $A_{ab},B_{cd}$,
    \[
        (LA_{ab}T)C(SB_{cd}R)=r^{-1}LA_{ab}B_{cd}R.
    \]
    Thus contraction gives $r^{-1}$ times the parent interval map.
    If that parent map is an isometry, the contracted map has input
    Gram $r^{-2}I$. The success probability is therefore the same
    for every input, including entangled inputs.
\end{theorem}
\begin{proof}\leanok
    \uses{thm:mpu_balanced_gram_contraction}
    The unnormalized contraction has squared Hilbert--Schmidt norm
    $r^2$. Divide by $r$, cancel the two invertible square roots,
    and take the input Gram of the resulting scalar multiple.
\end{proof}

\begin{definition}[Two-reflection subspace amplification]
    \label{def:mpu_exact_subspace_amplification}
    \lean{Matrix.subspaceReflection, Matrix.amplificationPlane,
        Matrix.subspaceAmplificationStep}
    \leanok
    For a projection $P$, put $R_P=I-2P$. For two maps $G,H$
    with common input and output spaces, put
    \[
        V_\theta=\sin\theta\,G+\cos\theta\,H,
        \qquad Q_\theta=-R_{V_\theta V_\theta^\dagger}R_P.
    \]
\end{definition}

\begin{theorem}[Unitary reflections and isometric planes]
    \label{thm:mpu_amplification_reflections}
    \lean{Matrix.subspaceReflection_mem_unitaryGroup,
        Matrix.amplificationPlane_conjTranspose_mul_self,
        Matrix.isStarProjection_mul_conjTranspose_of_isometry,
        Matrix.subspaceAmplificationStep_mem_unitaryGroup}
    \leanok
    \uses{def:mpu_exact_subspace_amplification}
    An orthogonal projection gives a unitary reflection.
    If $G^\dagger G=H^\dagger H=I$ and $G^\dagger H=0$, then
    $V_\theta$ is an isometry for every real $\theta$ and
    $V_\theta V_\theta^\dagger$ is an orthogonal projection.
    If $P$ is also an orthogonal projection, $Q_\theta$ is unitary
    on the entire output space.
\end{theorem}
\begin{proof}\leanok
    Square $I-2P$, use the orthogonal Gram identities, and use
    $V_\theta^\dagger V_\theta=I$ to square its image projection.
    A product of the two reflections, with scalar phase minus one,
    is unitary.
\end{proof}

\begin{theorem}[Exact rotation of an isometric subspace]
    \label{thm:mpu_exact_subspace_amplification}
    \lean{Matrix.subspaceAmplificationStep_mul_plane,
        Matrix.subspaceAmplificationStep_pow_mul_plane,
        Matrix.subspaceAmplificationStep_pow_eq_good}
    \leanok
    \uses{def:mpu_exact_subspace_amplification}
    Suppose $G^\dagger G=H^\dagger H=I$, $G^\dagger H=0$,
    $PG=G$ and $PH=0$. For all real $\theta,\phi$ and $k\geq0$,
    \[
        Q_\theta V_\phi=V_{\phi+2\theta},
        \qquad Q_\theta^k V_\phi=V_{\phi+2k\theta}.
    \]
    In particular, for $\ell\geq0$ and
    $\theta=\pi/(4\ell+2)$,
    $Q_\theta^\ell V_\theta=G$ exactly, with no residual phase.
\end{theorem}
\begin{proof}\leanok
    Expand the two reflections on the orthogonal images of $G,H$.
    The trigonometric addition identities give the one-step rotation.
    Iterate it by induction. The final angle is
    $(2\ell+1)\pi/(4\ell+2)=\pi/2$.
\end{proof}

\begin{definition}[Attenuation by one binary flag]
    \label{def:mpu_uniform_success_attenuation}
    \lean{Matrix.attenuationRotation, Matrix.attenuatedIsometry,
        Matrix.attenuatedSuccessProjection}
    \leanok
    For a real coefficient $t$, put $s=\sqrt{1-t^2}$ and define
    \[
        R_t=\begin{pmatrix}t&-s\\s&t\end{pmatrix},
        \qquad V_t=\begin{pmatrix}tV\\sV\end{pmatrix},
        \qquad P_t=\begin{pmatrix}P&0\\0&0\end{pmatrix}.
    \]
    The two row blocks represent the two values of a new binary flag.
\end{definition}

\begin{theorem}[Exact attenuation of uniform success]
    \label{thm:mpu_uniform_success_attenuation}
    \lean{Matrix.attenuationRotation_mem_unitaryGroup,
        Matrix.attenuatedIsometry_conjTranspose_mul_self,
        Matrix.attenuatedSuccessProjection_isStarProjection,
        Matrix.attenuatedIsometry_success_gram}
    \leanok
    \uses{def:mpu_uniform_success_attenuation}
    If $t^2\leq1$, then $R_t$ is unitary on both flag values.
    If $V^\dagger V=I$, also $V_t^\dagger V_t=I$.
    An orthogonal projection $P$ gives an orthogonal projection $P_t$.
    Moreover, for any uniform probability identity
    $V^\dagger PV=p^2I$,
    \[
        V_t^\dagger P_tV_t=(tp)^2I.
    \]
    Thus a nonnegative attenuation coefficient multiplies the
    original nonnegative success amplitude by $t$.
\end{theorem}
\begin{proof}\leanok
    The two flag coefficients satisfy $t^2+s^2=1$. Multiply the
    two-by-two rotation and the block matrices, using this identity
    for normalization and the first row block for the success Gram.
\end{proof}

\begin{definition}[Uniform success and failure maps]
    \label{def:mpu_uniform_postselection}
    \lean{Matrix.postselectionSuccess, Matrix.postselectionFailure,
        Matrix.postselectionAmplificationStep}
    \leanok
    \uses{def:mpu_exact_subspace_amplification}
    Given an isometry $V$, success projection $P$ and real
    success amplitude $0<p<1$, put
    \[
        G=p^{-1}PV,\qquad H=(1-p^2)^{-1/2}(I-P)V,
        \qquad Q=-R_{VV^\dagger}R_P.
    \]
\end{definition}

\begin{theorem}[Decomposition from uniform success probability]
    \label{thm:mpu_uniform_postselection_decomposition}
    \lean{Matrix.uniformPostselection_decomposition}
    \leanok
    \uses{def:mpu_uniform_postselection}
    Suppose $V^\dagger V=I$, $P$ is an orthogonal projection,
    $0<p<1$ and $V^\dagger PV=p^2I$. The maps $G,H$ above satisfy
    \[
        G^\dagger G=H^\dagger H=I,\qquad G^\dagger H=0,
        \qquad PG=G,\quad PH=0,
        \qquad pG+\sqrt{1-p^2}\,H=V.
    \]
    Thus both isometric subspaces follow from the uniform probability
    identity.
\end{theorem}
\begin{proof}\leanok
    Projection identities give the two projected Grams $p^2I$ and
    $(1-p^2)I$, and their mixed Gram is zero. Divide by their
    positive amplitudes. Adding the projected parts recovers $V$.
\end{proof}

\begin{theorem}[Exact amplification from the initial isometry]
    \label{thm:mpu_uniform_postselection_amplification}
    \lean{Matrix.postselectionAmplificationStep_mem_unitaryGroup,
        Matrix.amplificationPlane_postselection_eq_self,
        Matrix.postselectionAmplificationStep_pow_mul,
        Matrix.postselectionAmplificationStep_pow_eq_success,
        Matrix.projection_mul_eq_self_of_uniform_probability_one}
    \leanok
    \uses{def:mpu_uniform_postselection,
        def:mpu_exact_subspace_amplification}
    Suppose $V^\dagger V=I$ and $P$ is an orthogonal projection.
    The operator $Q$ is unitary on the entire output space.
    If $V^\dagger PV=\sin^2\theta\,I$ with $0<\theta<\pi/2$,
    then the derived success and failure maps give
    \[
        V=V_\theta,\qquad Q^kV=V_{\theta+2k\theta}.
    \]
    For $\theta=\pi/(4\ell+2)$ and $\ell\geq1$,
    $Q^\ell V=(\sin\theta)^{-1}PV$ exactly.
    If instead $V^\dagger PV=I$, then already $PV=V$.
\end{theorem}
\begin{proof}\leanok
    \uses{thm:mpu_amplification_reflections,
        thm:mpu_exact_subspace_amplification,
        thm:mpu_uniform_postselection_decomposition}
    Derive the orthogonal isometric decomposition. In the first
    quadrant, $\sqrt{1-\sin^2\theta}=\cos\theta$, so it is the
    initial amplification plane. Apply the rotation identities.
    At probability one the complementary projected Gram is zero,
    forcing the complementary map to vanish.
\end{proof}

\begin{definition}[Amplification rounds from the bond dimension]
    \label{def:mpu_exact_amplification_budget}
    \lean{MPUCircuit.amplificationRounds, MPUCircuit.amplificationAngle}
    \leanok
    For $r\geq0$, define
    \[
        \ell(r)=\lceil\pi r/4\rceil,
        \qquad \theta(r)=\frac{\pi}{4\ell(r)+2}.
    \]
\end{definition}

\begin{theorem}[A dimension bound for exact amplification]
    \label{thm:mpu_exact_amplification_budget}
    \lean{MPUCircuit.amplificationRounds_pos,
        MPUCircuit.amplificationRounds_le,
        MPUCircuit.amplificationAngle_pos,
        MPUCircuit.amplificationAngle_lt_pi_div_two,
        MPUCircuit.amplificationAngle_le_inv,
        MPUCircuit.amplificationAngle_sin_pos_le_inv}
    \leanok
    \uses{def:mpu_exact_amplification_budget}
    For $r\geq1$,
    \[
        1\leq\ell(r)\leq r,\qquad
        0<\theta(r)<\pi/2,\qquad
        \theta(r)\leq r^{-1},\qquad
        0<\sin\theta(r)\leq r^{-1}.
    \]
    Thus the merger success amplitude $1/r$ can be attenuated to
    the prescribed exact angle, with at most $r$ amplification rounds.
\end{theorem}
\begin{proof}\leanok
    Positivity and the ceiling inequalities give the lower bound on
    $\ell(r)$. Since $\pi\leq4$, its ceiling is at most $r$.
    Also $4\ell(r)+2\geq\pi r+2$, which bounds the angle by $1/r$.
    The angle is in $(0,\pi/2)$, and $\sin\theta\leq\theta$.
\end{proof}

\begin{definition}[Physical cut spaces]
    \label{def:mpu_minimal_physical_cut_spaces}
    \lean{MPSPreparation.cutCoefficientMatrix,
        MPSPreparation.CutPrefixConfig, MPSPreparation.CutSuffixConfig,
        MPSPreparation.extendCutPrefix, MPSPreparation.extendCutSuffix,
        MPSPreparation.cutPrefixRestriction, MPSPreparation.cutSuffixRestriction,
        MPSPreparation.cutColumnSpace, MPSPreparation.cutRank,
        MPSPreparation.cutSliceLinearMap, MPSPreparation.cutSlice,
        MPSPreparation.fullCutVector}
    \leanok
    For a coefficient tensor $\psi$ on a finite chain, let $C_k$ be
    its coefficient matrix with rows indexed by configurations before
    cut $k$ and columns indexed by the remaining configurations.
    Put $E_k=\operatorname{span}\{\text{columns of }C_k\}$ and
    $r_k=\operatorname{rank}C_k$. Fixing the letter at site $k$
    defines a coordinate slice from functions on the longer prefix
    to functions on the shorter prefix. The last-cut vector is $\psi$
    regarded as a function on the full prefix.
\end{definition}

\begin{theorem}[Slicing and ranks of physical cuts]
    \label{thm:mpu_minimal_physical_cut_spaces}
    \lean{MPSPreparation.cutCoefficientMatrix_adjacent,
        MPSPreparation.cutCoefficientMatrix_restrictions,
        MPSPreparation.extendCutPrefix_restriction,
        MPSPreparation.finrank_cutColumnSpace,
        MPSPreparation.cutSliceLinearMap_maps_columnSpace,
        MPSPreparation.cutSlice_apply_restriction,
        MPSPreparation.cutRank_pos,
        MPSPreparation.card_cutPrefixConfig_zero,
        MPSPreparation.card_cutSuffixConfig_last,
        MPSPreparation.cutRank_zero, MPSPreparation.cutRank_last,
        MPSPreparation.cutColumnSpace_zero_eq_top,
        MPSPreparation.cutCoefficientMatrix_last,
        MPSPreparation.fullCutVector_mem,
        MPSPreparation.fullCutVector_restriction,
        MPSPreparation.fullCutVector_ne_zero}
    \leanok
    \uses{def:mpu_minimal_physical_cut_spaces}
    Adjoining a letter on the prefix or suffix side gives the same
    coefficient, and joining the two restrictions of a configuration
    recovers $\psi$. Each coordinate slice maps $E_{k+1}$ into $E_k$,
    and $\dim E_k=r_k$. If $\psi\ne0$, every $r_k$ is positive;
    at the two endpoint cuts it is one. The initial cut space is the
    full one-dimensional space, and the last-cut vector is a nonzero
    member of the last cut space.
\end{theorem}
\begin{proof}\leanok
    The adjacent-cut identity follows by separating sites before,
    at, and after the new letter. Slices take coefficient columns
    to coefficient columns, and hence preserve their spans.
    A nonzero coefficient gives a nonzero entry at every cut.
    Each endpoint matrix has one row or one column.
\end{proof}

\begin{definition}[The chain of cut-coordinate maps]
    \label{def:mpu_minimal_cut_coordinate_chain}
    \lean{MPSPreparation.cutSiteMatrix,
        MPSPreparation.cutColumnBasis,
        MPSPreparation.cutEvaluationMap, MPSPreparation.cutEvaluationMatrix,
        MPSPreparation.cutPrefixChain, MPSPreparation.minimalCutChain}
    \leanok
    Choose a basis of each $E_k$ indexed by $r_k$ coordinates.
    The site matrix for a physical letter is the matrix of its slice
    $E_{k+1}\to E_k$. Evaluation at a prefix configuration is a
    linear functional on $E_k$. Padding the site matrices gives
    their products in a common matrix space, while their rectangular
    forms give an open-boundary chain of bond dimensions $r_k$.
\end{definition}

\begin{theorem}[Reconstruction in cut coordinates]
    \label{thm:mpu_minimal_cut_coordinate_reconstruction}
    \lean{MPSPreparation.exists_cutBasisFamily_with_endpoints,
        MPSPreparation.cutEvaluationMatrix_apply,
        MPSPreparation.cutEvaluationMap_step,
        MPSPreparation.cutEvaluationMatrix_step,
        MPSPreparation.cutEvaluationMatrix_eq_initial_mul_eval,
        MPSPreparation.initialCutEvaluation_row,
        MPSPreparation.initialCutEvaluation_mul_apply,
        MPSPreparation.cutPrefixChain_eval_entry,
        MPSPreparation.minimalCutChain_coeff}
    \leanok
    \uses{def:mpu_minimal_cut_coordinate_chain,
        thm:mpu_padded_bond_products}
    For a nonzero tensor of positive length, the endpoint bases can
    be chosen as $1$ and $\psi$. Successive slice matrices reconstruct
    prefix evaluation. Their full product between the zero endpoint
    coordinates is then exactly $\psi$ at the given configuration,
    provided each intermediate rank fits the padding dimension.
    Thus the resulting rectangular open-boundary chain has exactly
    the original coefficients, including their normalization and phase.
\end{theorem}
\begin{proof}\leanok
    Fix the two one-dimensional endpoint bases and choose bases
    at the other cuts. Composition of the evaluation maps gives
    the product formula. At the initial cut evaluation is the row
    selecting the zero coordinate; at the last cut its selected
    column evaluates $\psi$.
\end{proof}

\begin{theorem}[Existence of a minimal finite-chain representation]
    \label{thm:mpu_minimal_finite_chain_representation}
    \lean{MPSPreparation.exists_obcChainTensor_coeff_eq_bondDim_eq_cutRank}
    \leanok
    \uses{thm:mpu_minimal_cut_coordinate_reconstruction}
    Every nonzero coefficient tensor on $N\geq1$ sites whose
    consecutive physical cut ranks satisfy $r_k\leq D$ has an
    exact open-boundary matrix-product representation with bond
    dimension $r_k$ at each cut.
\end{theorem}
\begin{proof}\leanok
    Use the chain of coordinate slices with the endpoint bases fixed
    above. Its bond dimensions are the physical cut ranks by construction,
    and reconstruction gives every coefficient.
\end{proof}

\begin{definition}[Prefix and suffix factors at a physical cut]
    \label{def:mpu_minimal_cut_factors}
    \lean{MPSPreparation.cutPrefixMatrix, MPSPreparation.cutSuffixMatrix}
    \leanok
    \uses{def:mpu_minimal_cut_coordinate_chain}
    Let $F_k$ evaluate the chosen cut basis on each prefix configuration,
    and let $G_k$ contain the coordinates of each coefficient column
    in that basis.
\end{definition}

\begin{theorem}[Full-rank factorization at every physical cut]
    \label{thm:mpu_minimal_cut_factors}
    \lean{MPSPreparation.cutCoefficientMatrix_eq_prefix_mul_suffix,
        MPSPreparation.cutPrefixMatrix_rank,
        MPSPreparation.cutSuffixMatrix_rank,
        MPSPreparation.cutPrefixMatrix_rows_and_cutSuffixMatrix_cols_span,
        MPSPreparation.cutPrefixMatrix_step,
        MPSPreparation.cutSlice_column,
        MPSPreparation.cutSuffixMatrix_step,
        MPSPreparation.minimalCutChain_cutCoefficientMatrix_factorization,
        MPSPreparation.cutPrefixMatrix_zero,
        MPSPreparation.cutSuffixMatrix_last,
        MPSPreparation.exists_obcChainTensor_minimal_factorizations}
    \leanok
    \uses{def:mpu_minimal_cut_factors,
        thm:mpu_minimal_finite_chain_representation}
    At every cut, $C_k=F_kG_k$ and
    $\operatorname{rank}F_k=\operatorname{rank}G_k=r_k$.
    The prefix rows and suffix columns both span the bond space.
    Fixing the next letter gives the corresponding multiplication
    by its site matrix. The normalized endpoint factors are one.
    All these factors and spanning properties can be chosen together
    with the exact minimal open-boundary representation.
\end{theorem}
\begin{proof}\leanok
    Basis reconstruction gives the factorization. The basis evaluation
    map is injective, while coefficient columns span the cut space;
    both ranks are therefore $r_k$. Composition with a coordinate
    slice gives the site identities and the endpoint normalizations.
\end{proof}

\begin{definition}[A reversible prefix zero test]
    \label{def:mpu_reversible_zero_conjunction}
    \lean{MPSPreparation.zeroPrefixTest, MPSPreparation.zeroPrefixBit,
        MPSPreparation.zeroConjunctionPrev,
        MPSPreparation.zeroConjunctionIncrement,
        MPSPreparation.zeroConjunctionStep,
        MPSPreparation.zeroConjunctionPerm,
        MPSPreparation.zeroConjunctionState,
        MPSPreparation.zeroConjunctionStepSites,
        MPSPreparation.zeroConjunctionStepGateCount}
    \leanok
    For $a$ data qudits, use $a$ temporary qudits initialized in zero.
    At step $j$, increment temporary coordinate $j$ modulo $d$ when
    data coordinate $j$ is zero and either $j=0$ or the preceding
    temporary coordinate is one. This is a reversible permutation
    on the full register. Put $c_d=38(d^3)^6+1$; a step has gate
    budget $4ac_d$ after routing.
\end{definition}

\begin{theorem}[Conjunction on the entire logical space]
    \label{thm:mpu_reversible_zero_conjunction}
    \lean{MPSPreparation.zeroConjunctionPerm_apply_clean,
        MPSPreparation.zeroConjunctionIncrement_update,
        MPSPreparation.zeroPrefixBit_eq_one,
        MPSPreparation.zeroConjunctionStep_state,
        MPSPreparation.zeroConjunctionStep_isLocalPerm,
        MPSPreparation.isPairProduct_zeroConjunctionStep,
        MPSPreparation.isPairProduct_zeroConjunctionPerm}
    \leanok
    \uses{def:mpu_reversible_zero_conjunction,
        thm:mpu_quantitative_unitary_synthesis,
        thm:mpu_arbitrary_circuit_placement}
    For $d\geq2$ and every logical configuration $x$, the conjunction
    keeps $x$ and writes in temporary coordinate $j$ whether its
    first $j+1$ data coordinates are all zero. Each step acts on
    at most three sites. The full conjunction has an actual
    neighboring-pair circuit of at most $4c_da^2$ gates.
\end{theorem}
\begin{proof}\leanok
    Induct on the completed prefix. At the next step its temporary
    coordinate is still zero, and the previous Boolean value supplies
    the earlier prefix test. The controlled increment is independent
    of its target and supported on at most three sites. Synthesize
    that finite permutation and concatenate the routed steps.
\end{proof}

\begin{definition}[Reflection using a temporary zero-test register]
    \label{def:mpu_clean_zero_register_reflection}
    \lean{MPSPreparation.zeroRegisterReflection,
        MPSPreparation.zeroScratchPhaseLocal,
        MPSPreparation.zeroScratchPhase,
        MPSPreparation.zeroRegisterReflectionGateCount}
    \leanok
    \uses{def:mpu_reversible_zero_conjunction}
    For $a\geq1$, compute the prefix conjunction, change the phase
    by minus one when the final temporary coordinate is one, and
    undo the conjunction. Its gate budget is
    $K(d,a)=8c_da^2+4a$.
\end{definition}

\begin{theorem}[Exact initialized-register reflections]
    \label{thm:mpu_clean_zero_register_reflection}
    \lean{MPSPreparation.zeroRegisterReflection_mul_cleanEmbedding,
        MPSPreparation.zeroScratchPhaseLocal_mem_unitary,
        MPSPreparation.isPairProduct_zeroScratchPhase,
        MPSPreparation.isPairProduct_zeroRegisterReflection,
        MPSPreparation.zeroRegisterReflection_mem_unitary,
        MPSPreparation.zeroRegisterReflection_eq_diagonal,
        MPSPreparation.zeroRegisterReflection_apply_clean,
        MPSPreparation.zeroRegisterReflectionGateCount_eq,
        MPSPreparation.isPairProduct_embedOp_zeroRegisterReflection,
        MPSPreparation.embedOp_zeroRegisterReflection_apply_clean,
        MPSPreparation.isPairProduct_neg_one}
    \leanok
    \uses{def:mpu_clean_zero_register_reflection,
        thm:mpu_reversible_zero_conjunction}
    The constructed circuit is unitary on the full register and,
    for every logical input, returns the temporary register to zero
    while applying $I-2|0\rangle\langle0|$ to the tested register.
    This is an exact equality of linear maps and costs at most
    $K(d,a)$ neighboring-pair gates.
    Injective placement on an $n$-site chain costs at most $2nK(d,a)$;
    all other logical sites are unrestricted. For an empty tested
    register the reflection is the global phase $-I$, implemented
    by one neighboring-pair gate when $n\geq2$.
\end{theorem}
\begin{proof}\leanok
    On a clean temporary register the final Boolean value tests
    exactly whether the entire data register is zero. The diagonal
    phase changes no coordinate, so the inverse permutation erases
    every temporary value. Linear extension gives the complete
    initialized-subspace identity. The costs add the forward and
    inverse conjunctions and the one-site phase; arbitrary placement
    uses swap routing.
\end{proof}

\begin{theorem}[Products in padded bond coordinates]
    \label{thm:mpu_padded_bond_products}
    \lean{Matrix.zeroPad_mul, Matrix.zeroPad_mulVec}
    \leanok
    Extending rectangular bond matrices and vectors by zero preserves
    their products whenever the intermediate bond dimension is no
    larger than the common padding dimension. The outer dimensions
    need not satisfy this bound.
\end{theorem}
\begin{proof}\leanok
    In the product sum, terms outside the intermediate bond space
    vanish. The remaining terms are exactly the rectangular product.
\end{proof}

\begin{definition}[A two-copy dilation]
    \label{def:mpu_partial_isometry_dilation}
    \lean{Matrix.partialIsometryDilation}
    \leanok
    For a square matrix $A$, define
    \[
      \mathcal D_A=
      \begin{pmatrix}
        A&I-AA^\dagger\\
        I-A^\dagger A&-A^\dagger
      \end{pmatrix}.
    \]
\end{definition}

\begin{theorem}[Dilation of a partial isometry]
    \label{thm:mpu_partial_isometry_dilation}
    \lean{Matrix.partialIsometryDilation_mem_unitaryGroup,
        Matrix.partialIsometryDilation_mul_fromRows,
        Matrix.partialIsometryDilation_initial_gram,
        Matrix.partialIsometryDilation_success_gram}
    \leanok
    \uses{def:mpu_partial_isometry_dilation}
    If $AA^\dagger A=A$, then $\mathcal D_A$ is unitary. For every
    matrix of input columns $V$,
    \[
      \mathcal D_A\binom V0=
        \binom{AV}{(I-A^\dagger A)V}.
    \]
    The total output Gram is $V^\dagger V$, and projection onto
    the first output copy has Gram $(AV)^\dagger AV$.
    The first-copy identity holds for arbitrary $A$.
\end{theorem}
\begin{proof}\leanok
    Both $A^\dagger A$ and $AA^\dagger$ are projections.
    Block multiplication gives $\mathcal D_A^\dagger\mathcal D_A=I$,
    the displayed action, and its projected Gram.
\end{proof}

\begin{theorem}[A normalized row followed by a reset]
    \label{thm:mpu_normalized_row_reset}
    \lean{Matrix.vecMulVec_single_mul_conjTranspose_mul}
    \leanok
    If $\sum_x\overline{w_x}w_x=1$, the matrix
    $A=|z\rangle\sum_xw_x\langle x|$ satisfies
    $AA^\dagger A=A$.
\end{theorem}
\begin{proof}\leanok
    Multiply the two outer products and use the unit norm of the row.
\end{proof}

\begin{definition}[The normalized bond reset and its dilation]
    \label{def:mpu_normalized_bond_dilation}
    \lean{MPUCircuit.normalizedBondReset, MPUCircuit.normalizedBondDilation}
    \leanok
    \uses{def:mpu_partial_isometry_dilation}
    Let $w$ be the vectorization of the normalized balanced joining
    contraction. The bond reset sends this row functional to a
    specified pair of bond basis states. Its one-flag dilation
    is the two-copy dilation of that reset matrix.
\end{definition}

\begin{theorem}[The explicit unitary bond dilation]
    \label{thm:mpu_normalized_bond_dilation}
    \lean{MPUCircuit.normalizedBondDilation_mem_unitaryGroup,
        MPUCircuit.normalizedBondReset_mul_conjTranspose_mul,
        MPUCircuit.normalizedBondDilation_mul_fromRows,
        MPUCircuit.normalizedBondDilation_first_branch}
    \leanok
    \uses{def:mpu_normalized_bond_dilation,
        thm:mpu_normalized_row_reset,
        thm:mpu_partial_isometry_dilation}
    A positive definite balanced metric gives a unitary bond dilation.
    On a clean flag input, its first output branch applies exactly
    the normalized contraction and resets the bond pair to the
    specified basis state. Its second branch records the complementary
    initial component.
\end{theorem}
\begin{proof}\leanok
    The normalized contraction has vectorized norm one, so the bond
    reset is a partial isometry. Apply its explicit dilation and the
    first-copy action formula.
\end{proof}

\begin{definition}[Quantitative finite-register synthesis]
    \label{def:mpu_quantitative_unitary_synthesis}
    \lean{MPSPreparation.unitaryGateCount}
    \leanok
    Set
    \[
      G(d,n)=2n+3(d^n)^2(2n+1)(4n3^n).
    \]
\end{definition}

\begin{theorem}[Exact synthesis with a dimension bound]
    \label{thm:mpu_quantitative_unitary_synthesis}
    \lean{MPSPreparation.isPairProduct_unitaryGateCount,
        MPSPreparation.isPairProduct_ctrlOp_le_exponential,
        MPSPreparation.isPairProduct_twoLevel_le_exponential,
        MPSPreparation.unitaryGateCount_le_hilbertDimension_pow,
        MPSPreparation.isPairProduct_le_hilbertDimension_pow,
        MPSPreparation.exists_isPairProduct_apply_embedding_eq}
    \leanok
    \uses{def:mpu_quantitative_unitary_synthesis}
    Every unitary on $n\geq2$ qudits of dimension $d\geq2$
    has an exact neighboring-pair circuit of at most
    $G(d,n)\leq38(d^n)^6$ gates, including its global phase.
    An isometry together with an injective placement of its input
    basis extends to such a circuit with the prescribed columns.
    A two-level gate with $k$ controls has cost at most $4n3^k$;
    a two-level gate between arbitrary configurations has cost at
    most $(2n+1)(4n3^n)$.
\end{theorem}
\begin{proof}\leanok
    Eliminate matrix entries by two-level unitaries, express each
    by controlled one-site gates, and route its two-site factors.
    A separate one-site phase preserves exact matrix equality.
    Extension of orthonormal columns gives the isometry assertion.
\end{proof}

\begin{theorem}[Arbitrary placement of a circuit]
    \label{thm:mpu_arbitrary_circuit_placement}
    \lean{MPSPreparation.IsPairProduct.embedOp_injective}
    \leanok
    On a chain of $n$ sites with positive local dimension,
    injectively placing a neighboring-pair circuit of $K$ gates
    requires at most $2nK$ neighboring-pair gates.
\end{theorem}
\begin{proof}\leanok
    Each placed gate acts on two distinct sites. Swaps bring them
    together and restore their positions at cost at most $2n$.
\end{proof}

\begin{definition}[The attenuation of an inverse-rank amplitude]
    \label{def:mpu_exact_uniform_merging}
    \lean{MPUCircuit.mergingAttenuation}
    \leanok
    \uses{def:mpu_exact_amplification_budget}
    Set $t_r=r\sin\theta(r)$.
\end{definition}

\begin{theorem}[Exact merging at inverse-rank probability]
    \label{thm:mpu_exact_uniform_merging}
    \lean{Matrix.postselectionAmplificationStep_pow_eq_first_branch,
        MPUCircuit.mergingAttenuation_pos_le_one,
        Matrix.postselectionSuccess_attenuatedIsometry_eq_first_branch,
        Matrix.exists_exact_uniform_merging}
    \leanok
    \uses{def:mpu_exact_uniform_merging,
        thm:mpu_uniform_success_attenuation,
        thm:mpu_uniform_postselection_amplification}
    Let $V^\dagger V=I$, let $P$ be an orthogonal projection,
    and suppose $V^\dagger PV=r^{-2}I$, with $r\geq1$.
    Attenuation by $0<t_r\leq1$ and $\ell(r)\leq r$
    reflection rounds give exactly
    \[
      \binom{rPV}{0}.
    \]
    The equality includes the phase and holds on all coherent inputs.
\end{theorem}
\begin{proof}\leanok
    The attenuated amplitude is $t_r/r=\sin\theta(r)$.
    Exact amplification gives its normalized successful part.
    The two scalar normalizations cancel to leave $rPV$.
\end{proof}

\begin{theorem}[Circuits for image reflections]
    \label{thm:mpu_image_reflection_circuits}
    \lean{Matrix.subspaceReflection_image,
        MPSPreparation.IsPairProduct.subspaceReflection_image,
        MPSPreparation.IsPairProduct.pow,
        MPSPreparation.IsPairProduct.postselectionAmplificationStep}
    \leanok
    If a unitary $U$ has a $K$-gate circuit and reflection about
    $EE^\dagger$ has an $L$-gate circuit, then reflection about
    $(UE)(UE)^\dagger$ has a $(2K+L)$-gate circuit, since
    \[
      I-2(UE)(UE)^\dagger=U(I-2EE^\dagger)U^\dagger.
    \]
    A $K$-gate circuit repeated $k$ times costs at most $kK$.
    On $n\geq2$ sites, an amplification round with reflection
    costs $K,L$ costs at most $2n+K+L$, including its phase.
\end{theorem}
\begin{proof}\leanok
    Unitary conjugation gives the identity. Concatenate the forward
    circuit, the reflection, and the inverse circuit. Repetition
    concatenates copies; the prescribed global phase costs at most $2n$.
\end{proof}

\begin{theorem}[Balanced recursion with polynomial overhead]
    \label{thm:mpu_polynomial_overhead_recursion}
    \lean{MPUCircuit.cost_le_pow_of_geometric_overhead,
        MPUCircuit.cost_le_polynomial_of_geometric_overhead,
        MPUCircuit.cost_le_polynomial_clog_of_geometric_overhead}
    \leanok
    Suppose $C,q\geq0$, $r\geq q+1$, $T(0)\leq a$, and
    $T(j+1)\leq rT(j)+Cq^j$. Then
    \[
      T(L)\leq(a+C)r^L-Cq^L.
    \]
    If also $a\geq0$ and $r\leq2^k$, this is at most
    $(a+C)(2^L)^k$. In particular, polynomial overhead in the
    interval length is compatible with polynomial total cost
    at a fixed child-call factor. For integer $r$, one may take
    $k=\lceil\log_2 r\rceil$.
\end{theorem}
\begin{proof}\leanok
    Induction uses $(r-q-1)Cq^j\geq0$ at each step.
    The binary upper bound on $r$ gives the polynomial estimate.
\end{proof}

\begin{definition}[Paired coefficients of a finite operator]
    \label{def:mpu_paired_operator_coefficients}
    \lean{MPUCircuit.operatorCoefficientTensor, MPUCircuit.pairPhysicalConfig,
        MPUCircuit.physicalPairConfigEquiv, MPUCircuit.physicalCutConfigEquiv}
    \leanok
    Pair each output letter with its input letter, giving a coefficient
    tensor of local dimension $d^2$. Splitting physical configurations
    at a cut and separating their output-input pairs are bijections.
    The resulting coefficient cut rank is the consecutive operator
    Schmidt rank.
\end{definition}

\begin{theorem}[Minimal physical cuts of a finite unitary]
    \label{thm:mpu_minimal_physical_unitary_cuts}
    \lean{MPUCircuit.exists_minimal_obcChainTensor_of_unitary,
        MPUCircuit.operatorCoefficientTensor_apply,
        MPUCircuit.operatorCoefficientTensor_ne_zero,
        MPUCircuit.physicalCutConfigEquiv_pair}
    \leanok
    \uses{def:mpu_paired_operator_coefficients,
        thm:mpu_minimal_finite_chain_representation}
    Pairing and splitting commute, and the paired tensor recovers
    every original operator entry. A finite unitary on a positive
    physical alphabet has a nonzero coefficient tensor. If its
    consecutive operator Schmidt ranks are at most $D$ and its
    length is positive, it therefore has an exact open-boundary
    representation with those ranks as its bond dimensions.
\end{theorem}
\begin{proof}\leanok
    The configuration bijections act site by site. A zero paired
    tensor would make the entire unitary matrix zero, contradicting
    its unit Gram. Apply the minimal representation theorem.
\end{proof}

\begin{definition}[Operator factors in minimal cut bases]
    \label{def:mpu_minimal_operator_factors}
    \lean{MPUCircuit.minimalOperatorPrefixFactor,
        MPUCircuit.minimalOperatorSuffixFactor}
    \leanok
    Separate the output and input configurations in the prefix and
    suffix matrices of the paired coefficient tensor. The resulting
    factors are operator-valued prefix rows and suffix columns.
\end{definition}

\begin{theorem}[Simultaneous metrics from the finite unitary]
    \label{thm:mpu_metrics_from_finite_unitary}
    \lean{MPUCircuit.exists_cutBases_simultaneous_balanced_metrics_of_unitary,
        MPUCircuit.factorizedOperator_minimalOperatorFactors,
        MPUCircuit.minimalOperatorFactors_isIsometry,
        MPUCircuit.minimalOperatorFactors_spans,
        MPUCircuit.exists_simultaneous_balanced_minimalOperator_metrics}
    \leanok
    \uses{def:mpu_minimal_operator_factors,
        thm:mpu_minimal_cut_factors,
        thm:mpu_simultaneous_balanced_cut_metrics}
    Every finite unitary on a positive-length chain with $d>0$
    admits endpoint-compatible minimal cut bases and a simultaneous
    family of balanced positive metrics. At cut $j$, the metric
    $P_j$ belongs to the actual prefix Gram affine hull,
    $P_j$ and $Q_j=P_j^{-1}/r_j$ are positive definite,
    \[
      \operatorname{Tr}(XQ_j)=1
        \quad\text{for every }X\text{ in that affine hull},
      \qquad
      \operatorname{Tr}(Q_j^{-1}P_j^{-1})=r_j^2.
    \]
    No factorization, spanning, isometry, or metric witnesses
    need to be supplied.
\end{theorem}
\begin{proof}\leanok
    The constructed factors give the original unitary with its
    configurations bijectively relabelled, so their operator is
    isometric. Both spans follow from the minimal physical cut
    factorization and the bijection of paired configurations.
    Choose endpoint-compatible bases and apply simultaneous
    determinant balancing to these derived factors.
\end{proof}

\begin{theorem}[Dimensions of encoded bond registers]
    \label{thm:mpu_encoded_bond_register_dimensions}
    \lean{MPUCircuit.bond_register_dimension_bounds,
        MPUCircuit.nonempty_bond_register_embedding,
        MPUCircuit.exists_bond_register_embedding_zero,
        MPUCircuit.two_bond_register_dimension_le,
        MPUCircuit.leaf_packet_dimension_le}
    \leanok
    For $d\geq2$ and $r\geq1$, a bond basis embeds in
    $q=\lceil\log_d r\rceil$ qudits with its zero label represented
    by the all-zero configuration, and
    $r\leq d^q\leq dr$. Two encoded bonds of dimensions $r,s$
    and one physical site have padded dimension at most $d^3rs$.
    The same bound holds after increasing the packet to at least
    two sites for pair-gate synthesis.
\end{theorem}
\begin{proof}\leanok
    The ceiling logarithm gives both dimension inequalities.
    Embed the finite basis by cardinality and interchange its
    zero image with the all-zero configuration. Multiply the
    two flag bounds; the two-site exceptional packet also fits
    the displayed estimate.
\end{proof}

\begin{definition}[Initialized basis inclusion]
    \label{def:mpu_initialized_basis_inclusion}
    \lean{MPSPreparation.initializedBasisMatrix,
        MPSPreparation.zeroWorkspaceEmbedding}
    \leanok
    For an injection of logical basis vectors into a larger basis,
    let $J$ be its inclusion matrix. Appending an all-zero temporary
    register gives the zero-workspace inclusion as a special case.
\end{definition}

\begin{definition}[Implementation with a reusable initialized workspace]
    \label{def:mpu_clean_unitary_implementation}
    \lean{MPSPreparation.IsCleanImplementation}
    \leanok
    The ambient operator $U$ implements the logical operator $Z$
    with reusable initialized workspace when $UJ=JZ$. This identity
    applies to the entire logical space.
\end{definition}

\begin{theorem}[Composition and inverse calls with one workspace]
    \label{thm:mpu_clean_unitary_composition}
    \lean{MPSPreparation.IsCleanImplementation.mul,
        MPSPreparation.IsCleanImplementation.one,
        MPSPreparation.IsCleanImplementation.pow,
        MPSPreparation.IsCleanImplementation.conjTranspose,
        MPSPreparation.IsCleanImplementation.star,
        MPSPreparation.IsCleanImplementation.inv,
        MPSPreparation.IsCleanImplementation.mulVec,
        MPSPreparation.IsCleanImplementation.logical_mem_unitary,
        MPSPreparation.IsCleanImplementation.conjTranspose_of_isometry,
        MPSPreparation.IsCleanImplementation.kronecker_one}
    \leanok
    \uses{def:mpu_clean_unitary_implementation}
    Initialized-workspace implementations compose and repeat using
    the same workspace. If $J^\dagger J=I$ and $U$ is unitary,
    $UJ=JZ$ forces $Z$ to be unitary and gives
    $U^\dagger J=JZ^\dagger$. Thus inverse calls also return the
    workspace to its initialized state on every logical input.
    The identities persist after tensoring with an arbitrary
    finite reference identity.
\end{theorem}
\begin{proof}\leanok
    Multiply the two intertwining identities. Taking the Gram
    of $UJ=JZ$ gives $Z^\dagger Z=I$; ambient and logical unitarity
    then give the inverse identity. Tensor multiplication preserves
    the intertwining identity with a reference system.
\end{proof}

\begin{theorem}[Encoding and implementing isometric columns]
    \label{thm:mpu_initialized_isometry_synthesis}
    \lean{MPSPreparation.exists_isPairProduct_mul_initializedBasisMatrix_eq,
        MPSPreparation.initializedBasisMatrix_isIsometry,
        MPSPreparation.mul_initializedBasisMatrix,
        MPSPreparation.exists_isPairProduct_isCleanImplementation,
        MPSPreparation.isCleanImplementation_zeroRegisterReflection}
    \leanok
    \uses{def:mpu_initialized_basis_inclusion,
        thm:mpu_quantitative_unitary_synthesis,
        thm:mpu_clean_zero_register_reflection}
    A basis inclusion is an isometry, and multiplication by it selects
    the corresponding columns. For an isometry $V$, injective input
    and output encodings in an $n\geq2$ qudit register give an exact
    circuit $U$ with $UJ_{\rm in}=J_{\rm out}V$ and at most
    $38(d^n)^6$ gates. When $V$ is a logical unitary and the two
    inclusions coincide, this is an initialized-workspace implementation.
    The constructed zero-register reflection also satisfies this
    complete logical-space identity.
\end{theorem}
\begin{proof}\leanok
    Encoded columns remain orthonormal. Extend them to a unitary
    with their prescribed input columns and synthesize that finite
    unitary. The zero-register assertion is its already proved
    exact clean-input matrix identity.
\end{proof}

\begin{definition}[Intervals in minimal cut coordinates]
    \label{def:mpu_minimal_cut_intervals}
    \lean{MPSPreparation.CutIntervalConfig,
        MPSPreparation.joinCutPrefix,
        MPSPreparation.joinCutSuffix,
        MPSPreparation.cutPrefixIntervalConfigEquiv,
        MPSPreparation.cutIntervalSliceLinearMap,
        MPSPreparation.cutIntervalSlice,
        MPSPreparation.cutIntervalMatrix,
        MPSPreparation.joinCutInterval,
        MPUCircuit.minimalOperatorInterval}
    \leanok
    For cuts $j\leq k$, an interval configuration assigns physical
    letters to the sites $j,\ldots,k-1$. Fixing these letters restricts
    the column space at cut $k$ to the column space at cut $j$.
    In the chosen cut bases, denote this restriction by $A_{j,k}$.
    Separating every physical letter into output and input letters
    gives the corresponding operator interval.
\end{definition}

\begin{theorem}[One metric family for every interval]
    \label{thm:mpu_metrics_for_all_minimal_intervals}
    \lean{MPSPreparation.cutCoefficientMatrix_interval,
        MPSPreparation.cutIntervalSliceLinearMap_maps_columnSpace,
        MPSPreparation.cutPrefixMatrix_interval,
        MPSPreparation.cutIntervalSlice_column,
        MPSPreparation.cutSuffixMatrix_interval,
        MPSPreparation.joinCutPrefix_assoc,
        MPSPreparation.cutIntervalSlice_comp,
        MPSPreparation.cutIntervalMatrix_comp,
        MPSPreparation.cutIntervalMatrix_adjacent,
        MPUCircuit.prefixInputGram_physical_reindex,
        MPUCircuit.densityPrefixGrams_physical_reindex,
        MPUCircuit.prefixGramAffineHull_physical_reindex,
        MPUCircuit.pairPhysicalConfig_joinCutPrefix,
        MPUCircuit.minimalOperatorPrefix_concat,
        MPUCircuit.prefixGramAffineHull_minimalOperatorPrefix_concat,
        MPUCircuit.minimalOperatorInterval_isometry_of_cut_metrics,
        MPUCircuit.pairPhysicalConfig_joinCutInterval,
        MPUCircuit.minimalOperatorInterval_comp,
        MPUCircuit.exists_simultaneous_balanced_interval_isometries_of_unitary}
    \leanok
    Let $U$ be a finite unitary, with $d>0$ and $N>0$.
    Its minimal cut bases admit positive definite metrics which
    normalize every interval $[j,k]$, simultaneously for all
    $0\leq j\leq k\leq N$. The interval matrices satisfy
    $A_{j,l}=A_{j,k}A_{k,l}$, with physical configurations concatenated.
    The statement includes inputs entangled among the interval sites;
    it requires no externally chosen interval isometries.
\end{theorem}
\begin{proof}\leanok
    Prefix and suffix restrictions give the interval factorization
    in the minimal cut bases. Physical reindexing preserves the
    affine Gram spaces. Their transfer identity, applied to the
    single family of balanced cut metrics, makes every weighted
    interval an isometry. Composition follows by applying successive
    restrictions to a basis vector.
\end{proof}

\begin{definition}[Physical dilation of a joining contraction]
    \label{def:mpu_balanced_interval_dilation}
    \lean{Matrix.basisResetOutput,
        MPUCircuit.intervalChildrenRegrouping,
        MPUCircuit.vectorizedIntervalColumns,
        MPUCircuit.regroupedIntervalChildren,
        MPUCircuit.liftedNormalizedBondReset,
        MPUCircuit.balancedIntervalDilation,
        MPUCircuit.balancedIntervalChildren,
        MPUCircuit.dilatedBalancedIntervalChildren}
    \leanok
    Write the two child isometries as a tensor product and regroup
    their outputs into physical coordinates, outer bond coordinates,
    and the two joining bond coordinates. Apply the normalized
    bond-reset dilation only to the joining coordinates and its
    binary flag. The remaining coordinates carry the identity.
\end{definition}

\begin{theorem}[Exact joining with erased bond and binary flags]
    \label{thm:mpu_balanced_interval_dilation}
    \lean{Matrix.basisResetOutput_gram,
        Matrix.reindex_rows_gram,
        Matrix.partialIsometryDilation_one_kronecker,
        MPUCircuit.regroupedIntervalChildren_eq_reindex,
        MPUCircuit.regroupedIntervalChildren_gram,
        MPUCircuit.liftedNormalizedBondReset_mul_children,
        MPUCircuit.balancedIntervalDilation_eq_local,
        MPUCircuit.liftedNormalizedBondReset_mul_conjTranspose_mul,
        MPUCircuit.balancedIntervalDilation_mem_unitaryGroup,
        MPUCircuit.dilatedBalancedIntervalChildren_gram,
        MPUCircuit.dilatedBalancedIntervalChildren_success_gram,
        MPUCircuit.dilatedBalancedIntervalChildren_normalized_success,
        MPUCircuit.balancedInterval_merging_exact}
    \leanok
    Suppose the two weighted child maps and their parent map are
    isometries, formed with the same positive definite joining metric
    on a bond of dimension $r>0$. The actual bond dilation is unitary.
    Its successful branch is the parent isometry with both joining
    bond labels reset to a specified basis vector, scaled by $1/r$.
    Its success Gram matrix is $r^{-2}I$. After the prescribed
    attenuation and $\ell_r$ exact amplification rounds, the result
    is the parent isometry, with the joining bonds and both binary
    flags reset. The global phase is retained.
\end{theorem}
\begin{proof}\leanok
    The normalized contraction row has norm one. Its partial-isometry
    dilation is unitary, and its first branch contracts precisely
    the two joining indices. Tensor-product Gram identities give
    the constant successful probability. The exact uniform-merging
    identity then erases both binary flags.
\end{proof}

\begin{definition}[A finite bond-dilation packet]
    \label{def:mpu_bond_dilation_packet}
    \lean{MPUCircuit.bondDilationPacketSites}
    \leanok
    For a joining bond of dimension $r$, use a qudit packet of
    $\max\{2,\lceil\log_d(2r^2)\rceil+1\}$ sites to encode the two
    copies of the joining-pair space in the dilation.
\end{definition}

\begin{theorem}[An actual circuit for the bond dilation]
    \label{thm:mpu_bond_dilation_packet}
    \lean{MPUCircuit.bondDilationPacket_dimension_le,
        MPUCircuit.exists_bondDilationPacket_embedding_zero,
        MPUCircuit.exists_normalizedBondDilation_circuit}
    \leanok
    For $d\geq2$, $r>0$, and a positive definite joining metric,
    the packet admits an encoding whose prescribed reset vector is
    physical zero and an exact neighboring-pair circuit implementing
    the bond dilation on its entire encoded space. Its dimension is
    at most $2d^3r^2$, and its gate count is at most
    $38(2d^3r^2)^6$.
\end{theorem}
\begin{proof}\leanok
    Choose a cardinality embedding and exchange the prescribed reset
    vector with zero. Extend the encoded logical unitary to the
    whole packet, then apply quantitative unitary synthesis.
\end{proof}

\begin{definition}[Reflections on selected initialized flags]
    \label{def:mpu_selected_flag_reflection}
    \lean{MPSPreparation.selectedZeroRegisterSites,
        MPSPreparation.selectedZeroRegisterLogical,
        MPSPreparation.selectedZeroRegisterReflection,
        MPSPreparation.selectedZeroRegisterReflectionGateCount,
        MPSPreparation.selectedZeroRegisterPoolSites,
        MPSPreparation.selectedZeroRegisterReflectionPool,
        MPSPreparation.selectedZeroRegisterReflectionPoolGateCount}
    \leanok
    An injective choice of $a$ sites in an $n$-site logical register
    defines the diagonal reflection with phase $-1$ precisely when
    those sites are all zero. A common pool of $A\geq a$ initialized
    scratch sites supplies its reversible conjunction circuit.
    Denote the zero-register gate bound by $K_d(a)$.
\end{definition}

\begin{theorem}[Selected flags with one common scratch pool]
    \label{thm:mpu_selected_flag_reflection}
    \lean{MPSPreparation.selectedZeroRegisterSites_castAdd,
        MPSPreparation.selectedZeroRegisterSites_natAdd,
        MPSPreparation.selectedZeroRegisterLogical_mem_unitary,
        MPSPreparation.selectedZeroRegisterLogical_eq_one_sub,
        MPSPreparation.isPairProduct_selectedZeroRegisterReflection,
        MPSPreparation.isCleanImplementation_selectedZeroRegisterReflection,
        MPSPreparation.selectedZeroRegisterLogical_empty,
        MPSPreparation.exists_isPairProduct_isCleanImplementation_selectedZeroRegister,
        MPSPreparation.selectedZeroRegisterPoolSites_castAdd,
        MPSPreparation.selectedZeroRegisterPoolSites_natAdd,
        MPSPreparation.isPairProduct_selectedZeroRegisterReflectionPool,
        MPSPreparation.isCleanImplementation_selectedZeroRegisterReflectionPool,
        MPSPreparation.exists_isPairProduct_isCleanImplementation_selectedZeroRegisterPool}
    \leanok
    For $d\geq2$ and $n\geq2$, the selected-flags reflection has an
    actual circuit on $n+A$ sites. It returns all $A$ scratch sites
    to zero on every logical input and costs at most
    $2(n+A)K_d(a)$ gates when $a>0$. For $a=0$, one neighboring-pair
    gate implements the required phase $-1$. Unselected logical
    sites are unrestricted.
\end{theorem}
\begin{proof}\leanok
    Place the reversible zero conjunction on the selected flags and
    the first $a$ scratch sites, apply its phase, and reverse it.
    Its exact clean-input identity gives the logical reflection,
    with every unused scratch coordinate unchanged. The empty case
    is the scalar unitary $-I$.
\end{proof}

\begin{definition}[Attenuation on one physical qudit]
    \label{def:mpu_physical_success_attenuation}
    \lean{MPUCircuit.attenuationQuditRotation}
    \leanok
    Place the real binary attenuation rotation on levels zero and
    one of a qudit, and let it be the identity on the other levels.
\end{definition}

\begin{theorem}[A circuit for the prescribed attenuation]
    \label{thm:mpu_physical_success_attenuation}
    \lean{MPUCircuit.attenuationQuditRotation_mem_unitary,
        MPUCircuit.attenuationQuditRotation_mulVec_zero,
        MPUCircuit.isPairProduct_siteOp_attenuationQuditRotation,
        MPUCircuit.isPairProduct_siteOp_mergingAttenuation}
    \leanok
    The rotation is unitary whenever $t^2\leq1$, and sends zero to
    $t\lvert0\rangle+\sqrt{1-t^2}\lvert1\rangle$ exactly.
    At any site of an $n\geq2$ qudit chain it costs at most $2n$
    neighboring-pair gates. The prescribed merger attenuation
    $t=r\sin(\pi/(4\ell_r+2))$ satisfies the required inequality.
\end{theorem}
\begin{proof}\leanok
    The binary rotation is extended by the identity. Its first
    column gives the asserted amplitudes. Route its single-site
    support along the chain. The amplification budget proves
    feasibility of the prescribed angle.
\end{proof}

\begin{definition}[Registers indexed by internal physical cuts]
    \label{def:mpu_common_register_layout}
    \lean{MPUCircuit.bondRegisterWidth,
        MPUCircuit.internalCutEmbedding,
        MPUCircuit.BondRegisterSite,
        MPUCircuit.AmplificationFlagSite,
        MPUCircuit.AuxiliarySite,
        MPUCircuit.LogicalSite,
        MPUCircuit.GlobalSite,
        MPUCircuit.auxiliarySiteCount,
        MPUCircuit.logicalSiteCount,
        MPUCircuit.globalSiteCount,
        MPUCircuit.auxiliarySiteEquivFin,
        MPUCircuit.logicalSiteEquivFin,
        MPUCircuit.globalSiteEquivFin,
        MPUCircuit.physicalSiteEmbedding,
        MPUCircuit.bondSiteEmbedding,
        MPUCircuit.amplificationFlagEmbedding,
        MPUCircuit.auxiliarySiteEmbedding,
        MPUCircuit.logicalSiteEmbedding,
        MPUCircuit.scratchSiteEmbedding,
        MPUCircuit.physicalLogicalSites,
        MPUCircuit.auxiliaryLogicalSites,
        MPUCircuit.logicalGlobalSites,
        MPUCircuit.scratchGlobalSites,
        MPUCircuit.bondRegisterEncoding}
    \leanok
    Put $q=\lceil\log_d D\rceil$. At each internal physical cut
    allocate two $q$-site bond registers and two amplification flags.
    The logical register consists of these sites and the physical
    chain. Append a scratch pool with as many sites as the logical
    auxiliary register. Consecutive indexing preserves the order
    physical sites, logical auxiliary sites, then scratch sites.
\end{definition}

\begin{theorem}[Register counts and compatible bond encodings]
    \label{thm:mpu_common_register_layout}
    \lean{MPUCircuit.internalCutEmbedding_bounds,
        MPUCircuit.card_bondRegisterSite,
        MPUCircuit.card_amplificationFlagSite,
        MPUCircuit.card_auxiliarySite,
        MPUCircuit.card_logicalSite,
        MPUCircuit.card_globalSite,
        MPUCircuit.globalSiteCount_eq,
        MPUCircuit.globalSiteCount_le,
        MPUCircuit.physicalLogicalSites_apply,
        MPUCircuit.auxiliaryLogicalSites_apply,
        MPUCircuit.logicalGlobalSites_apply,
        MPUCircuit.scratchGlobalSites_apply,
        MPUCircuit.bondDim_le_registerCapacity,
        MPUCircuit.registerCapacity_le_mul_bondDim,
        MPUCircuit.bondRegisterEncoding_zero}
    \leanok
    The logical auxiliary register has $A=2(N-1)(q+1)$ sites,
    and the complete register has
    \[
      N+2A=N+4(N-1)(q+1)\leq5N(q+1).
    \]
    For $D>0$ one has $D\leq d^q\leq dD$. Every positive bond
    dimension $r\leq D$ admits an injective encoding in its assigned
    $q$ qudits, with bond label zero represented by physical zero.
    For $N=1$ the layout has one site; pair-gate synthesis may append
    one initialized padding site.
\end{theorem}
\begin{proof}\leanok
    Count the finite products and disjoint unions defining the site
    sets. Ceiling-logarithmic capacity gives the dimension inequalities.
    A cardinality embedding followed by a transposition of two
    configurations gives the required zero-preserving encoding.
\end{proof}

\begin{definition}[Joining in the given child bond encodings]
    \label{def:mpu_compatible_bond_dilation}
    \lean{MPUCircuit.joiningBondPair,
        MPUCircuit.joiningDilationFlag,
        MPUCircuit.compatibleBondDilationCfg,
        MPUCircuit.compatibleBondDilationEmbedding}
    \leanok
    Encode the two joining bond labels using the given child bond
    encodings, followed by the dilation flag and one attenuation
    qudit. The logical bond dilation acts as the identity on every
    level of the attenuation qudit. The order of the two bond labels
    agrees with the vectorized joining contraction.
\end{definition}

\begin{theorem}[Compatible physical joining circuit]
    \label{thm:mpu_compatible_bond_dilation}
    \lean{MPUCircuit.compatibleBondDilationCfg_left,
        MPUCircuit.compatibleBondDilationCfg_right,
        MPUCircuit.compatibleBondDilationCfg_dilation,
        MPUCircuit.compatibleBondDilationCfg_attenuation,
        MPUCircuit.compatibleBondDilationCfg_injective,
        MPUCircuit.compatibleBondDilationEmbedding_zero,
        MPUCircuit.exists_compatibleNormalizedBondDilation_circuit,
        MPUCircuit.compatibleBondDilationPacket_dimension_le,
        MPUCircuit.exists_compatibleNormalizedBondDilation_circuit_bounded}
    \leanok
    Let a positive bond of dimension $r$ have a zero-preserving
    encoding in $q$ qudits, and let its joining metric be positive
    definite. There is an actual circuit on $2q+2$ qudits implementing
    the bond dilation in precisely this product encoding, preserving
    the entire attenuation qudit. Its reset vector is all zero.
    For $q=\lceil\log_dD\rceil$ and $D>0$, the packet dimension is
    at most $d^4D^2$ and its gate count is at most $38(d^4D^2)^6$.
\end{theorem}
\begin{proof}\leanok
    The product encoding is injective and maps the specified reset
    labels and flags to zero. Extend the encoded full logical unitary,
    including the identity on the attenuation qudit, to the whole
    packet. Quantitative synthesis and the register-capacity estimate
    give the stated gate bound.
\end{proof}

\begin{definition}[Separate leaf bond encodings]
    \label{def:mpu_clean_leaf_encoding}
    \lean{MPUCircuit.leafOutputEmbedding,
        MPUCircuit.leafInputEmbedding}
    \leanok
    A one-site interval output consists of a physical letter and
    two bond labels. Encode the two labels separately in their
    assigned registers, allowing width zero at an endpoint cut. Its initialized input contains only
    the physical letter, with both bond registers in zero.
\end{definition}

\begin{theorem}[A clean circuit for every leaf isometry]
    \label{thm:mpu_clean_leaf_circuit}
    \lean{MPUCircuit.exists_leaf_logical_unitary,
        MPUCircuit.exists_exact_clean_leaf_circuit,
        MPUCircuit.leaf_workspace_dimension_le,
        MPUCircuit.leaf_workspace_gateCount_le,
        MPUCircuit.leaf_workspace_dimension_le_of_width_le}
    \leanok
    Every isometric one-site interval map admits a full logical
    unitary extension in these given product encodings, and an actual
    neighboring-pair circuit using one initialized scratch qudit.
    The scratch qudit returns to zero on every logical state.
    When both bond-register widths are at most
    $q=\lceil\log_dD\rceil$ and $D>0$, the dimension including
    scratch is at most $d^4D^2$, and the gate cost is at most
    $38(d^4D^2)^6$. The assertion also covers $q=0$.
\end{theorem}
\begin{proof}\leanok
    Encode the orthonormal output columns and extend them to a full
    logical unitary with the prescribed initialized input columns.
    Extend this full unitary along the one-scratch-qudit inclusion
    and synthesize the resulting unitary. The complete clean identity
    ensures cleanup of adjoint calls as well.
\end{proof}

\begin{theorem}[Exact amplification with reused workspace]
    \label{thm:mpu_exact_clean_amplification_circuit}
    \lean{MPUCircuit.exists_exact_clean_amplification_circuit}
    \leanok
    Let $J$ be an initialized isometric inclusion. Suppose an actual
    child circuit of cost $K$ implements a logical unitary $Z$
    cleanly along $J$, and actual reflection circuits of costs $B,L$
    implement $I-2EE^\dagger$ and $I-2P$ cleanly along the same $J$.
    Here $E$ is an isometry and $P$ an orthogonal projection. If
    \[
      (ZE)^\dagger P(ZE)
       =\sin^2\!\left(\frac{\pi}{4\ell+2}\right) I,
      \qquad \ell>0,
    \]
    the amplified circuit has cost
    $\ell(1+2K+B+L)+K$ and acts exactly as normalized postselection
    on the columns $JE$. It also implements a full logical unitary
    cleanly along $J$, so its adjoint may use the same workspace.
    The supplied circuits are the induction data for this assembly
    statement.
\end{theorem}
\begin{proof}\leanok
    Ambient unitarity and the clean identity imply full logical
    unitarity of $Z$, hence cleanup of the inverse child call.
    Each amplification round is the child circuit, the initial
    reflection, its inverse, the successful reflection, and the
    scalar phase $-1$. Products and powers preserve the complete
    clean identity. The exact subspace-amplification formula gives
    the normalized successful columns with their phase retained.
\end{proof}

\begin{definition}[Placement of a circuit and its initialized workspace]
    \label{def:mpu_clean_circuit_placement}
    \lean{MPSPreparation.cleanImplementationSites}
    \leanok
    Choose injective placements of the local logical sites in the
    full logical register and of the local scratch sites in the
    common scratch pool. Their disjoint union places the entire
    local circuit, preserving the logical and scratch blocks.
\end{definition}

\begin{theorem}[Clean placement into a larger shared pool]
    \label{thm:mpu_clean_circuit_placement}
    \lean{MPSPreparation.cleanImplementationSites_castAdd,
        MPSPreparation.cleanImplementationSites_natAdd,
        MPSPreparation.initializedBasisMatrix_zeroWorkspace_mul_apply,
        MPSPreparation.IsCleanImplementation.embedOp,
        MPSPreparation.isPairProduct_isCleanImplementation_embedOp}
    \leanok
    Suppose a local circuit satisfies its complete clean identity
    on every logical input. Its placed circuit satisfies the same
    identity on the full logical register and restores the entire
    larger scratch pool to zero. If the local circuit has $K$ gates,
    its routed circuit on $n+A$ sites has at most $2(n+A)K$ gates.
    No additional cleanup hypothesis is required for the placed circuit.
\end{theorem}
\begin{proof}\leanok
    The initialized basis inclusion selects the zero-scratch columns.
    In these columns, the placed matrix identity is precisely the
    local clean identity, with identity factors on every unused site.
    Injective site placement gives the stated routing bound.
\end{proof}

\begin{theorem}[An explicit logarithmic branching exponent]
    \label{thm:mpu_explicit_amplified_recursion_bound}
    \lean{MPUCircuit.amplification_gateCount_le,
        MPUCircuit.prescribed_amplification_gateCount_le,
        MPUCircuit.gateCount_add_overhead_le_pow,
        MPUCircuit.gateCount_le_polynomial_of_balanced_height,
        MPUCircuit.amplification_branching_exponent_le}
    \leanok
    If the two child circuits each cost at most $T$, the joining
    primitive costs at most $M$, and each reflection costs at most $F$,
    exact amplification at a bond of dimension at most $D$ costs at most
    \[
      (4D+2)T+(2D+1)M+D(1+2F).
    \]
    A numerical recurrence $T_{j+1}\leq RT_j+C$ with $R\geq2$
    and $T_0\leq a$ satisfies, at balanced height
    $L=\lceil\log_2N\rceil$,
    \[
      T_L\leq(a+C)(2N)^{\lceil\log_2R\rceil}.
    \]
    For $D>0$, the branching exponent satisfies
    $\lceil\log_2(4D+2)\rceil\leq\lceil\log_2D\rceil+3$.
    These assertions are numerical consequences of the stated recurrence.
\end{theorem}
\begin{proof}\leanok
    Substitute the joining and reflection bounds into the exact
    amplification gate count. Induct on $j$ using
    $T_j+C\leq(a+C)R^j$. Ceiling-logarithmic bounds give
    $R\leq2^{\lceil\log_2R\rceil}$ and $2^L\leq2N$.
    Finally $4D+2\leq8D$ for $D\geq1$.
\end{proof}

\begin{definition}[Coordinates of the complete interval]
    \label{def:mpu_full_minimal_interval_coordinates}
    \lean{MPUCircuit.fullCutIntervalConfigEquiv}
    \leanok
    The configurations of the interval between the two endpoint cuts
    are canonically identified with physical configurations of the
    complete chain. Both endpoint bond spaces have their normalized
    one-dimensional bases.
\end{definition}

\begin{theorem}[The normalized root is the original unitary]
    \label{thm:mpu_normalized_minimal_interval_root}
    \lean{MPUCircuit.prefixInputGram_minimalOperatorPrefix_zero,
        MPUCircuit.prefixInputGram_minimalOperatorPrefix_last,
        MPUCircuit.prefixGramAffineHull_minimalOperatorPrefix_zero,
        MPUCircuit.prefixGramAffineHull_minimalOperatorPrefix_last,
        MPUCircuit.eq_one_of_mem_minimalOperatorPrefixHull_zero,
        MPUCircuit.eq_one_of_mem_minimalOperatorPrefixHull_last,
        MPUCircuit.minimalOperatorInterval_full_apply,
        MPUCircuit.weightedMinimalOperatorInterval_full_apply,
        MPUCircuit.weightedMinimalOperatorInterval_full_submatrix,
        MPUCircuit.exists_normalized_minimalInterval_family_of_unitary}
    \leanok
    For every finite unitary with $d>0$ and $N>0$, the actual endpoint
    affine Gram spaces of the normalized minimal representation are
    singleton identity spaces. The coherent metric family therefore
    has $P_0=P_N=I$. Its weighted complete interval is exactly the
    original unitary, entry by entry, including its complex phase.
    This family and all its interval isometries are derived from
    unitarity; no metric or interval witness is supplied.
\end{theorem}
\begin{proof}\leanok
    The empty prefix has Gram equal to its input trace. At the final
    cut, global unitarity gives the same identity. Every normalized
    endpoint Gram is thus the identity, as is its entire affine hull.
    Minimal-coordinate reconstruction identifies the complete interval
    with the original matrix. Identity endpoint weights preserve it.
\end{proof}

\begin{theorem}[The source transpose convention for a balanced metric]
    \label{thm:mpu_balanced_source_conditioning}
    \lean{MPUCircuit.trace_transpose_inv_mul_transpose_inv,
        MPUCircuit.trace_transpose_dualGramMetric_inv_mul_transpose_inv,
        MPUCircuit.sqrt_re_trace_transpose_dualGramMetric_inv_mul_transpose_inv}
    \leanok
    For positive definite $P$ on a bond of dimension $r>0$ and
    $Q=P^{-1}/r$, the source's transpose expression satisfies
    \[
      \operatorname{Tr}\bigl[(Q^{-1})^T(P^{-1})^T\bigr]=r^2,
      \qquad
      \sqrt{\operatorname{Re}\operatorname{Tr}
        \bigl[(Q^{-1})^T(P^{-1})^T\bigr]}=r.
    \]
    These metrics give an explicit candidate of joining norm $r$.
    The assertion does not bound the appendix's infimum restricted
    to physical cap Grams in \cite{Styliaris2025MPUCircuits}.
\end{theorem}
\begin{proof}\leanok
    Reverse the matrix product under transpose and use invariance
    of trace under transpose. The balanced inverse-trace identity
    gives $r^2$. Its real part is nonnegative and its positive
    square root is $r$. The interval and flag conventions are
    expanded explicitly in \cite{RankTwoMPUCircuits2026}, Section~5.
\end{proof}

\begin{definition}[The physical registers of an interval]
    \label{def:mpu_interval_register_support}
    \lean{MPUCircuit.intervalLogicalSupport,
        MPUCircuit.joiningFlagSupport,
        MPUCircuit.cutBondRegisterWidth,
        MPUCircuit.cutBondSiteEmbedding,
        MPUCircuit.LeafLogicalSite,
        MPUCircuit.leafLogicalSiteEmbedding,
        MPUCircuit.leafPacketSites,
        MPUCircuit.intervalConsecutiveSupport,
        MPUCircuit.intervalInteriorAuxiliarySupport,
        MPUCircuit.joiningAuxiliarySupport,
        MPUCircuit.IsIntervalInteriorInitialized,
        MPUCircuit.leafConsecutiveSites}
    \leanok
    The support of an interval consists of its physical sites, all
    strictly interior bond and amplification sites, the inward-facing
    bond register at its left boundary, and the inward-facing bond
    register at its right boundary. Endpoint bond registers have
    width zero. The interior is initialized when every strictly
    interior auxiliary site is in its designated initial configuration.
    A leaf packet has one physical site and its two available
    boundary bond registers.
\end{definition}

\begin{theorem}[Disjoint children and the exact parent register]
    \label{thm:mpu_interval_register_support}
    \lean{MPUCircuit.intervalLogicalSupport_disjoint,
        MPUCircuit.joiningFlagSupport_disjoint_left,
        MPUCircuit.joiningFlagSupport_disjoint_right,
        MPUCircuit.intervalLogicalSupport_split,
        MPUCircuit.cutBondRegisterWidth_le,
        MPUCircuit.cutBondRegisterWidth_internal,
        MPUCircuit.cutBondSiteEmbedding_disjoint,
        MPUCircuit.cutBondSiteEmbedding_mem_left,
        MPUCircuit.cutBondSiteEmbedding_mem_right,
        MPUCircuit.leafLogicalSiteEmbedding_mem,
        MPUCircuit.leafPacketSites_mem,
        MPUCircuit.intervalConsecutiveSupport_disjoint,
        MPUCircuit.interval_supportedOperators_commute,
        MPUCircuit.cutBondSiteEmbedding_range_iff,
        MPUCircuit.leafLogicalSiteEmbedding_range,
        MPUCircuit.leafPacketSites_range,
        MPUCircuit.cutBondRegisterWidth_zero,
        MPUCircuit.cutBondRegisterWidth_last,
        MPUCircuit.intervalInteriorAuxiliarySupport_split,
        MPUCircuit.intervalInteriorAuxiliarySupport_leaf,
        MPUCircuit.isIntervalInteriorInitialized_split,
        MPUCircuit.intervalLogicalSupport_full,
        MPUCircuit.isIntervalInteriorInitialized_leaf,
        MPUCircuit.leafConsecutiveSites_range,
        MPUCircuit.leaf_embedOp_mem_supportedOperators}
    \leanok
    For adjacent nonempty child intervals $[j,m)$ and $[m,k)$,
    their logical supports are disjoint. Their union together with
    the two amplification flags at $m$ is exactly the parent support.
    Operators supported on the two child intervals commute.
    The parent interior is initialized precisely when both child
    interiors and all joining auxiliaries are initialized.
    A leaf has no interior auxiliary sites, and its packet placement
    is injective with range exactly its interval support. Thus no
    neighboring bond register is borrowed at an endpoint.
\end{theorem}
\begin{proof}\leanok
    Separate physical sites, bond sites, and flags. The two boundary
    sides at a joining cut are different registers; its flags belong
    to neither child. These observations give both disjointness and
    the stated union. The explicit leaf placement has the same range.
    Disjoint support implies commutation, and the partition of
    interior auxiliary sites gives the initialization equivalence.
\end{proof}

\begin{definition}[The inclusion with selected flags initialized]
    \label{def:mpu_initialized_flag_inclusion}
    \lean{MPSPreparation.zeroFlagEmbedding}
    \leanok
    Given an injective selection of qudit flags, include arbitrary
    configurations on the remaining sites by assigning zero to each
    selected flag. Denote its basis-inclusion matrix by $E$.
\end{definition}

\begin{theorem}[The constructed reflection is the initialized-range reflection]
    \label{thm:mpu_initialized_flag_projection}
    \lean{MPSPreparation.initializedBasisMatrix_mul_conjTranspose,
        MPSPreparation.zeroFlagEmbedding_apply_selected,
        MPSPreparation.mem_range_zeroFlagEmbedding_iff,
        MPSPreparation.zeroFlagEmbedding_range_projection,
        MPSPreparation.isStarProjection_zeroFlagEmbedding_range,
        MPSPreparation.selectedZeroRegisterLogical_eq_subspaceReflection,
        MPSPreparation.exists_isPairProduct_isCleanImplementation_zeroFlagReflection}
    \leanok
    For any injective basis encoding, its range projection is the
    diagonal indicator of its encoded configurations. Consequently,
    the selected-flag inclusion has
    \[
      EE^\dagger=\operatorname{diag}
      \bigl(\mathbf 1_{\text{all selected flags zero}}\bigr).
    \]
    This is an orthogonal projection, and the reversible selected-zero
    circuit reflects it exactly. For $d\geq2$, an $n\geq2$ logical
    register and a scratch pool of size $A$ at least the number of
    selected flags, an actual circuit implements $I-2EE^\dagger$
    with the selected-reflection gate bound and returns the whole
    pool to zero on every logical input. No reflection identity or
    circuit decomposition is assumed.
\end{theorem}
\begin{proof}\leanok
    Multiplying an injective basis inclusion by its adjoint sums
    orthogonal basis projections over its image. The image of the
    selected-flag inclusion is exactly the stated zero-flag set.
    Substitute this projection identity into the already constructed
    shared-pool selected-reflection circuit.
\end{proof}

\begin{definition}[Bond encodings at the endpoint cuts]
    \label{def:mpu_endpoint_bond_encodings}
    \lean{MPUCircuit.cutBondRegisterEncoding,
        MPUCircuit.minimalCutBondRegisterEncoding}
    \leanok
    Encode the label space of every internal cut in its common bond
    register, and encode a one-dimensional endpoint label space in
    the unique configuration of a zero-site register. In each case,
    the distinguished label zero is represented by physical zero.
\end{definition}

\begin{theorem}[Actual minimal bonds fit the allocated registers]
    \label{thm:mpu_endpoint_bond_encodings}
    \lean{MPUCircuit.cutBondDim_le_registerCapacity,
        MPUCircuit.cutBondRegisterEncoding_zero,
        MPUCircuit.minimalCutRank_eq_one_of_endpoint,
        MPUCircuit.minimalCutBondRegisterEncoding_zero,
        MPUCircuit.minimalCutBondDim_le_registerCapacity,
        MPUCircuit.minimalCutBondRegisterEncoding_endpoint,
        MPUCircuit.exists_minimalCut_bondEncodings_of_unitary}
    \leanok
    A finite unitary whose minimal physical cut ranks are at most
    $D$ admits simultaneous zero-preserving encodings at all cuts
    in precisely their allocated registers. Both endpoint ranks are
    derived to be one; no endpoint rank witness is required.
    Consequently the entire endpoint label space is represented by
    the unique empty configuration.
\end{theorem}
\begin{proof}\leanok
    Nonvanishing of the operator coefficient tensor gives rank one
    at the empty and complete cuts. Its remaining cut ranks fit the
    ceiling-logarithmic capacity. Choose basis embeddings and exchange
    the image of the distinguished zero label with physical zero.
\end{proof}

\begin{definition}[The actual balanced interval tree]
    \label{def:mpu_balanced_interval_tree}
    \lean{MPUCircuit.balancedIntervalTree,
        MPUCircuit.IsIntervalPartition,
        MPUCircuit.intervalLeafSites,
        MPUCircuit.intervalTreeCuts}
    \leanok
    An interval of length one is a leaf. Split every longer interval
    at its midpoint into consecutive children of lengths
    $\lfloor n/2\rfloor$ and $\lceil n/2\rceil$. Label each internal
    node by its actual joining cut. A consecutive interval partition
    records these nonempty child intervals and their single-site leaves.
\end{definition}

\begin{theorem}[Every physical site and joining cut occurs once]
    \label{thm:mpu_balanced_interval_tree}
    \lean{MPUCircuit.midpoint_child_lengths,
        MPUCircuit.midpoint_child_balance,
        MPUCircuit.clog_two_midpoint,
        MPUCircuit.balancedIntervalTree_one,
        MPUCircuit.balancedIntervalTree_split,
        MPUCircuit.midpoint_child_clog_lt,
        MPUCircuit.balancedIntervalTree_height,
        MPUCircuit.IsIntervalPartition.leafSites,
        MPUCircuit.IsIntervalPartition.cuts,
        MPUCircuit.balancedIntervalTree_isIntervalPartition,
        MPUCircuit.balancedIntervalTree_leafSites,
        MPUCircuit.balancedIntervalTree_cuts,
        MPUCircuit.balancedIntervalTree_leafSites_nodup,
        MPUCircuit.balancedIntervalTree_cuts_nodup,
        MPUCircuit.mem_balancedIntervalTree_cuts,
        MPUCircuit.count_balancedIntervalTree_cut,
        MPUCircuit.balancedIntervalTree_numLeaves,
        MPUCircuit.balancedIntervalTree_numNodes,
        MPUCircuit.exists_balancedIntervalPartition}
    \leanok
    For every $N>0$, midpoint splitting constructs a genuine binary
    partition of the $N$ consecutive physical sites, without a
    supplied tree. Its leaves are precisely those $N$ sites, in
    physical order; its internal nodes are precisely the $N-1$
    internal cuts, also in physical order. Every such cut occurs
    exactly once. Its height is exactly $\lceil\log_2N\rceil$,
    including lengths which are not powers of two.
\end{theorem}
\begin{proof}\leanok
    Strong induction on the interval length constructs its two
    strictly shorter nonempty children. Concatenate their site and
    cut ranges. The larger child has ceiling logarithm one less
    than the parent, which proves the exact height formula.
\end{proof}

\begin{definition}[The gate count on the actual interval tree]
    \label{def:mpu_actual_tree_gate_count}
    \lean{MPUCircuit.treeGateCount,
        MPUCircuit.amplificationTreeOverhead,
        MPUCircuit.amplificationTreeGateCount}
    \leanok
    Give every leaf cost $K$. At a node whose children have costs
    $T_1,T_2$, let $T_0=T_1+T_2+M$ and assign cost
    $\ell(1+2T_0+B+L)+T_0$, where $M$ bounds the joining operation,
    $B,L$ bound the two reflections, and $\ell$ is the number of
    amplification rounds at that node. Define
    $C=(2D+1)M+D(1+B+L)$.
\end{definition}

\begin{theorem}[Numerical amplification count on the constructed tree]
    \label{thm:mpu_actual_tree_gate_count}
    \lean{MPUCircuit.treeGateCount_add_overhead_le_height,
        MPUCircuit.amplificationTreeGateCount_const_eq_treeGateCount,
        MPUCircuit.amplification_node_gateCount_le,
        MPUCircuit.amplificationTreeGateCount_le_treeGateCount,
        MPUCircuit.amplificationTreeGateCount_add_overhead_le_height,
        MPUCircuit.amplificationTreeGateCount_le_height,
        MPUCircuit.amplificationTreeGateCount_balanced_le_polynomial}
    \leanok
    If every amplification count satisfies $\ell\leq D$, then the
    recursively defined count on any binary tree of height $h$
    satisfies
    \[
      T+C\leq(K+C)(4D+2)^h.
    \]
    On the actual midpoint partition of $N>0$ consecutive sites,
    \[
      T\leq(K+C)(2N)^{\lceil\log_2(4D+2)\rceil}.
    \]
    These are bounds on the stated numerical recursion; existence
    of the corresponding recursive circuits is a separate assertion.
    Setting every round count equal to $D$ gives exactly the recursion
    with child multiplier $2D+1$ and additive term $C$.
\end{theorem}
\begin{proof}\leanok
    Each node has cost at most $(2D+1)(T_1+T_2)+C$.
    Keeping the reserve $C$ in the induction invariant proves the
    height bound. The height of the midpoint tree is
    $\lceil\log_2N\rceil$, and the ceiling-logarithmic power comparison
    gives the displayed polynomial bound.
\end{proof}

\begin{definition}[Weighted single-site intervals]
    \label{def:mpu_minimal_leaf_intervals}
    \lean{MPUCircuit.singleSiteIntervalConfigEquiv,
        MPUCircuit.minimalLeafInterval}
    \leanok
    Identify the physical configurations of a single-site interval
    with its local physical alphabet. Reindex its weighted interval
    isometry so that the output coordinates are, in order, its
    physical output, its right bond label, and its left bond label.
\end{definition}

\begin{theorem}[Coherent interval metrics and actual leaf circuits]
    \label{thm:mpu_minimal_leaf_circuits}
    \lean{MPUCircuit.minimalLeafInterval_isometry_of_interval_isometry,
        MPUCircuit.exists_normalized_minimalInterval_family_with_leaf_circuits_of_unitary}
    \leanok
    For $d\geq2$ and $N>0$, a finite unitary with consecutive operator
    Schmidt ranks at most $D$ admits one minimal open representation,
    one family of balanced positive cut metrics, and exact circuits
    for every weighted single-site interval of that same family.
    The endpoint metrics equal one and the full weighted interval
    equals the original unitary, including its complex phase.
    Each leaf circuit implements a unitary extension on its entire
    padded logical packet, using one workspace qudit which it returns
    to zero on every logical input. For left and right register widths
    $q_-,q_+$, its gate count is at most
    $38\bigl(d^{2+q_-+q_+}\bigr)^6$.
\end{theorem}
\begin{proof}\leanok
    Choose normalized bases at the two endpoint cuts and derive the
    balanced metrics and all interval isometries from unitarity.
    Reindex each one-site interval and extend its isometry to a
    unitary on the padded packet. Apply the exact unitary synthesis
    with one initialized workspace qudit. Endpoint bond labels use
    zero-site registers.
\end{proof}

\begin{definition}[Leaf placement in a common workspace]
    \label{def:mpu_minimal_leaf_placement}
    \lean{MPUCircuit.leafWorkspaceSites,
        MPUCircuit.leafCleanPlacementSites}
    \leanok
    For a chain of at least two sites, choose one qudit in the common
    workspace. Embed each leaf packet in its physical site and its
    two boundary bond registers, and place its workspace qudit in
    the chosen common site.
\end{definition}

\begin{theorem}[Actual leaf circuits in the complete register]
    \label{thm:mpu_minimal_leaf_placement}
    \lean{MPUCircuit.auxiliarySiteCount_pos_of_two_le,
        MPUCircuit.placed_leaf_circuit_isCleanImplementation,
        MPUCircuit.exists_normalized_minimalInterval_family_with_placed_leaf_circuits_of_unitary}
    \leanok
    Let $d\geq2$ and $N\geq2$. A finite unitary with consecutive
    operator Schmidt ranks at most $D$ admits coherent normalized
    interval metrics and exact leaf circuits placed in the complete
    register. Every logical leaf unitary acts only on its interval
    support. Its complete circuit returns the entire common workspace
    to zero on every logical input, with gate count at most
    \[
      76(d^4D^2)^6\,n_{\mathrm{tot}},
    \]
    where $n_{\mathrm{tot}}$ is the total number of physical, bond,
    flag, and workspace qudits.
\end{theorem}
\begin{proof}\leanok
    Place the local circuit using the two site embeddings. The clean
    implementation identity is preserved under placement. A local
    gate uses at most $2n_{\mathrm{tot}}$ gates after routing.
    The ceiling-logarithmic capacity bound gives local packet
    dimension at most $d^4D^2$. Positivity of the minimal cut ranks
    supplies $D>0$, and $N\geq2$ supplies the common workspace site.
\end{proof}

\begin{definition}[One-site padding]
    \label{def:mpu_single_site_padding}
    \lean{MPUCircuit.singleSitePaddedOperator}
    \leanok
    For a one-site operator $U$, define its two-site padding by
    applying $U$ to the first site and the identity to the second.
\end{definition}

\begin{theorem}[The one-site endpoint is one pair gate]
    \label{thm:mpu_single_site_padding}
    \lean{MPUCircuit.singleSitePaddedOperator_mem_unitary,
        MPUCircuit.singleSitePaddedOperator_isNeighbourGate,
        MPUCircuit.singleSitePaddedOperator_isPairProduct,
        MPUCircuit.singleSitePaddedOperator_isCleanImplementation,
        MPUCircuit.exists_singleSite_clean_pair_circuit}
    \leanok
    Every one-site unitary in positive physical dimension has an exact
    implementation by one neighboring two-qudit gate with one padding
    qudit. The padding qudit remains exactly zero for every physical
    input, and the complete scalar phase is retained. In fact the
    padded operator is $U\otimes I$ on the full two-site space.
\end{theorem}
\begin{proof}\leanok
    The padded operator is unitary and supported on the sole neighboring
    pair. Its initialized columns are the columns of $U$ tensored with
    the zero basis vector of the second site.
\end{proof}

\begin{definition}[Product encodings for an exact joining step]
    \label{def:mpu_encoded_merging}
    \lean{MPUCircuit.appendBasisEmbedding,
        MPUCircuit.compatibleBondFlagEmbedding,
        MPUCircuit.oneSiteBasisEmbedding,
        MPUCircuit.joiningChildBasisEmbedding,
        MPUCircuit.joiningSuccessProjection,
        MPUCircuit.mergingFlagSites,
        MPUCircuit.mergingPreparationGateCount,
        MPUCircuit.normalizedJoiningParent,
        MPUCircuit.encodedJoiningParent,
        MPUCircuit.encodedUniformMergingGateCount}
    \leanok
    Concatenate a spectator basis encoding with the two bond registers
    and their dilation and attenuation flags. Each flag is an entire
    physical qudit, with the two distinguished values zero and one.
    For a positive cut metric $P$ of dimension $r$, let $M_P$ be its
    normalized bond-reset operator. For joint child columns $V$, define
    the normalized parent by $r(I\otimes M_P)V$ and encode its output
    with both flags zero. Define the joining cost by the actual
    dilation and attenuation costs, including their routing, followed
    by the two selected-zero reflections and the amplification rounds.
\end{definition}

\begin{theorem}[An actual exact circuit for the encoded joining step]
    \label{thm:mpu_encoded_merging}
    \lean{MPUCircuit.embedOp_suffix_eq_reindex,
        MPUCircuit.initializedBasisMatrix_appendBasisEmbedding,
        MPUCircuit.isCleanImplementation_embedOp_suffix,
        MPUCircuit.embedOp_prefix_eq_reindex,
        MPUCircuit.isCleanImplementation_embedOp_prefix,
        MPUCircuit.compatibleBondDilationEmbedding_eq_append,
        MPUCircuit.isCleanImplementation_oneSiteBasis,
        MPUCircuit.isCleanImplementation_compatible_attenuation,
        MPUCircuit.exists_compatible_merging_preparation_circuit,
        MPUCircuit.exists_placed_compatible_merging_preparation_circuit,
        MPUCircuit.joiningSuccessProjection_merging_preparation,
        MPUCircuit.diagonal_mul_initializedBasisMatrix,
        MPUCircuit.appendBasisEmbedding_merging_flags,
        MPUCircuit.mergingFlagProjection_mul_encoded_basis,
        MPUCircuit.encoded_merging_preparation_success,
        MPUCircuit.isCleanImplementation_zeroWorkspace_empty,
        MPUCircuit.exists_clean_encoded_merging_preparation,
        MPUCircuit.success_gram_of_encoded_columns,
        MPUCircuit.exists_encoded_uniform_merging_circuit}
    \leanok
    Let $d\geq2$, $r>0$, and $P>0$. Encode each bond label in $q$
    qudits with its distinguished label represented by physical zero.
    Suppose that two child circuits implement full logical unitaries
    with the same initialized workspace, that their initialized
    product columns equal the prescribed encoding of $V$, and that
    $r(I\otimes M_P)V$ is an isometry. Let $b$ be the number of flags
    initialized at the parent input and let the common workspace
    contain at least $\max(b,2)$ qudits.
    Then an actual pair-gate circuit implements the encoded normalized
    parent exactly on those input columns. It also implements a full
    logical unitary, and its adjoint implements the adjoint logical
    unitary with the same workspace returned exactly to zero on every
    logical input. If $K_1,K_2$ are the child costs and $M$ is the
    joining cost, its gate count is
    \[
      \ell\bigl(1+2(K_1+K_2+M)+R_b+R_2\bigr)+K_1+K_2+M,
      \qquad \ell=\lceil\pi r/4\rceil,
    \]
    where $R_b,R_2$ are the actual common-workspace reflection costs.
    The success probability is derived from the normalized parent
    isometry; it is not an additional assumption.
    The same full logical unitary is supported in any set containing
    both child supports, the initialized input flags, and the joining
    packet, including on padded logical states.
\end{theorem}
\begin{proof}\leanok
    Construct the bond dilation and the attenuation rotation on the
    prescribed product basis. Their successful physical branch is
    $\sin(\pi/(4\ell+2))$ times the encoded normalized parent.
    Its isometry identity gives the constant squared success amplitude.
    Construct both reflections by testing the selected initialized
    flags in the common workspace. Exact amplification gives the parent
    columns with their full phase, and the all-input clean identities
    justify every inverse child call.
    The preparation acts only on the joining packet. Conjugating the
    initialized reflection by that supported preparation and the child
    unitaries, then multiplying the two reflections and taking powers,
    preserves the claimed logical support.
\end{proof}

This joining theorem is an induction step with explicit child columns
and parent isometry as hypotheses. Identification of those columns with
arbitrary MPU intervals, and the complete recursive existence statement,
remain separate assertions.

\begin{definition}[The initialized auxiliaries of an interval]
    \label{def:mpu_interval_initialization}
    \lean{MPUCircuit.intervalAuxiliaryInitializationSupport,
        MPUCircuit.logicalSupportSites}
    \leanok
    The auxiliaries initialized at an interval input are the auxiliary
    sites in its logical support, including both boundary bond
    registers. Enumerate any finite logical support by a finite site
    injection, using the fixed consecutive ordering of logical sites.
\end{definition}

\begin{theorem}[Every interval flag selection fits the common workspace]
    \label{thm:mpu_interval_initialization}
    \lean{MPUCircuit.card_logicalAuxiliarySupport_le,
        MPUCircuit.intervalInteriorAuxiliarySupport_subset_auxiliary_range,
        MPUCircuit.joiningAuxiliarySupport_subset_auxiliary_range,
        MPUCircuit.intervalAuxiliaryInitializationSupport_subset_auxiliary_range,
        MPUCircuit.card_intervalInteriorAuxiliarySupport_le,
        MPUCircuit.card_joiningAuxiliarySupport_le,
        MPUCircuit.card_intervalAuxiliaryInitializationSupport_le,
        MPUCircuit.logicalSupportSites_range,
        MPUCircuit.logicalSupportSites_initialized_iff}
    \leanok
    Every set of logical auxiliary sites has cardinality at most the
    total number of allocated auxiliaries. In particular, the interior
    auxiliary sites of an interval, its two joining bond registers and
    two joining flags, and all auxiliaries initialized at its input
    fit in the common workspace. Their finite selections are
    constructed, and testing the selected consecutive coordinates
    for zero is equivalent to testing the original logical support.
\end{theorem}
\begin{proof}\leanok
    Each support is contained in the image of the auxiliary-site
    injection. Inject it into the full auxiliary register and compare
    cardinalities. A finite enumeration transported through the logical
    coordinate equivalence gives the required site selection and its
    exact zero-test identity.
\end{proof}

\begin{definition}[Joint configurations for disjoint child intervals]
    \label{def:mpu_joint_child_columns}
    \lean{MPUCircuit.AgreeOffPair,
        MPUCircuit.OutsidePlacedConfig,
        MPUCircuit.placedConfigEquiv,
        MPUCircuit.jointPlacedConfigEquiv,
        MPUCircuit.jointPlacedBasisEmbedding}
    \leanok
    Decompose a complete logical configuration into its two disjoint
    child packets and its remaining sites. Concatenate prescribed
    input or output basis encodings on the child packets with a
    prescribed spectator encoding on the remaining sites.
\end{definition}

\begin{theorem}[Joint child columns follow from individual columns]
    \label{thm:mpu_joint_child_columns}
    \lean{MPUCircuit.mul_embedOp_disjoint_apply,
        MPUCircuit.placedConfigEquiv_fst,
        MPUCircuit.placedConfigEquiv_snd,
        MPUCircuit.jointPlacedConfigEquiv_left,
        MPUCircuit.jointPlacedConfigEquiv_right,
        MPUCircuit.jointPlacedConfigEquiv_snd,
        MPUCircuit.agreeOffPair_iff_outside_eq,
        MPUCircuit.reindex_mul_embedOp_of_disjoint,
        MPUCircuit.initializedBasisMatrix_jointPlacedBasisEmbedding,
        MPUCircuit.jointPlacedColumns_of_individual}
    \leanok
    Suppose the two child packets are disjoint and their individual
    column identities are $XE_g=E_{g'}V$ and $YE_h=E_{h'}W$.
    Placing the child operators in the complete register gives the
    joint column identity with output map $V\otimes W\otimes I$,
    for every encoded spectator state. No extra joint column identity
    is assumed. In these coordinates the product of the placed
    operators is exactly $X\otimes Y\otimes I$.
\end{theorem}
\begin{proof}\leanok
    Expand the matrix product of the two disjoint placed operators.
    Agreement of the remaining coordinates reduces the sum to the two
    independent child matrix entries. Relabel the configuration by the
    two packets and the remaining sites, and apply the two individual
    column identities inside their tensor product.
\end{proof}

\begin{definition}[Uniform numerical budgets]
    \label{def:mpu_uniform_numeric_resources}
    \lean{MPUCircuit.uniformPacketGateCount,
        MPUCircuit.uniformNodeGateCount,
        MPUCircuit.uniformReflectionGateCount,
        MPUCircuit.uniformCircuitResourceCoefficient,
        MPUCircuit.uniformCircuitTreeGateCount}
    \leanok
    Put $q=\lceil\log_dD\rceil$, let $g$ be the complete register size,
    and let $A$ be the common workspace size. A common packet budget is
    \[
      M_0=\bigl(76d^{24}D^{12}+8(D+1)\bigr)g.
    \]
    Reserve an additional $16g^2$ for coordinate permutations, and use
    $M=M_0+16g^2$ as a common leaf and node overhead. The common
    reflection budget is $F=1+K_d(A)\,2g$, where
    $K_d(a)=8(38d^{18}+1)a^2+4a$.
    Define the midpoint-tree count using leaf cost $M$, joining cost
    $M$, and two reflection costs $F$. Finally put
    \[
    c(d,D)=(2D+2)\bigl(5(76d^{24}D^{12}+8(D+1))+400\bigr)
          +D(30400d^{18}+803).
    \]
\end{definition}

\begin{theorem}[An explicit uniform numerical resource bound]
    \label{thm:mpu_uniform_numeric_resources}
    \lean{MPUCircuit.bondRegisterWidth_le_bondDim,
        MPUCircuit.placedLeafGateCount_le_uniform,
        MPUCircuit.joiningPacketGateCount_le_uniform,
        MPUCircuit.selectedZeroRegisterReflectionPoolGateCount_le_uniform,
        MPUCircuit.uniformReflectionGateCount_le_cubic,
        MPUCircuit.uniformCircuitTreeGateCount_le_polynomial,
        MPUCircuit.uniformCircuitTreeGateCount_le_polynomial_bond}
    \leanok
    The placed leaf budget and the stated routed joining-packet budget
    are at most $M_0$ for $d\geq2$ and $D>0$. Every selected reflection
    testing at most $A$ flags, including the empty selection, has cost
    at most $F$. For $N>0$, the latter satisfies
    \[
      F\leq(15200d^{18}+401)\bigl(N(q+1)\bigr)^3.
    \]
    If every node uses at most $D$ amplification rounds, the stated
    numerical tree count satisfies
    \[
      T\leq c(d,D)(D+1)^3
          (2N)^{\lceil\log_2D\rceil+6}.
    \]
    The coefficient $c(d,D)(D+1)^3$ is polynomial in $D$ for fixed
    $d$, of degree at most sixteen. This theorem bounds the numerical
    recursion and does not assume complete recursive circuit existence.
\end{theorem}
\begin{proof}\leanok
    The capacity bound gives joining-packet dimension at most $d^4D^2$
    and $q\leq D$. Monotonicity of $K_d$ gives the uniform reflection
    bound. The allocation gives $A\leq2N(q+1)$ and $g\leq5N(q+1)$.
    Thus every node overhead is at most a coefficient depending only
    on $d,D$ times $\bigl(N(q+1)\bigr)^3$. Substitute these estimates
    in the actual balanced-tree recurrence and use
    $\lceil\log_2(4D+2)\rceil\leq\lceil\log_2D\rceil+3$.
\end{proof}

\begin{theorem}[Selected-zero reflections act only on their flags]
    \label{thm:mpu_selected_reflection_support}
    \lean{MPSPreparation.embedOp_diagonal,
        MPSPreparation.selectedZeroRegisterLogical_eq_embedOp,
        MPSPreparation.selectedZeroRegisterLogical_mem_supportedOperators,
        MPSPreparation.selectedZeroRegisterLogical_mem_supportedOperators_of_range_subset}
    \leanok
    The reflection which changes sign precisely when a selected set
    of qudits is zero is supported on those qudits. All states of
    the remaining logical qudits are allowed. This includes the empty
    selection, whose reflection is the scalar minus identity.
\end{theorem}
\begin{proof}\leanok
    A diagonal operator placed on a selected packet depends only on
    the selected configuration, and has a Kronecker delta on the
    remaining coordinates. Apply this identity to the local diagonal
    all-zero reflection.
\end{proof}

\begin{theorem}[Exact amplification preserves interval support]
    \label{thm:mpu_amplification_support}
    \lean{MPSPreparation.subspaceReflection_mem_supportedOperators,
        MPSPreparation.pow_mem_supportedOperators,
        MPSPreparation.subspaceReflection_image_mem_supportedOperators,
        MPSPreparation.postselectionAmplificationStep_mem_supportedOperators,
        MPSPreparation.postselectionAmplification_mem_supportedOperators,
        MPSPreparation.postselectionAmplification_mem_supportedOperators_of_projection_support}
    \leanok
    Suppose $Z$ is a unitary supported on $S$, and the reflections
    about the initialized image $EE^\dagger$ and the successful
    projection $P$ are also supported on $S$. Then the full logical
    amplification operator
    \[
      \bigl[-(I-2ZEE^\dagger Z^\dagger)(I-2P)\bigr]^\ell Z
    \]
    is supported on $S$, with its exact scalar phase retained.
    Support of the two projections themselves suffices to derive
    support of their reflections. No initialization of the remaining
    logical sites is imposed.
\end{theorem}
\begin{proof}\leanok
    Conjugating the initialized reflection by $Z$ gives the displayed
    image reflection. Operators supported on $S$ are closed under
    adjoints, products, powers, addition, and scalar multiplication,
    including multiplication by minus one.
\end{proof}

\begin{definition}[Canonical input and output encodings of an interval]
    \label{def:mpu_interval_column_encoding}
    \lean{MPUCircuit.intervalPacketCard,
        MPUCircuit.intervalPacketSites,
        MPUCircuit.intervalPacketLogicalSite,
        MPUCircuit.intervalPhysicalZeroConfig,
        MPUCircuit.intervalInputEmbedding,
        MPUCircuit.intervalBoundaryConfig,
        MPUCircuit.intervalOutputEmbedding}
    \leanok
    Enumerate an interval's logical support as its physical packet.
    Its input encoding records the physical input configuration and
    sets every auxiliary to zero. Its output encoding records its
    physical output, right bond label, and left bond label, using
    prescribed injective encodings in the two boundary registers and
    setting every strict-interior auxiliary to zero.
\end{definition}

\begin{theorem}[Actual interval encodings preserve every recorded coordinate]
    \label{thm:mpu_interval_column_encoding}
    \lean{MPUCircuit.intervalPacketSites_range,
        MPUCircuit.intervalPacketSites_disjoint,
        MPUCircuit.intervalPacketLogicalSite_mem,
        MPUCircuit.intervalPacketLogicalSite_surjective,
        MPUCircuit.intervalPhysicalZeroConfig_physical,
        MPUCircuit.intervalBoundaryConfig_right,
        MPUCircuit.intervalBoundaryConfig_left,
        MPUCircuit.intervalBoundaryConfig_physical,
        MPUCircuit.intervalBoundaryConfig_interior}
    \leanok
    The input and output maps are injective and agree with all of
    their recorded physical and boundary coordinates. The output
    has every strict-interior auxiliary exactly zero. Both maps
    include the zero-width endpoint registers. Consecutive child
    packets are disjoint and their site enumerations have precisely
    the prescribed interval supports.
\end{theorem}
\begin{proof}\leanok
    Read the physical and two boundary registers to recover the
    encoded coordinates. The two boundary embeddings have disjoint
    images, and the remaining auxiliary entries are set to zero.
    The interval-support identities give the enumeration and
    disjointness statements.
\end{proof}

\begin{definition}[Physical configurations split at the joining cut]
    \label{def:mpu_actual_parent_split}
    \lean{MPUCircuit.cutIntervalSplitEquiv}
    \leanok
    Split the physical configuration of an interval $[j,k)$ into
    its two consecutive restrictions to $[j,m)$ and $[m,k)$.
\end{definition}

\begin{theorem}[The computed joining parent is the actual MPU interval]
    \label{thm:mpu_actual_parent_isometry}
    \lean{MPUCircuit.normalizedJoiningParent_balancedIntervalChildren,
        MPUCircuit.weightedMinimalOperatorInterval_comp_eq_reindex,
        MPUCircuit.normalizedJoiningParent_minimalOperatorIntervals,
        MPUCircuit.normalizedJoiningParent_minimalOperatorIntervals_isIsometry,
        MPUCircuit.exists_normalized_joining_parent_isometries_of_unitary}
    \leanok
    The normalized joining contraction of the two balanced adjacent
    intervals equals the actual weighted parent interval, after
    splitting its physical coordinates, with the joining label pair
    reset to zero. Consequently the computed parent is an isometry
    whenever the left cut metric belongs to its prefix Gram affine
    space and the right dual metric normalizes its affine space.
    For every finite unitary on $N>0$ sites in positive physical
    dimension, one endpoint-normalized minimal basis and coherent
    positive metric family satisfying these conditions exists, and
    all its joining parents are isometries simultaneously. No
    parent-isometry hypothesis is supplied in this existence statement.
\end{theorem}
\begin{proof}\leanok
    The normalized bond reset cancels the two weighted metrics at
    the common cut. Multiplication of the actual minimal interval
    factors gives the parent factor. Relabel its physical coordinates
    by the two consecutive children and apply the previously derived
    weighted-interval isometry identity. The simultaneous normalized
    metric family follows from global unitarity.
\end{proof}

\begin{definition}[Placed basis encodings with spectators]
    \label{def:mpu_supported_child_basis}
    \lean{MPUCircuit.placedBasisEmbedding}
    \leanok
    Encode a child packet together with a basis encoding on its
    remaining logical sites. Taking the identity basis encoding
    on those sites allows every outside logical configuration.
\end{definition}

\begin{theorem}[Supported global children supply their joint columns]
    \label{thm:mpu_supported_child_columns}
    \lean{MPUCircuit.initializedBasisMatrix_placedBasisEmbedding,
        MPUCircuit.reindex_embedOp_placedConfigEquiv,
        MPUCircuit.placedColumns_iff_localColumns,
        MPUCircuit.jointPlacedColumns_of_supported_individual}
    \leanok
    Let two global operators have disjoint child supports. Their
    individual initialized-column identities, each retaining every
    outside logical configuration, imply their joint column identity
    with output map $V\otimes W\otimes I$ for every spectator encoding.
    Local child matrices and the joint column identity are derived;
    neither is an additional hypothesis.
\end{theorem}
\begin{proof}\leanok
    Support supplies local matrices. Relabel the complete configuration
    by a child packet and its complement; the global identity is then
    the local identity tensored with an identity on that complement.
    Recover the local identity by taking a fixed diagonal spectator
    entry, and apply the disjoint-child product-column theorem.
\end{proof}

\begin{definition}[Logical site permutations fixing the workspace]
    \label{def:mpu_workspace_fixed_site_permutation}
    \lean{MPSPreparation.workspaceFixedSitePerm}
    \leanok
    Extend a permutation of the logical sites by the identity on
    the appended workspace sites, transporting the disjoint sum to
    the consecutive site coordinates.
\end{definition}

\begin{theorem}[Coordinate regrouping has additive circuit cost]
    \label{thm:mpu_quantitative_site_permutation}
    \lean{Equiv.Perm.swapFactorsAux_length_le,
        Equiv.Perm.swapFactors_length_le_card,
        MPSPreparation.permOp_list_prod,
        MPSPreparation.isPairProduct_permOp,
        MPSPreparation.IsPairProduct.conjugate_permOp,
        MPSPreparation.workspaceFixedSitePerm_castAdd,
        MPSPreparation.workspaceFixedSitePerm_natAdd,
        MPSPreparation.isCleanImplementation_permOp_workspaceFixedSitePerm,
        MPSPreparation.IsCleanImplementation.conjugate_workspaceFixedSitePerm,
        MPSPreparation.isPairProduct_isCleanImplementation_conjugate_workspaceFixedSitePerm}
    \leanok
    Every permutation of $n$ qudit sites has an actual neighboring
    pair-gate circuit of at most $2n^2$ gates. Conjugating a circuit
    of cost $K$ by that permutation gives cost at most $K+4n^2$,
    including its exact scalar phase. A logical site permutation
    fixing the appended workspace preserves every all-input clean
    implementation identity. For a register of total size $n+A$,
    this conjugation costs at most $K+4(n+A)^2$.
\end{theorem}
\begin{proof}\leanok
    The finite permutation factors into at most $n$ transpositions,
    each routed with at most $2n$ neighboring swaps. Use the
    permutation circuit and its adjoint on the two sides of the
    original circuit. The extended permutation changes only logical
    coordinates and fixes each initialized workspace site, which
    gives the exact clean identity under conjugation.
\end{proof}

\begin{definition}[Actual interval columns]
    \label{def:mpu_actual_interval_columns}
    \lean{MPUCircuit.minimalIntervalInputEmbedding}
    \lean{MPUCircuit.minimalIntervalOutputEmbedding}
    \lean{MPUCircuit.IsMinimalIntervalColumnImplementation}
    \leanok
    Fix the minimal cut bases and one balanced metric family of the
    given unitary. The input of an interval consists of its physical
    configuration with all interval auxiliaries initialized to zero,
    together with an arbitrary configuration outside the interval.
    Its output consists of the weighted interval amplitudes in the
    prescribed cut-label encodings, with all strictly interior
    auxiliaries zero and the outside configuration unchanged.
    An interval implementation satisfies this complete column identity.
\end{definition}

\begin{theorem}[Joint columns from actual child implementations]
    \label{thm:mpu_actual_interval_child_columns}
    \lean{MPUCircuit.minimalIntervalChildColumns_of_supported_individual}
    \leanok
    For adjacent nonempty intervals, suppose each logical operator
    has its interval support and satisfies the actual interval column
    identity. Their product has the tensor product of the two weighted
    child maps as its initialized columns, with the identity on every
    configuration outside their combined support. No joint column
    identity is an additional assumption.
\end{theorem}
\begin{proof}\leanok
    The child supports are disjoint. Their individual column identities
    therefore determine the columns of their product by finite matrix
    multiplication. The prescribed cut-label encodings occur in the
    resulting product without a further change of bond basis.
\end{proof}

\begin{definition}[Actual interval circuit implementations]
    \label{def:mpu_actual_interval_circuit}
    \lean{MPUCircuit.IsMinimalIntervalCircuitImplementation}
    \leanok
    Fix one minimal cut basis and balanced metric family. An interval
    circuit implementation of budget $K$ consists of a full logical
    unitary supported on the interval and a neighboring-pair circuit
    of at most $K$ gates. The circuit implements that logical unitary
    with the common workspace returned to zero on every logical
    input. Its initialized interval columns are the actual weighted
    interval columns, including arbitrary outside configurations.
\end{definition}

\begin{theorem}[Interval recursion from actual leaf and merging circuits]
    \label{thm:mpu_actual_interval_recursion}
    \lean{MPUCircuit.IsMinimalIntervalCircuitImplementation.mono,
        MPUCircuit.exists_minimalInterval_circuit_of_partition}
    \leanok
    Suppose single-site implementations of budget $K$ and an
    adjacent-interval merging rule have been established for the same
    fixed basis and metric family. If that rule gives budget
    $b(K_1+K_2)+M$ from child budgets $K_1,K_2$, then every consecutive
    interval partition has an actual circuit with exactly its numerical
    tree budget. All logical support, workspace cleanup, and actual
    interval column identities are retained. Increasing a budget
    preserves these identities.
\end{theorem}
\begin{proof}\leanok
    Induct on the consecutive interval partition. Its leaves use the
    given single-site circuits. Each node uses the given merging rule
    on the two actual child implementations, whose endpoint cuts
    coincide. The resulting gate budget is the defining numerical
    tree recursion. This conditional statement requires the leaf and
    merging constructions to be established separately.
\end{proof}

\begin{definition}[The ordered joining packet]
    \label{def:mpu_actual_joining_packet}
    \lean{MPUCircuit.joiningFlagPacketSites,
        MPUCircuit.joiningBondPacketSites,
        MPUCircuit.joiningPacketLogicalSites,
        MPUCircuit.joiningPacketSites}
    \leanok
    At an internal cut, order the joining packet as the left child's
    right bond register, the right child's left bond register, the
    dilation flag, and the attenuation flag. Its consecutive site
    inclusion is determined by the allocated logical registers.
\end{definition}

\begin{theorem}[Coordinates of the actual joining packet]
    \label{thm:mpu_actual_joining_packet}
    \lean{MPUCircuit.joiningBondPacketSites_disjoint_flags,
        MPUCircuit.joiningPacketLogicalSites_left,
        MPUCircuit.joiningPacketLogicalSites_right,
        MPUCircuit.joiningPacketLogicalSites_flag,
        MPUCircuit.exists_joiningPacket_suffix_coordinates}
    \leanok
    The ordered packet inclusion is injective and has the prescribed
    restrictions to both bond registers and both flags. There is a
    permutation of the full logical register sending its final packet
    to precisely these actual sites in this order. Both the remaining
    prefix size and this permutation are derived from the site
    inclusions; neither is an additional geometric assumption.
\end{theorem}
\begin{proof}\leanok
    The allocated bond and flag sites are disjoint. Their injective
    enumerations give the packet inclusion. Its cardinality does not
    exceed the logical register size. Extend the correspondence
    between the final packet inclusion and the actual packet inclusion
    to a permutation of the finite site set.
\end{proof}

\begin{definition}[Prefix of a product-compatible basis encoding]
    \label{def:mpu_prefix_basis_encoding}
    \lean{MPUCircuit.prefixBasisEmbedding}
    \leanok
    Given an injective full basis encoding with a prescribed suffix
    depending only on its bond label, fix one bond label and restrict
    the full encoding to the prefix registers. This defines the
    prefix basis inclusion.
\end{definition}

\begin{theorem}[Factorization of a product-compatible basis encoding]
    \label{thm:mpu_prefix_basis_encoding}
    \lean{MPUCircuit.appendBasisEmbedding_prefixBasisEmbedding}
    \leanok
    Suppose the suffix of the full encoding is the prescribed bond
    encoding, and its prefix is independent of the bond label.
    Then the prefix inclusion is injective, and its product with the
    prescribed suffix is exactly the full encoding.
\end{theorem}
\begin{proof}\leanok
    Prefix equality, together with equality of the fixed suffix,
    gives equality of full encoded configurations. Full injectivity
    therefore gives prefix injectivity. Restriction to the two
    coordinate blocks proves the product identity.
\end{proof}

\begin{definition}[Joint output coordinates of adjacent intervals]
    \label{def:mpu_actual_joint_output_encoding}
    \lean{MPUCircuit.initializedComplementEmbedding,
        MPUCircuit.intervalJointOutputEmbedding}
    \leanok
    Combine the prescribed child output inclusions on their disjoint
    interval supports with an arbitrary outside basis inclusion.
    Setting a selected set of outside coordinates to zero gives
    its initialized complement inclusion.
\end{definition}

\begin{theorem}[Restrictions of the actual joint output encoding]
    \label{thm:mpu_actual_joint_output_encoding}
    \lean{MPUCircuit.jointPlacedBasisEmbedding_left,
        MPUCircuit.jointPlacedBasisEmbedding_right,
        MPUCircuit.jointPlacedBasisEmbedding_outside,
        MPUCircuit.initializedComplementEmbedding_zero,
        MPUCircuit.intervalJointOutputEmbedding_left_support,
        MPUCircuit.intervalJointOutputEmbedding_right_support,
        MPUCircuit.intervalBoundaryConfig_right_eq_away,
        MPUCircuit.intervalBoundaryConfig_left_eq_away,
        MPUCircuit.intervalJointOutputEmbedding_eq_away_joining_bonds}
    \leanok
    The joint output restricts to the prescribed left and right
    child configurations and to its prescribed outside configuration.
    Changing the two joining bond labels affects only their two
    allocated bond registers. Every other coordinate, including
    initialized interior auxiliaries and outside spectators, is
    independent of those labels.
\end{theorem}
\begin{proof}\leanok
    Restrict the joint product inclusion to its disjoint coordinate
    blocks. The interval boundary configuration reads a bond label
    only at its allocated boundary register. The stated independence
    follows by distinguishing those boundary sites from all others.
\end{proof}

\begin{definition}[Leaf coordinates of the actual interval]
    \label{def:mpu_actual_leaf_coordinates}
    \lean{MPUCircuit.leafIntervalSiteEquiv,
        MPUCircuit.leafIntervalConfigEquiv}
    \leanok
    For a one-site interval, identify its consecutive packet with
    the physical site followed by its right and left boundary bond
    registers. The induced configuration bijection preserves the
    prescribed physical and boundary coordinates.
\end{definition}

\begin{theorem}[Actual interval columns of the existing leaf circuits]
    \label{thm:mpu_actual_leaf_columns}
    \lean{MPUCircuit.intervalPacketLogicalSite_leafIntervalSiteEquiv,
        MPUCircuit.leafConsecutiveSites_eq_comp_intervalPacketSites,
        MPUCircuit.embedOp_reindex_siteEquiv,
        MPUCircuit.leafIntervalConfigEquiv_leafInput,
        MPUCircuit.leafIntervalConfigEquiv_leafOutput,
        MPUCircuit.initializedBasisMatrix_reindex_of_mapping,
        MPUCircuit.leafColumns_reindex_interval,
        MPUCircuit.minimalLeafInterval_reindex,
        MPUCircuit.isMinimalIntervalColumnImplementation_leaf_of_columns,
        MPUCircuit.exists_normalized_minimalInterval_family_with_placed_leaf_columns_of_unitary}
    \leanok
    For $d\geq2$ and $N\geq2$, a finite unitary of consecutive cut
    ranks at most $D$ admits one normalized minimal basis and metric
    family whose placed leaf circuits satisfy the actual interval
    column identities. The previously constructed circuits and logical
    operators remain exactly the same, with their gate counts, support,
    and all-input workspace cleanup. The outside configurations are
    arbitrary. No separate interval column identity is assumed.
\end{theorem}
\begin{proof}\leanok
    The packet coordinate bijection fixes the actual physical and
    boundary sites. Transport the input and output basis inclusions
    through it, together with the weighted single-site map. Reindexing
    the local operator commutes with its placement in the logical
    register. Thus the existing local columns give the actual interval
    columns without a further gate or change of unitary extension.
\end{proof}

\begin{definition}[Physical inclusions at the root]
    \label{def:mpu_actual_root_inclusions}
    \lean{MPUCircuit.physicalLogicalConfigEmbedding,
        MPUCircuit.logicalGlobalConfigEmbedding,
        MPUCircuit.physicalGlobalConfigEmbedding}
    \leanok
    Include the physical chain into the logical register with its
    auxiliaries zero. Include the logical register into the complete
    register with the common scratch pool zero. Their composition
    includes the physical chain with every auxiliary zero.
\end{definition}

\begin{theorem}[The actual root gives the physical unitary]
    \label{thm:mpu_actual_root_columns}
    \lean{MPUCircuit.intervalPacketSites_full_range,
        MPUCircuit.outsidePlacedConfig_full_subsingleton,
        MPUCircuit.minimalIntervalInputEmbedding_full_apply,
        MPUCircuit.minimalIntervalOutputEmbedding_full_apply,
        MPUCircuit.minimalIntervalRootColumns_eq_physical_of_fullInterval,
        MPUCircuit.minimalIntervalRootColumns_eq_physical,
        MPUCircuit.physicalGlobalConfigEmbedding_apply,
        MPUCircuit.initializedBasisMatrix_physicalGlobalConfigEmbedding,
        MPUCircuit.physicalGlobalColumns_of_clean_logicalColumns,
        MPUCircuit.minimalIntervalRootColumns_eq_physical_global}
    \leanok
    Let a logical operator satisfy the actual full-interval column
    identity of the given finite unitary $U$. Normalize the minimal
    endpoint bases and take the endpoint metrics in their actual
    prefix Gram spaces. If a complete circuit $C$ implements that
    logical operator with the common scratch pool returned to zero
    on every logical input, then
    \[
        CJ=JU,
    \]
    where $J$ initializes every logical auxiliary and scratch qudit.
    The physical output and the scalar phase follow from the given
    unitary; they are not additional assumptions.
\end{theorem}
\begin{proof}\leanok
    The full packet contains every logical site, so its outside
    configuration is unique. Both endpoint bond registers have width
    zero. The actual endpoint Gram spaces and normalized bases make
    the full weighted interval equal to $U$. The two root basis
    inclusions therefore become the physical inclusion with logical
    auxiliaries zero. Compose its column identity with the full
    logical clean-workspace identity to obtain the stated equality.
\end{proof}

\begin{definition}[Complete circuit registers]
    \label{def:mpu_complete_register}
    \lean{MPUCircuit.circuitPaddingCount,
        MPUCircuit.completeCircuitSiteCount,
        MPUCircuit.completePhysicalConfigEmbedding}
    \leanok
Append one padding qudit when $N=1$, and none otherwise, to the physical register,
logical auxiliaries, and reusable scratch pool. The physical basis inclusion
initializes all three groups to zero.
\end{definition}

\begin{theorem}[Complete circuit registers]
    \label{thm:mpu_complete_register}
    \lean{MPUCircuit.circuitPaddingCount_le_one,
        MPUCircuit.completeCircuitSiteCount_le,
        MPUCircuit.completePhysicalConfigEmbedding_apply}
    \leanok
The complete register has at most $5N(q+1)+1$ sites, where $q$ is the
bond-register width. Its physical basis inclusion is the physical configuration
followed by zeros on the logical auxiliaries, scratch pool, and padding.
\end{theorem}
\begin{proof}\leanok
The padding count is at most one. Apply the common register bound and expand
the three successive initialized inclusions.
\end{proof}

\begin{definition}[Initialized joining flags]
    \label{def:mpu_joint_flags}
    \lean{MPUCircuit.intervalJointOutsideFlagSet,
        MPUCircuit.intervalJointInitializedOutsideEmbedding}
    \leanok
View the two joining flags as sites outside both child packets. Initialize exactly
these flags and retain arbitrary configurations on all other outside sites.
\end{definition}

\begin{theorem}[Initialized joining flags]
    \label{thm:mpu_joint_flags}
    \lean{MPUCircuit.joiningFlag_outside_childPackets,
        MPUCircuit.intervalJointOutputEmbedding_joiningFlag_zero}
    \leanok
The joining flags lie outside both child packets. The joint output inclusion,
with these outside flags initialized, has both joining flags equal to zero.
\end{theorem}
\begin{proof}\leanok
The joining flag sites are disjoint from both interval packets. Their outside
inclusion sets exactly these two coordinates to zero.
\end{proof}

\begin{theorem}[The actual active joining packet]
    \label{thm:mpu_actual_packet_encoding}
    \lean{MPUCircuit.intervalJointOutputEmbedding_joiningPacket_active}
    \leanok
The actual joint child output restricted to the joining packet is precisely
the active compatible bond code, with the shared right and left labels in
contraction order and both flags zero.
\end{theorem}
\begin{proof}\leanok
Evaluate the prescribed boundary configurations on the two shared bond
registers, then use the initialized joining-flag identities.
\end{proof}

\begin{definition}[Derived product coordinates]
    \label{def:mpu_actual_suffix_factorization}
    \lean{MPUCircuit.outerSpectatorJoiningRegrouping,
        MPUCircuit.pullbackBasisEmbedding,
        MPUCircuit.regroupedIntervalJointOutputEmbedding}
    \leanok
Regroup the physical outputs, outer cut labels, and outside configurations
before the two joining labels. Pull a basis inclusion back along a site
bijection to place the joining packet at the end.
\end{definition}

\begin{theorem}[Derived product coordinates]
    \label{thm:mpu_actual_suffix_factorization}
    \lean{MPUCircuit.pullbackBasisEmbedding_apply,
        MPUCircuit.suffixCoordinates_prefix_not_mem,
        MPUCircuit.joiningPacketLogicalSites_range_left,
        MPUCircuit.joiningPacketLogicalSites_range_right,
        MPUCircuit.regroupedIntervalJointOutputEmbedding_apply,
        MPUCircuit.exists_regroupedIntervalJointOutput_product_coordinates}
    \leanok
The actual joint output admits a derived site bijection and a product code:
a prefix encoding independent of the joining labels, followed by the compatible
joining-packet code. The prefix encoding is injective. Neither the site
bijection nor its factorization is supplied as a hypothesis.
\end{theorem}
\begin{proof}\leanok
Use the actual ordered joining packet and its complementary site bijection.
Every prefix site is independent of the joining labels. Injectivity of the full
joint output code therefore implies injectivity of its prefix factor.
\end{proof}

\begin{theorem}[The encoded merger in original coordinates]
    \label{thm:mpu_permuted_encoded_merger}
    \lean{MPSPreparation.permOp_conjTranspose,
        MPSPreparation.permOp_conjTranspose_mul_mul,
        MPSPreparation.permOp_conjugate_mem_supportedOperators,
        MPSPreparation.permOp_mul_rectangular,
        MPSPreparation.conjugate_workspaceFixedSitePerm_columns,
        MPUCircuit.exists_permuted_encoded_uniform_merging_circuit}
    \leanok
Suppose two full clean child circuits have the stated encoded joint columns,
logical support inside a prescribed set, and an isometric normalized parent.
A site bijection putting the joining packet last yields a merger in the original
coordinates. Its logical support, initialized columns, and scalar phase are
preserved. Each child conjugation and the final return costs at most $4g^2$
neighboring-pair gates, where $g$ is the complete logical and scratch site count.
\end{theorem}
\begin{proof}\leanok
Conjugate both child circuits into the suffix coordinates. Apply the supported
encoded merger, then conjugate its logical and physical operators back.
The site permutation fixes the scratch pool. Its basis action transports the
rectangular initialized columns exactly.
\end{proof}

\begin{theorem}[Complete merger gate bound]
    \label{thm:mpu_actual_merger_resources}
    \lean{MPUCircuit.permutedEncodedUniformMergingGateCount_le_uniform}
    \leanok
For $N\geq2$, $0<r\leq D$, and a selected initial flag count within the allocated
pool, the encoded merger including both child conjugations and its final return
has gate count at most
\[
 (2D+1)(K_1+K_2)+(2D+1)M+D(1+2F),
\]
where $M$ is the common node budget and $F$ is the common reflection budget.
The quadratic routing reserve is included in $M$.
\end{theorem}
\begin{proof}\leanok
Bound the joining preparation by the common packet budget, both reflections
by $F$, and the prescribed number of amplification rounds by $r\leq D$.
Expand the gate formula. The node reserve absorbs the three conjugations.
\end{proof}

\begin{definition}[Joining with arbitrary spectator states]
    \label{def:mpu_joining_spectators}
    \lean{MPUCircuit.joiningSpectatorColumns}
    \leanok
Tensor the two weighted interval columns with the identity on the remaining
outside configurations, and move those configurations before the joining labels.
\end{definition}

\begin{theorem}[Joining with arbitrary spectator states]
    \label{thm:mpu_joining_spectators}
    \lean{MPUCircuit.joiningSpectatorColumns_eq_reindex,
        MPUCircuit.normalizedJoiningParent_joiningSpectatorColumns,
        MPUCircuit.normalizedJoiningParent_joiningSpectatorColumns_isIsometry,
        MPUCircuit.regroupedWeightedChildren_joiningSpectatorColumns}
    \leanok
The normalized joining contraction commutes with the spectator identity after
this regrouping. Thus its isometry follows from the actual normalized parent
isometry. The regrouped child tensor is precisely the actual weighted child
Kronecker tensor, with every spectator coordinate retained.
\end{theorem}
\begin{proof}\leanok
Expand the joining contraction entrywise. The spectator Kronecker delta
passes through its finite sum. Reindex the actual child product and use the
actual parent isometry.
\end{proof}

\begin{definition}[Coordinates of basis inclusions]
    \label{def:mpu_basis_encoding_coordinates}
    \lean{MPUCircuit.basisEncodingEquivOfRangeEq}
    \leanok
Two injective basis encodings with the same image determine a unique bijection
between their label sets, carrying each encoded configuration to itself.
\end{definition}

\begin{theorem}[Coordinates of basis inclusions]
    \label{thm:mpu_basis_encoding_coordinates}
    \lean{MPUCircuit.basisEncodingEquivOfRangeEq_apply,
        MPUCircuit.initializedBasisMatrix_trans,
        MPUCircuit.initializedBasisMatrix_equivOfRangeEq,
        MPUCircuit.initializedBasisMatrix_prodMap,
        MPUCircuit.initializedBasisMatrix_append_joiningChild}
    \leanok
The matrices of initialized basis inclusions respect composition and products.
The derived label bijection transports equal-range encodings exactly. Appending
a compatible active joining code is the corresponding product basis inclusion.
\end{theorem}
\begin{proof}\leanok
Equal images and injectivity determine the label bijection. Matrix
multiplication of the basis inclusions reduces to their Kronecker deltas.
\end{proof}

\begin{theorem}[Transport of initialized flag spaces]
    \label{thm:mpu_initialized_encoding_coordinates}
    \lean{MPUCircuit.mem_range_pullbackBasisEmbedding_iff,
        MPUCircuit.range_pullbackBasisEmbedding_eq_zeroFlagEmbedding}
    \leanok
Pulling an initialized input encoding back along a site bijection transports
its image to the zero-flag inclusion on the transported selected sites.
This gives the complete initialized input subspace, retaining arbitrary
configurations outside the selected flags.
\end{theorem}
\begin{proof}\leanok
Membership in the pulled-back image is membership in the original image
after the site bijection. The zero conditions are precisely its transported
selected sites.
\end{proof}

\begin{definition}[Actual child and parent inputs]
    \label{def:mpu_actual_input_regrouping}
    \lean{MPUCircuit.intervalJointRemainingSiteEquiv,
        MPUCircuit.intervalJointInputConfigEquiv,
        MPUCircuit.intervalJointInputEmbedding}
    \leanok
The remaining sites outside two adjacent child packets and the joining flags
are identified with the sites outside the parent packet. Join the two physical
input configurations and transport the arbitrary remaining configuration along
this site bijection.
\end{definition}

\begin{theorem}[Actual child and parent inputs]
    \label{thm:mpu_actual_input_regrouping}
    \lean{MPUCircuit.mem_range_finEmbeddingAppend_iff,
        MPUCircuit.mem_range_intervalPacketSites_iff,
        MPUCircuit.not_mem_parentPacket_iff,
        MPUCircuit.intervalPhysicalZeroConfig_join_left,
        MPUCircuit.intervalPhysicalZeroConfig_join_right,
        MPUCircuit.intervalPhysicalZeroConfig_joiningFlag_zero,
        MPUCircuit.intervalPacket_restriction_apply,
        MPUCircuit.intervalJointInitializedOutsideEmbedding_parentOutside,
        MPUCircuit.intervalJointInputEmbedding_eq_parent_input,
        MPUCircuit.placedBasisEmbedding_restriction,
        MPUCircuit.minimalIntervalInputEmbedding_zero_auxiliary,
        MPUCircuit.mem_range_minimalIntervalInputEmbedding_iff,
        MPUCircuit.range_minimalIntervalInputEmbedding_eq_zeroFlagEmbedding,
        MPUCircuit.range_intervalJointInputEmbedding_eq_parent,
        MPUCircuit.range_intervalJointInputEmbedding_eq_zeroFlagEmbedding,
        MPUCircuit.minimalIntervalJointInitializedChildColumns_of_supported_individual}
    \leanok
The joint initialized child input is exactly the parent initialized input under
the derived configuration bijection. Their images equal the zero-flag subspace
selected by the parent auxiliary sites. The two actual child column identities,
together with their support conditions, therefore give the joint initialized
child columns; no input bijection or joint-column identity is assumed.
\end{theorem}
\begin{proof}\leanok
The parent support is the union of the two child supports and the joining
flags. Its complement gives the remaining-site bijection. Evaluate the two
input configurations on these disjoint sets. Compose the supported child
identities with the resulting initialized inclusion.
\end{proof}

\begin{theorem}[Actual joining reset gives the parent output]
    \label{thm:mpu_actual_output_reset}
    \lean{MPUCircuit.cutBondSiteEmbedding_mem_intervalInterior,
        MPUCircuit.intervalBoundaryConfig_join_zero_left,
        MPUCircuit.intervalBoundaryConfig_join_zero_right,
        MPUCircuit.intervalJointOutputEmbedding_zero_eq_parent_support,
        MPUCircuit.intervalJointOutputEmbedding_zero_eq_parent}
    \leanok
Suppose the shared cut code sends its distinguished zero label to the zero
bond configuration. The actual joint output with both shared labels reset to
that label is exactly the parent output inclusion, after joining the physical
outputs and transporting the arbitrary remaining configuration. This equality
includes both joining flags and every interior auxiliary.
\end{theorem}
\begin{proof}\leanok
Evaluate on the left child, right child, joining flags, and parent complement.
The reset code makes both shared bond configurations zero, the flags are
initialized to zero, and the outer labels and arbitrary outside configuration
agree with the parent boundary configuration.
\end{proof}

\begin{theorem}[Operator cut ranks of ordinary MPOs]
    \label{thm:mpu_periodic_operator_cut_rank}
    \lean{MPUCircuit.cutRank_traceBoundary_le,
        MPUCircuit.cutRank_mpv_le,
        MPUCircuit.operatorCoefficientTensor_mpo_eq_mpv,
        MPUCircuit.cutRank_operatorCoefficientTensor_mpo_le,
        MPUCircuit.cutRank_openBoundary_le,
        MPUCircuit.cutRank_operatorCoefficientTensor_traceBoundary_le,
        MPUCircuit.cutRank_operatorCoefficientTensor_openBoundary_le,
        MPUCircuit.cutRank_operatorCoefficientTensor_mpo_le_all}
    \leanok
For an ordinary periodic MPO of virtual dimension $\chi$, every consecutive
operator coefficient cut has rank at most $\chi^2$. The same bound holds for an
arbitrary trace boundary. Scalar left and right open boundaries give the bound
$\chi$. These statements require no unitarity, normalization, minimality, or
canonical-form assumption.
\end{theorem}
\begin{proof}\leanok
Split the ordered word at the physical cut. A trace contraction factors
through its two open virtual indices; scalar open boundaries factor through
one. The cut matrix rank is bounded by the dimension of this intermediate
space. For an MPO, its paired-letter coefficient tensor is exactly the
corresponding ordinary matrix-product tensor.
\end{proof}

\begin{theorem}[Transport and cancellation of exact columns]
    \label{thm:mpu_column_coordinate_transport}
    \lean{MPUCircuit.matrix_columns_reindex,
        MPUCircuit.normalizedJoiningParent_submatrix_columns,
        MPSPreparation.IsCleanImplementation.logical_columns_of_ambient}
    \leanok
Physical and input-output coordinate bijections transport exact initialized
column identities. Input-column reindexing commutes with the normalized joining
parent. A full clean implementation through an isometric initialized inclusion
also determines the exact logical columns from its ambient columns.
\end{theorem}
\begin{proof}\leanok
Use the matrix reindexing product identity. Reindexing input columns passes
through the joining finite sum. To recover logical columns, multiply the
ambient identity by the adjoint of the initialized inclusion.
\end{proof}

\begin{theorem}[The reset label in the output inclusion]
    \label{thm:mpu_joining_reset_columns}
    \lean{MPUCircuit.initializedBasisMatrix_mul_basisResetOutput,
        MPUCircuit.joiningSpectatorColumns_basisResetOutput}
    \leanok
Multiplication of an initialized basis inclusion by the reset columns selects
its prescribed reset label. The same identity retains every spectator
configuration in the regrouped joining tensor.
\end{theorem}
\begin{proof}\leanok
Expand the basis inclusion. Its Kronecker delta leaves only the reset label;
the spectator delta is unchanged.
\end{proof}

\begin{theorem}[The joining packet lies in the parent interval]
    \label{thm:mpu_joining_packet_support}
    \lean{MPUCircuit.joiningPacketLogicalSites_mem_interval,
        MPUCircuit.joiningPacketSites_range_subset_interval}
    \leanok
The two shared bond registers and both joining flags lie inside the parent
interval support. Consequently the range of the actual ordered joining packet
is contained in that support.
\end{theorem}
\begin{proof}\leanok
The joining cut is strictly between the parent endpoints. Its two bond
registers belong to the respective interval boundary and interior sets, and
its flags belong to the parent interior. Apply the site-coordinate bijection.
\end{proof}

\begin{theorem}[The one-site case and empty padding]
    \label{thm:mpu_single_site_complete}
    \lean{MPUCircuit.completeCircuitSiteCount_one,
        MPUCircuit.circuitPaddingCount_eq_zero_of_two_le,
        MPUCircuit.completeCircuitSiteCount_eq_global_of_two_le,
        MPUCircuit.exists_clean_pair_circuit_of_siteCount_eq,
        MPUCircuit.completePhysicalConfigEmbedding_pullback_of_two_le,
        MPUCircuit.exists_singleSite_complete_clean_pair_circuit,
        MPUCircuit.exists_complete_clean_pair_circuit_of_global}
    \leanok
Every one-site unitary has a one-gate neighboring-pair circuit on the physical
site and one padding qudit, with the padding returned exactly to zero. When
$N\geq2$, the additional padding is empty: any exact circuit on the physical,
logical auxiliary, and scratch register transfers to the complete register
with the same gate count and physical columns.
\end{theorem}
\begin{proof}\leanok
At one site the complete register has two sites. Transport the padded physical
unitary along this equality of site counts. At longer lengths, append the
empty padding configuration and transport along the resulting site equality.
\end{proof}

\begin{theorem}[Derived support conditions for an interval merger]
    \label{thm:mpu_actual_merging_support}
    \lean{MPUCircuit.intervalConsecutiveSupport_left_subset,
        MPUCircuit.intervalConsecutiveSupport_right_subset,
        MPUCircuit.interval_children_mem_parent_supportedOperators,
        MPUCircuit.interval_initializedSites_range_subset_reindexed_parent,
        MPUCircuit.interval_joiningSuffix_range_subset_reindexed_parent}
    \leanok
Each of two adjacent child supports is contained in their parent support.
The actual selected initialized sites remain in the parent support after a
site-coordinate change. If that change puts the actual joining packet last,
its entire suffix range also lies in the transported parent support.
\end{theorem}
\begin{proof}\leanok
Use the consecutive support decomposition at the joining cut. The selected
sites are the actual parent auxiliaries, and the suffix range is the image
of the actual joining packet.
\end{proof}

\begin{theorem}[Joint child columns in the derived encoded coordinates]
    \label{thm:mpu_encoded_joint_columns}
    \lean{MPUCircuit.encodedJointColumns_of_coordinate_identities,
        MPUCircuit.normalizedJoiningParent_input_reindex_isIsometry,
        MPUCircuit.normalizedJoiningParent_spectator_reindex_isIsometry}
    \leanok
Given exact input and output coordinate identities, transport the actual
joint child columns to the compatible encoded joining columns. The
normalized parent remains an isometry under input reindexing and addition
of arbitrary spectator identity factors.
\end{theorem}
\begin{proof}\leanok
Compose the initialized basis inclusions and reindex matrix products. Apply
the normalized joining contraction before reindexing the input; its finite
sum is unchanged. The spectator identity preserves the isometry.
\end{proof}

\begin{theorem}[The encoded contraction gives the actual parent columns]
    \label{thm:mpu_encoded_actual_parent_columns}
    \lean{MPUCircuit.initializedBasisMatrix_pullback_rows,
        MPUCircuit.encodedJoiningParent_pullback_columns,
        MPUCircuit.regroupedIntervalJointOutputEmbedding_reset_apply,
        MPUCircuit.initializedBasisMatrix_regroupedOutput_mul_normalizedParent,
        MPUCircuit.encodedJoiningParent_minimalIntervalColumns}
    \leanok
Suppose the actual joint output inclusion has its derived product-coordinate
description and the joining metric is positive definite. Returning the encoded
normalized joining contraction to the original physical coordinates gives
exactly the weighted parent interval columns. Both shared bond labels and
joining flags are reset to zero. Arbitrary outside configurations are retained.
The parent-column and reset equalities are consequences, rather than additional
assumptions.
\end{theorem}
\begin{proof}\leanok
    \uses{thm:mpu_column_coordinate_transport,
        thm:mpu_joining_reset_columns,
        thm:mpu_actual_parent_isometry}
Transport the initialized basis inclusion and cancel the input-coordinate
bijection. The normalized joining identity gives the weighted parent tensor.
The concrete zero-label identity identifies the reset output inclusion with
the parent inclusion, including its unchanged outside register.
\end{proof}

\begin{theorem}[Joint initialized columns determine the parent invariant]
    \label{thm:mpu_parent_invariant_from_joint_columns}
    \lean{MPUCircuit.isMinimalIntervalColumnImplementation_of_joint_initialized_columns}
    \leanok
An operator satisfying the actual parent columns on the joint child input
satisfies every initialized parent column, with the outside configuration
unrestricted.
\end{theorem}
\begin{proof}\leanok
The derived joint input-coordinate bijection covers the parent input inclusion.
Evaluate the joint column identity at the inverse image of each parent column.
\end{proof}

\begin{theorem}[The actual minimal-interval circuit merger]
    \label{thm:mpu_actual_minimal_interval_merger}
    \lean{MPUCircuit.exists_minimalInterval_merging_circuit}
    \leanok
Let the minimal cut metrics be positive definite elements of their prefix Gram
affine hulls, with their dual metrics normalizing every element of those hulls.
Two exact supported implementations of adjacent intervals, with gate counts
$K_L,K_R$ and one common auxiliary workspace returned to zero on every logical
input, give an exact supported parent implementation. Its initialized columns
are the actual weighted parent interval, including arbitrary outside states.
Its gate count is at most
\[
  (2D+1)(K_L+K_R)+(2D+1)K+D(1+2F),
\]
where $K$ and $F$ are the uniform joining and reflection bounds. The same
workspace is returned to zero on every logical input.
\end{theorem}
\begin{proof}\leanok
    \uses{thm:mpu_permuted_encoded_merger,
        thm:mpu_actual_merging_support,
        thm:mpu_encoded_joint_columns,
        thm:mpu_encoded_actual_parent_columns,
        thm:mpu_parent_invariant_from_joint_columns,
        thm:mpu_uniform_numeric_resources}
Derive the shared cut encoding, product coordinates, selected initialized
sites, and normalized parent isometry from the actual minimal intervals.
Apply exact clean amplification in these coordinates and return to the original
site order. The preceding two column results establish the parent invariant.
The uniform joining estimate gives the displayed recurrence.
\end{proof}

\begin{theorem}[An explicit exact circuit bound for lengths at least two]
    \label{thm:mpu_explicit_exact_bounded_cut_circuit}
    \lean{MPUCircuit.exists_clean_exact_circuit_of_bounded_cutRank_of_two_le}
    \leanok
Let $d\geq2$ and $N\geq2$, and let $U$ be a unitary whose consecutive operator
cut ranks are at most $D$. There is an exact neighboring two-qudit circuit with
at most
\[
  c(d,D)(D+1)^3(2N)^{\lceil\log_2 D\rceil+6}
\]
gates, where $c(d,D)$ is the explicit polynomial coefficient in the uniform
resource estimate. All initialized logical auxiliary and scratch qudits are
returned to zero, and the physical action has the exact phase of $U$.
\end{theorem}
\begin{proof}\leanok
    \uses{thm:mpu_actual_minimal_interval_merger,
        thm:mpu_actual_interval_recursion,
        thm:mpu_actual_root_columns,
        thm:mpu_uniform_numeric_resources}
Obtain the minimal representation, normalized metrics, and leaf circuits from
$U$. Supply the actual merger to the midpoint-tree induction, use the uniform
resource estimate, and identify the root columns with $U$.
\end{proof}

\begin{theorem}[Exact circuits for arbitrary MPUs]
    \label{thm:mpu_general_exact_polynomial_circuits}
    \lean{MPUCircuit.exists_clean_exact_circuit_of_bounded_cutRank,
        MPUCircuit.exists_clean_exact_circuit_of_periodic_mpo_unitary,
        MPUCircuit.exists_clean_operatorNorm_approximate_circuit_of_bounded_cutRank,
        MPUCircuit.exists_clean_operatorNorm_approximate_circuit_of_periodic_mpo_unitary}
    \leanok
    Fix $d\geq2$. Let $U$ be a unitary on $N\geq1$ qudits of
    consecutive operator Schmidt ranks at most $D$.
    It has an exact circuit of arbitrary one- and two-qudit gates,
    with gate count and depth
    \[
      \operatorname{poly}_d(D)
        (N+1)^{O_d(\log(D+1))}
    \]
    and $O_d(N\log(D+1))$ auxiliary qudits which begin and end
    exactly in zero. The same form of bound holds on a line.
    In particular, both exact and arbitrarily accurate operator-norm
    circuit implementations have polynomial resources at fixed $D$.
\end{theorem}
\begin{proof}\leanok
    \uses{thm:mpu_affine_gram_balancing,
        thm:mpu_balanced_interval_minimal_cut_ranks,
        thm:mpu_simultaneous_balanced_cut_metrics,
        thm:mpu_endpoint_gram_normalization,
        thm:mpu_balanced_merger_success_amplitude,
        thm:mpu_amplification_reflections,
        thm:mpu_exact_subspace_amplification,
        thm:mpu_uniform_postselection_amplification,
        thm:mpu_uniform_success_attenuation,
        thm:mpu_exact_amplification_budget,
        thm:mpu_constant_branching_cost,
        thm:mpu_metrics_for_all_minimal_intervals,
        thm:mpu_common_register_layout,
        thm:mpu_clean_leaf_circuit,
        thm:mpu_compatible_bond_dilation,
        thm:mpu_selected_flag_reflection,
        thm:mpu_clean_circuit_placement,
        thm:mpu_exact_clean_amplification_circuit,
        thm:mpu_explicit_amplified_recursion_bound,
        thm:mpu_actual_tree_gate_count,
        thm:mpu_minimal_leaf_placement,
        thm:mpu_encoded_merging,
        thm:mpu_interval_initialization,
        thm:mpu_joint_child_columns,
        thm:mpu_uniform_numeric_resources,
        thm:mpu_amplification_support,
        thm:mpu_interval_column_encoding,
        thm:mpu_actual_parent_isometry,
        thm:mpu_supported_child_columns,
        thm:mpu_quantitative_site_permutation,
        thm:mpu_actual_interval_child_columns,
        thm:mpu_actual_minimal_interval_merger,
        thm:mpu_explicit_exact_bounded_cut_circuit,
        thm:mpu_actual_root_columns,
        thm:mpu_single_site_complete,
        thm:mpu_complete_register,
        thm:mpu_periodic_operator_cut_rank}
    Take a minimal open-boundary MPO representation, with cut dimensions
    $r_j\leq D$ and $r_0=r_N=1$. Let $f_{\mathbf a,\mathbf b}$ be
    its prefix rows at cut $j$. For each prefix density matrix $\rho$,
    define
    \[
      P_j(\rho)=\sum_{\mathbf a,\mathbf b,\mathbf b'}
        \rho_{\mathbf b',\mathbf b}
        f_{\mathbf a,\mathbf b}^\dagger
        f_{\mathbf a,\mathbf b'}.
    \]
    These are positive semidefinite. The maximally mixed prefix Gram
    is positive definite by minimality, as is the maximally mixed
    suffix Gram $\overline Q_j$. Global unitarity gives
    $\operatorname{Tr}(P_j(\rho)\overline Q_j)=1$.
    Apply the lemma to the real affine hull $\mathcal C_j$ of all
    prefix Grams. Obtain $P_j^*>0$ and
    $Q_j^*=(P_j^*)^{-1}/r_j$ normalizing every element of
    $\mathcal C_j$. At the endpoint cuts the affine space is $\{1\}$,
    so both metrics are one.

    For an interval from cut $j$ to cut $k$, write its tensor product
    as $T_{\mathbf a,\mathbf b}$. Any interval density matrix $\sigma$
    defines the positive transfer
    \[
      \mathcal T_\sigma(P)=
        \sum_{\mathbf a,\mathbf b,\mathbf b'}
          \sigma_{\mathbf b',\mathbf b}
          T_{\mathbf a,\mathbf b}^\dagger
          P T_{\mathbf a,\mathbf b'}.
    \]
    The identity
    $\mathcal T_\sigma(P_j(\rho))=P_k(\rho\otimes\sigma)$
    gives $\mathcal T_\sigma(\mathcal C_j)\subseteq\mathcal C_k$.
    Hence $\operatorname{Tr}(Q_k^*\mathcal T_\sigma(P_j^*))=1$
    for every density matrix $\sigma$. This proves the complete
    interval isometry identity by polarization, including all
    off-diagonal input entries. With $L_j=(P_j^*)^{1/2}$ and
    $R_j=(Q_j^*)^{1/2}$, its amplitudes are $L_jTR_k$.

    At a joining cut the contraction is
    $R_j^{-1}L_j^{-1}=\sqrt{r_j}\,I$, of norm $r_j\leq D$.
    The normalized contraction is a Bell bra and has constant
    success amplitude $1/r_j$ on the child images. Exact subspace
    amplification uses $O(D)$ child invocations per merge.
    One-site isometry extensions and Bell transformations have
    polynomial gate cost in $D$. Image reflections conjugate the initialized-auxiliary reflection
    by the full preparation, including both children, joining dilation,
    and attenuation. Their reversible conjunctions reuse one common
    $O(N\log(D+1))$ scratch pool. The balanced recursion therefore
    has polynomial cost at fixed $D$, with exponent $O(\log(D+1))$.
    Every inner flag is erased exactly; the root interval is $U$
    with no outer flags. Routing on a line adds at most a linear
    factor in the total number of registers.
    The complete proof and bounds are in
    \cite{RankTwoMPUCircuits2026}, Section~5.
\end{proof}

The metric is chosen in the positive semidefinite slice of an affine
space, which may contain real affine combinations with negative
coefficients. Its positivity supplies a physical isometry; the affine
identities supply its normalization. This removes the spectral condition
from the general nonuniform circuit statement of
\cite{Styliaris2025MPUCircuits}. It asserts circuit existence with
arbitrary gates, without an additional claim about computing exact
gate parameters from finite-precision input data. A periodic bond-$D$
presentation has an open presentation of bond at most $D^2$, and is
therefore covered at fixed $D$.
